% !TEX root = ../../main.tex
% Appendix E: Symbolic-pipeline outputs for the surveyed theories.

\section{Symbolic outputs}%
\label{SymbolicOutputs}

% Per-theory equation/constraint/term counts used in figure captions.
% Auto-generated by scripts/listings/compute_appendix_counts.py.
% Auto-generated by scripts/listings/compute_appendix_counts.py.
% Do not edit by hand. Re-run the script to refresh after listing edits.
%
% Counts are per-theory and reflect the FULL pipeline-emitted system
% (commented-out / pruned equations are counted alongside visible ones).

\newcommand{\BaselineEomCount}{6}
\newcommand{\BaselineConstraintCount}{0}
\newcommand{\BaselineTermCount}{19}

\newcommand{\EulerHeisenbergEomCount}{4}
\newcommand{\EulerHeisenbergConstraintCount}{0}
\newcommand{\EulerHeisenbergTermCount}{28}

\newcommand{\ReferenceEomCount}{38}
\newcommand{\ReferenceConstraintCount}{31}
\newcommand{\ReferenceTermCount}{161}

\newcommand{\MechanismEomCount}{38}
\newcommand{\MechanismConstraintCount}{25}
\newcommand{\MechanismTermCount}{380}

\newcommand{\NonminimalEomCount}{38}
\newcommand{\NonminimalConstraintCount}{31}
\newcommand{\NonminimalTermCount}{377}

\newcommand{\GeneralNonminimalEomCount}{38}
\newcommand{\GeneralNonminimalConstraintCount}{22}
\newcommand{\GeneralNonminimalTermCount}{806}

\newcommand{\ParityOddEomCount}{38}
\newcommand{\ParityOddConstraintCount}{22}
\newcommand{\ParityOddTermCount}{1854}

\newcommand{\DarkPhotonPlasmaEomCount}{18}
\newcommand{\DarkPhotonPlasmaConstraintCount}{3}
\newcommand{\DarkPhotonPlasmaTermCount}{168}


We collect representative symbolic outputs for the surveyed theories.
Each theory is linearised around its own background and then
dimensionally reduced to plane waves propagating along $\hat{z}$
(\cref{Linearisation}). The simulations initialise the gravitational
wave on $h_{xy}$, whose conversion product is $A_x$ at the rate
$\sin^{2}(\kappa \Bzero\,t / 2)$~\cite{boccaletti1970conversion}.
We show representative equations for the conversion-relevant
channel; equations of motion for the orthogonal polarisation, gauge
modes, and constraint fields are computed but not always displayed here.

We present Hamiltonian densities restricted to the graviton
and photon self-energy sector: torsion contributions are suppressed
and interaction terms are excluded,
so the displayed $\mathscr{H} \supset \dots$ is a
proper subset of the full canonical density (\cf~\cref{DiracBergmann}).

We also display the linearised kinetic matrix
$\mathcal{K}(\partial_t, \partial_z)$ of~\cref{KineticMatrix}.
Theories with velocity-pair fields
admit first-time-derivative contributions; these appear as
additional columns labelled $\dot{X}$ for each $q$-field $X$
whose velocity is referenced on a right-hand side. For those
theories the matrix is rectangular: $n_{\mathrm{eq}}$ rows (one
per equation of motion) and $n_q + n_v$ columns (one per dynamical
field plus one per velocity-pair reference); contracting row $i$
with the augmented state $\xi = (q, \dot{q})^\mathsf{T}$ gives
equation $i$ as $\mathcal{K}_{ij}\xi_j = 0$.

% --- E.1  Einstein--Maxwell Gertsenshtein baseline ------------------
\subsection{Einstein--Maxwell Gertsenshtein}%
\label{SymbolicBaseline}
Linearised around a uniform magnetic background
$\bm{B} = \Bzero\,\hat{x}$. The pure Einstein--Maxwell Lagrangian
gives a block-diagonal kinetic sub-matrix in the transverse photon
channels (\cref{Validation}). The symbolic output comprises
\BaselineEomCount{} equations of motion (\BaselineConstraintCount{}
constraints) with \BaselineTermCount{} RHS terms.

\medskip
\noindent\adjustbox{max width=\linewidth}{$\displaystyle
\setlength{\jot}{1pt}%
\begin{gathered}
% !TEX root = ../../../main.tex
  \begin{aligned}
    \mathscr{L} &= \tfrac{1}{\kappa^2} \mathcal{R} - \tfrac{1}{4} \tensor{\mathcal{F}}{_a_b} \tensor{g}{^a^c} \tensor{g}{^b^d} \tensor{\mathcal{F}}{_c_d}
  \end{aligned} \\[1ex]
  \begin{aligned}
    \partial_t^{2} A_{t} &= - \partial_z^2 A_{t} \\
    \partial_t^{2} A_{x} &= B_{0} \partial_z h_{xy} + \partial_z^2 A_{x} \\
    \partial_t^{2} A_{y} &= B_{0} \partial_z h_{yy} + \partial_z^2 A_{y} \\
    \partial_t^{2} A_{z} &= \partial_z^2 A_{z}
  \end{aligned} \\[1ex]
  \begin{aligned}
    -\tfrac{1}{\kappa^{2}} \, \partial_t^{2} h_{xy} &= \tfrac{B_{0}^2}{2} h_{xy} + B_{0} \partial_z A_{x} -\tfrac{1}{\kappa^{2}} \partial_z^2 h_{xy} \\
    -\tfrac{1}{\kappa^{2}} \, \partial_t^{2} h_{yy} &= \tfrac{B_{0}^2}{2} h_{yy} + B_{0} \partial_z A_{y} -\tfrac{1}{\kappa^{2}} \partial_z^2 h_{yy}
  \end{aligned} \\[1ex]
  \begin{aligned}
    \mathscr{H} &\supset \tfrac{B_{0}^2}{4} h_{xy}^2 + \tfrac{B_{0}^2}{4} h_{yy}^2 - \tfrac{1}{2} \partial_z A_{t}^2 + \tfrac{1}{2} \partial_z A_{x}^2 + \tfrac{1}{2} \partial_z A_{y}^2 \\
    &\quad + \tfrac{1}{2} \partial_z A_{z}^2 + \tfrac{1}{2 \, \kappa^2} \partial_z h_{xy}^2 + \tfrac{1}{2 \, \kappa^2} \partial_z h_{yy}^2 - \tfrac{1}{2} \dot{A}_{t}^2 \\
  &\quad + \tfrac{1}{2} \dot{A}_{x}^2 + \tfrac{1}{2} \dot{A}_{y}^2 + \tfrac{1}{2} \dot{A}_{z}^2 + \tfrac{1}{2 \, \kappa^2} \dot{h}_{xy}^2 + \tfrac{1}{2 \, \kappa^2} \dot{h}_{yy}^2
  \end{aligned}

\end{gathered}
$}
\medskip

The off-diagonal entries of the kinetic matrix below are the
Gertsenshtein couplings
$\mathcal{K}_{h_{xy}A_x} = \mathcal{K}_{A_xh_{xy}} = -\Bzero\,\partial_z$
(and the symmetric pair $h_{yy}\!\leftrightarrow\!A_y$).

\medskip
\noindent\adjustbox{max width=\linewidth}{$\displaystyle
\setlength{\jot}{0.5pt}%
\mathcal{K}(\partial_t,\partial_z) = 
\begin{array}{c|cccccc}
   & A_{t} & A_{x} & A_{y} & A_{z} & h_{xy} & h_{yy} \\
  \hline
  A_{t} & \partial_t^2 - \partial_z^2 & 0 & 0 & 0 & 0 & 0 \\
  A_{x} & 0 & \partial_t^2 - \partial_z^2 & 0 & 0 & -B_{0} \, \partial_z & 0 \\
  A_{y} & 0 & 0 & \partial_t^2 - \partial_z^2 & 0 & 0 & -B_{0} \, \partial_z \\
  A_{z} & 0 & 0 & 0 & \partial_t^2 - \partial_z^2 & 0 & 0 \\
  h_{xy} & 0 & -B_{0} \, \partial_z & 0 & 0 & -\tfrac{1}{\kappa^{2}} \, \partial_t^2 - \tfrac{B_{0}^2}{2} + \tfrac{1}{\kappa^{2}} \, \partial_z^2 & 0 \\
  h_{yy} & 0 & 0 & -B_{0} \, \partial_z & 0 & 0 & -\tfrac{1}{\kappa^{2}} \, \partial_t^2 - \tfrac{B_{0}^2}{2} + \tfrac{1}{\kappa^{2}} \, \partial_z^2
\end{array}

$}
\medskip

% --- E.2  Euler--Heisenberg QED 1-loop -------------------------------
\subsection{Euler--Heisenberg QED}%
\label{SymbolicEulerHeisenberg}
The Euler--Heisenberg corrections
of~\cref{LagEH} (\cref{Validation}) yield a symbolic output
comprises \EulerHeisenbergEomCount{} equations of motion
(\EulerHeisenbergConstraintCount{} constraints) with
\EulerHeisenbergTermCount{} RHS terms.

\medskip
\noindent\adjustbox{max width=\linewidth}{$\displaystyle
\setlength{\jot}{1pt}%
\begin{gathered}
  \begin{aligned}
    \mathscr{L} &= -\tfrac{1}{4} \tensor{\mathcal{F}}{_a_b} \tensor{\mathcal{F}}{^a^b} + \tfrac{\rho}{8} (\tensor{\mathcal{F}}{_a_b} \tensor{g}{^a^c} \tensor{g}{^b^d} \tensor{\mathcal{F}}{_c_d})^2 + \tfrac{\sigma}{8} (\tensor{\epsilon}{^a^b^c^d} \tensor{\mathcal{F}}{_a_b} \tensor{\mathcal{F}}{_c_d})^2
  \end{aligned} \\[1ex]
  \begin{aligned}
    \partial_t^{2} A_{t} &= - \partial_z^2 A_{t} + 2 \, B_{0}^2 \, \rho \partial_z^2 A_{t} -2 \, B_{0}^2 \, \rho \partial_z \dot{A}_{z} \\
    \left(-1 + 2 \, B_{0}^2 \, \rho - 16 \, B_{0}^2 \, \sigma\right) \, \partial_t^{2} A_{x} &= - \partial_z^2 A_{x} + 2 \, B_{0}^2 \, \rho \partial_z^2 A_{x} \\
    \left(-1 + 2 \, B_{0}^2 \, \rho\right) \, \partial_t^{2} A_{y} &= - \partial_z^2 A_{y} + 6 \, B_{0}^2 \, \rho \partial_z^2 A_{y} \\
    \left(-1 + 2 \, B_{0}^2 \, \rho\right) \, \partial_t^{2} A_{z} &= - \partial_z^2 A_{z} + 2 \, B_{0}^2 \, \rho \partial_z \dot{A}_{t}
  \end{aligned} \\[1ex]
  \begin{aligned}
    \mathscr{H} &= -\tfrac{1}{2} \partial_z A_{t}^2 + B_{0}^2 \, \rho \partial_z A_{t}^2 + \tfrac{1}{2} \partial_z A_{x}^2 - B_{0}^2 \, \rho \partial_z A_{x}^2 + \tfrac{1}{2} \partial_z A_{y}^2 \\
    &\quad- 3 \, B_{0}^2 \, \rho \partial_z A_{y}^2 + \tfrac{1}{2} \partial_z A_{z}^2 - \tfrac{1}{2} \dot{A}_{t}^2 + \tfrac{1}{2} \dot{A}_{x}^2 - B_{0}^2 \, \rho \dot{A}_{x}^2 \\
  &\quad + 8 \, B_{0}^2 \, \sigma \dot{A}_{x}^2 + \tfrac{1}{2} \dot{A}_{y}^2 - B_{0}^2 \, \rho \dot{A}_{y}^2 + \tfrac{1}{2} \dot{A}_{z}^2 - B_{0}^2 \, \rho \dot{A}_{z}^2
  \end{aligned}

\end{gathered}
$}
\medskip


% --- E.3  Minimally-coupled Einstein--Cartan ------------------------
\subsection{Minimally-coupled Einstein--Cartan}%
\label{SymbolicReference}
Three quadratic torsion invariants with
couplings $\beta_{1,2,3}$ are added to Einstein--Maxwell; the torsion
components enter only in algebraic equations with LHS $=0$ (\cref{SectorSurvey}).
The symbolic output comprises \ReferenceEomCount{} equations of motion
(\ReferenceConstraintCount{} constraints) with \ReferenceTermCount{}
RHS terms.

\medskip
\noindent\adjustbox{max width=\linewidth}{$\displaystyle
\setlength{\jot}{0.2pt}%
\begin{gathered}
% !TEX root = ../../../main.tex
  \begin{aligned}
    \mathscr{L} &= \tfrac{1}{\kappa^2} \tilde{\mathcal{R}} + \alpha_{1} \tensor{T}{_a_b_c} \tensor{g}{^a^d} \tensor{g}{^b^e} \tensor{g}{^c^f} \tensor{T}{_d_e_f} + \alpha_{2} \tensor{T}{_a_b_c} \tensor{g}{^b^d} \tensor{g}{^a^e} \tensor{g}{^c^f} \tensor{T}{_d_e_f} + \alpha_{3} \tensor{T}{^b_b_a} \tensor{g}{^a^d} \tensor{T}{^c_c_d} - \tfrac{1}{4} \tensor{\mathcal{F}}{_a_b} \tensor{g}{^a^c} \tensor{g}{^b^d} \tensor{\mathcal{F}}{_c_d}
  \end{aligned} \\[1ex]
  \begin{aligned}
    \partial_t^{2} A_{t} &= -B_{0} \partial_z h_{ty} - \partial_z^2 A_{t} \\
    \partial_t^{2} A_{x} &= B_{0} \partial_z h_{xy} + \partial_z^2 A_{x} \\
    \partial_t^{2} A_{y} &= \tfrac{B_{0}}{2} \partial_z h_{tt} -\tfrac{1}{2} B_{0} \partial_z h_{xx} + \tfrac{B_{0}}{2} \partial_z h_{yy} + \tfrac{B_{0}}{2} \partial_z h_{zz} + \partial_z^2 A_{y} -B_{0} \partial_t h_{tz} \\
    \partial_t^{2} A_{z} &= \partial_z^2 A_{z} + B_{0} \partial_t h_{ty}
  \end{aligned} \\[1ex]
  \begin{aligned}
    0 &= \tfrac{B_{0}^2}{8} h_{tt} + \tfrac{B_{0}^2}{8} h_{xx} -\tfrac{1}{8} B_{0}^2 h_{yy} -\tfrac{1}{8} B_{0}^2 h_{zz} -\tfrac{1}{2} B_{0} \partial_z A_{y} + \tfrac{1}{2 \, \kappa^2} \partial_z^2 h_{xx} + \tfrac{1}{2 \, \kappa^2} \partial_z^2 h_{yy} \\
    0 &= -\tfrac{1}{2} B_{0}^2 h_{tx} -\tfrac{1}{\kappa^{2}} \partial_z^2 h_{tx} + \tfrac{1}{\kappa^{2}} \partial_z \dot{h}_{xz} \\
    0 &= \tfrac{B_{0}^2}{2} h_{ty} + B_{0} \partial_z A_{t} -\tfrac{1}{\kappa^{2}} \partial_z^2 h_{ty} -B_{0} \partial_t A_{z} + \tfrac{1}{\kappa^{2}} \partial_z \dot{h}_{yz} \\
    0 &= \tfrac{B_{0}^2}{2} h_{tz} + B_{0} \partial_t A_{y} -\tfrac{1}{\kappa^{2}} \partial_z \dot{h}_{xx} -\tfrac{1}{\kappa^{2}} \partial_z \dot{h}_{yy} \\
    0 &= \tfrac{B_{0}^2}{8} h_{tt} + \tfrac{B_{0}^2}{8} h_{xx} + \tfrac{B_{0}^2}{8} h_{yy} + \tfrac{B_{0}^2}{8} h_{zz} + \tfrac{B_{0}}{2} \partial_z A_{y} + \tfrac{1}{2 \, \kappa^2} \partial_z^2 h_{tt} - \tfrac{1}{2 \, \kappa^2} \partial_z^2 h_{yy} - \tfrac{1}{\kappa^{2}} \partial_z \dot{h}_{tz}  + \tfrac{1}{2 \, \kappa^2} \partial_t^2 h_{yy} + \tfrac{1}{2 \, \kappa^2} \partial_t^2 h_{zz} \\
    -\tfrac{1}{\kappa^{2}} \, \partial_t^{2} h_{xy} &= \tfrac{B_{0}^2}{2} h_{xy} + B_{0} \partial_z A_{x} -\tfrac{1}{\kappa^{2}} \partial_z^2 h_{xy} \\
    -\tfrac{1}{\kappa^{2}} \, \partial_t^{2} h_{xz} &= \tfrac{B_{0}^2}{2} h_{xz} -\tfrac{1}{\kappa^{2}} \partial_z \dot{h}_{tx} \\
    0 &= -\tfrac{1}{8} B_{0}^2 h_{tt} + \tfrac{B_{0}^2}{8} h_{xx} - \tfrac{3 \, B_{0}^2}{8} h_{yy} - \tfrac{1}{8} B_{0}^2 h_{zz} - \tfrac{1}{2} B_{0} \partial_z A_{y} + \tfrac{1}{2 \, \kappa^2} \partial_z^2 h_{tt} - \tfrac{1}{2 \, \kappa^2} \partial_z^2 h_{xx} - \tfrac{1}{\kappa^{2}} \partial_z \dot{h}_{tz} + \tfrac{1}{2 \, \kappa^2} \partial_t^2 h_{xx} + \tfrac{1}{2 \, \kappa^2} \partial_t^2 h_{zz} \\
    -\tfrac{1}{\kappa^{2}} \, \partial_t^{2} h_{yz} &= \tfrac{B_{0}^2}{2} h_{yz} -\tfrac{1}{\kappa^{2}} \partial_z \dot{h}_{ty} \\
    0 &= -\tfrac{1}{8} B_{0}^2 h_{tt} + \tfrac{B_{0}^2}{8} h_{xx} -\tfrac{1}{8} B_{0}^2 h_{yy} -\tfrac{3 \, B_{0}^2}{8} h_{zz} -\tfrac{1}{2} B_{0} \partial_z A_{y} + \tfrac{1}{2 \, \kappa^2} \partial_t^2 h_{xx} + \tfrac{1}{2 \, \kappa^2} \partial_t^2 h_{yy}
  \end{aligned} \\[1ex]
  \begin{aligned}
    0 &= 4 \, \alpha_{1} \mathcal{T}_{ttx} + 2 \, \alpha_{2} \mathcal{T}_{ttx} + 2 \, \alpha_{3} \mathcal{T}_{ttx} + \tfrac{1}{2 \, \kappa^2} \mathcal{T}_{ttx} -2 \, \alpha_{3} \mathcal{T}_{yxy} + \tfrac{3}{\kappa^2} \mathcal{T}_{yxy} -2 \, \alpha_{3} \mathcal{T}_{zxz} + \tfrac{3}{\kappa^2} \mathcal{T}_{zxz} \\
    0 &= 4 \, \alpha_{1} \mathcal{T}_{tty} + 2 \, \alpha_{2} \mathcal{T}_{tty} + 2 \, \alpha_{3} \mathcal{T}_{tty} + \tfrac{1}{2 \, \kappa^2} \mathcal{T}_{tty} -2 \, \alpha_{3} \mathcal{T}_{zyz} + \tfrac{3}{\kappa^2} \mathcal{T}_{zyz} + 2 \, \alpha_{3} \mathcal{T}_{xxy} -\tfrac{3}{\kappa^2} \mathcal{T}_{xxy} \\
    0 &= 4 \, \alpha_{1} \mathcal{T}_{xxz} + 2 \, \alpha_{2} \mathcal{T}_{xxz} + 2 \, \alpha_{3} \mathcal{T}_{xxz} + \tfrac{1}{2 \, \kappa^2} \mathcal{T}_{xxz} + 2 \, \alpha_{3} \mathcal{T}_{yyz} -\tfrac{3}{\kappa^2} \mathcal{T}_{yyz} + 2 \, \alpha_{3} \mathcal{T}_{ttz} -\tfrac{3}{\kappa^2} \mathcal{T}_{ttz} \\
    0 &= 4 \, \alpha_{1} \mathcal{T}_{xyz} + \tfrac{2}{\kappa^2} \mathcal{T}_{xyz} + 2 \, \alpha_{2} \mathcal{T}_{yxz} + \tfrac{3}{2 \, \kappa^2} \mathcal{T}_{yxz} -2 \, \alpha_{2} \mathcal{T}_{zxy} -\tfrac{3}{2 \, \kappa^2} \mathcal{T}_{zxy} \\
    0 &= -4 \, \alpha_{1} \mathcal{T}_{ytx} -\tfrac{2}{\kappa^2} \mathcal{T}_{ytx} -2 \, \alpha_{2} \mathcal{T}_{txy} -\tfrac{3}{2 \, \kappa^2} \mathcal{T}_{txy} -2 \, \alpha_{2} \mathcal{T}_{xty} -\tfrac{3}{2 \, \kappa^2} \mathcal{T}_{xty} \\
    0 &= -4 \, \alpha_{1} \mathcal{T}_{yty} -2 \, \alpha_{2} \mathcal{T}_{yty} -2 \, \alpha_{3} \mathcal{T}_{yty} -\tfrac{1}{2 \, \kappa^2} \mathcal{T}_{yty} -2 \, \alpha_{3} \mathcal{T}_{ztz} + \tfrac{3}{\kappa^2} \mathcal{T}_{ztz} -2 \, \alpha_{3} \mathcal{T}_{xtx} + \tfrac{3}{\kappa^2} \mathcal{T}_{xtx} \\
    0 &= -4 \, \alpha_{1} \mathcal{T}_{ytz} -\tfrac{2}{\kappa^2} \mathcal{T}_{ytz} -2 \, \alpha_{2} \mathcal{T}_{zty} -\tfrac{3}{2 \, \kappa^2} \mathcal{T}_{zty} + 2 \, \alpha_{2} \mathcal{T}_{tyz} + \tfrac{3}{2 \, \kappa^2} \mathcal{T}_{tyz} \\
    0 &= -2 \, \alpha_{3} \mathcal{T}_{ttx} + \tfrac{3}{\kappa^2} \mathcal{T}_{ttx} + 4 \, \alpha_{1} \mathcal{T}_{yxy} + 2 \, \alpha_{2} \mathcal{T}_{yxy} + 2 \, \alpha_{3} \mathcal{T}_{yxy} + \tfrac{1}{2 \, \kappa^2} \mathcal{T}_{yxy} + 2 \, \alpha_{3} \mathcal{T}_{zxz} -\tfrac{3}{\kappa^2} \mathcal{T}_{zxz} \\
    0 &= 2 \, \alpha_{2} \mathcal{T}_{xyz} + \tfrac{3}{2 \, \kappa^2} \mathcal{T}_{xyz} + 4 \, \alpha_{1} \mathcal{T}_{yxz} + \tfrac{2}{\kappa^2} \mathcal{T}_{yxz} + 2 \, \alpha_{2} \mathcal{T}_{zxy} + \tfrac{3}{2 \, \kappa^2} \mathcal{T}_{zxy} \\
    0 &= 2 \, \alpha_{3} \mathcal{T}_{xxz} -\tfrac{3}{\kappa^2} \mathcal{T}_{xxz} + 4 \, \alpha_{1} \mathcal{T}_{yyz} + 2 \, \alpha_{2} \mathcal{T}_{yyz} + 2 \, \alpha_{3} \mathcal{T}_{yyz} + \tfrac{1}{2 \, \kappa^2} \mathcal{T}_{yyz} + 2 \, \alpha_{3} \mathcal{T}_{ttz} -\tfrac{3}{\kappa^2} \mathcal{T}_{ttz} \\
    0 &= -4 \, \alpha_{1} \mathcal{T}_{ztx} -\tfrac{2}{\kappa^2} \mathcal{T}_{ztx} -2 \, \alpha_{2} \mathcal{T}_{txz} -\tfrac{3}{2 \, \kappa^2} \mathcal{T}_{txz} -2 \, \alpha_{2} \mathcal{T}_{xtz} -\tfrac{3}{2 \, \kappa^2} \mathcal{T}_{xtz} \\
    0 &= -2 \, \alpha_{2} \mathcal{T}_{ytz} -\tfrac{3}{2 \, \kappa^2} \mathcal{T}_{ytz} -4 \, \alpha_{1} \mathcal{T}_{zty} -\tfrac{2}{\kappa^2} \mathcal{T}_{zty} -2 \, \alpha_{2} \mathcal{T}_{tyz} -\tfrac{3}{2 \, \kappa^2} \mathcal{T}_{tyz} \\
    0 &= 2 \, \alpha_{3} \mathcal{T}_{xxz} -\tfrac{3}{\kappa^2} \mathcal{T}_{xxz} + 2 \, \alpha_{3} \mathcal{T}_{yyz} -\tfrac{3}{\kappa^2} \mathcal{T}_{yyz} + 4 \, \alpha_{1} \mathcal{T}_{ttz} + 2 \, \alpha_{2} \mathcal{T}_{ttz} + 2 \, \alpha_{3} \mathcal{T}_{ttz} + \tfrac{1}{2 \, \kappa^2} \mathcal{T}_{ttz} \\
    0 &= -2 \, \alpha_{3} \mathcal{T}_{yty} + \tfrac{3}{\kappa^2} \mathcal{T}_{yty} -4 \, \alpha_{1} \mathcal{T}_{ztz} -2 \, \alpha_{2} \mathcal{T}_{ztz} -2 \, \alpha_{3} \mathcal{T}_{ztz} -\tfrac{1}{2 \, \kappa^2} \mathcal{T}_{ztz} -2 \, \alpha_{3} \mathcal{T}_{xtx} + \tfrac{3}{\kappa^2} \mathcal{T}_{xtx} \\
    0 &= -2 \, \alpha_{2} \mathcal{T}_{xyz} -\tfrac{3}{2 \, \kappa^2} \mathcal{T}_{xyz} + 2 \, \alpha_{2} \mathcal{T}_{yxz} + \tfrac{3}{2 \, \kappa^2} \mathcal{T}_{yxz} + 4 \, \alpha_{1} \mathcal{T}_{zxy} + \tfrac{2}{\kappa^2} \mathcal{T}_{zxy} \\
    0 &= -2 \, \alpha_{3} \mathcal{T}_{ttx} + \tfrac{3}{\kappa^2} \mathcal{T}_{ttx} + 2 \, \alpha_{3} \mathcal{T}_{yxy} -\tfrac{3}{\kappa^2} \mathcal{T}_{yxy} + 4 \, \alpha_{1} \mathcal{T}_{zxz} + 2 \, \alpha_{2} \mathcal{T}_{zxz} + 2 \, \alpha_{3} \mathcal{T}_{zxz} + \tfrac{1}{2 \, \kappa^2} \mathcal{T}_{zxz} \\
    0 &= -2 \, \alpha_{3} \mathcal{T}_{tty} + \tfrac{3}{\kappa^2} \mathcal{T}_{tty} + 4 \, \alpha_{1} \mathcal{T}_{zyz} + 2 \, \alpha_{2} \mathcal{T}_{zyz} + 2 \, \alpha_{3} \mathcal{T}_{zyz} + \tfrac{1}{2 \, \kappa^2} \mathcal{T}_{zyz} -2 \, \alpha_{3} \mathcal{T}_{xxy} + \tfrac{3}{\kappa^2} \mathcal{T}_{xxy} \\
    0 &= -2 \, \alpha_{2} \mathcal{T}_{ytx} -\tfrac{3}{2 \, \kappa^2} \mathcal{T}_{ytx} -4 \, \alpha_{1} \mathcal{T}_{txy} -\tfrac{2}{\kappa^2} \mathcal{T}_{txy} + 2 \, \alpha_{2} \mathcal{T}_{xty} + \tfrac{3}{2 \, \kappa^2} \mathcal{T}_{xty} \\
    0 &= -2 \, \alpha_{2} \mathcal{T}_{ztx} -\tfrac{3}{2 \, \kappa^2} \mathcal{T}_{ztx} -4 \, \alpha_{1} \mathcal{T}_{txz} -\tfrac{2}{\kappa^2} \mathcal{T}_{txz} + 2 \, \alpha_{2} \mathcal{T}_{xtz} + \tfrac{3}{2 \, \kappa^2} \mathcal{T}_{xtz} \\
    0 &= 2 \, \alpha_{2} \mathcal{T}_{ytz} + \tfrac{3}{2 \, \kappa^2} \mathcal{T}_{ytz} -2 \, \alpha_{2} \mathcal{T}_{zty} -\tfrac{3}{2 \, \kappa^2} \mathcal{T}_{zty} -4 \, \alpha_{1} \mathcal{T}_{tyz} -\tfrac{2}{\kappa^2} \mathcal{T}_{tyz} \\
    0 &= -2 \, \alpha_{3} \mathcal{T}_{yty} + \tfrac{3}{\kappa^2} \mathcal{T}_{yty} -2 \, \alpha_{3} \mathcal{T}_{ztz} + \tfrac{3}{\kappa^2} \mathcal{T}_{ztz} -4 \, \alpha_{1} \mathcal{T}_{xtx} -2 \, \alpha_{2} \mathcal{T}_{xtx} -2 \, \alpha_{3} \mathcal{T}_{xtx} -\tfrac{1}{2 \, \kappa^2} \mathcal{T}_{xtx} \\
    0 &= -2 \, \alpha_{2} \mathcal{T}_{ytx} -\tfrac{3}{2 \, \kappa^2} \mathcal{T}_{ytx} + 2 \, \alpha_{2} \mathcal{T}_{txy} + \tfrac{3}{2 \, \kappa^2} \mathcal{T}_{txy} -4 \, \alpha_{1} \mathcal{T}_{xty} -\tfrac{2}{\kappa^2} \mathcal{T}_{xty} \\
    0 &= -2 \, \alpha_{2} \mathcal{T}_{ztx} -\tfrac{3}{2 \, \kappa^2} \mathcal{T}_{ztx} + 2 \, \alpha_{2} \mathcal{T}_{txz} + \tfrac{3}{2 \, \kappa^2} \mathcal{T}_{txz} -4 \, \alpha_{1} \mathcal{T}_{xtz} -\tfrac{2}{\kappa^2} \mathcal{T}_{xtz} \\
    0 &= 2 \, \alpha_{3} \mathcal{T}_{tty} -\tfrac{3}{\kappa^2} \mathcal{T}_{tty} -2 \, \alpha_{3} \mathcal{T}_{zyz} + \tfrac{3}{\kappa^2} \mathcal{T}_{zyz} + 4 \, \alpha_{1} \mathcal{T}_{xxy} + 2 \, \alpha_{2} \mathcal{T}_{xxy} + 2 \, \alpha_{3} \mathcal{T}_{xxy} + \tfrac{1}{2 \, \kappa^2} \mathcal{T}_{xxy}
  \end{aligned} \\[1ex]
  \begin{aligned}
    \mathscr{H} &\supset -\tfrac{1}{16} B_{0}^2 h_{tt}^2 + \tfrac{B_{0}^2}{4} h_{tx}^2 - \tfrac{1}{4} B_{0}^2 h_{ty}^2 - \tfrac{1}{4} B_{0}^2 h_{tz}^2 - \tfrac{1}{8} B_{0}^2 h_{tt} \, h_{xx} - \tfrac{1}{16} B_{0}^2 h_{xx}^2 + \tfrac{B_{0}^2}{4} h_{xy}^2 + \tfrac{B_{0}^2}{4} h_{xz}^2 + \tfrac{B_{0}^2}{8} h_{tt} \, h_{yy} - \tfrac{1}{8} B_{0}^2 h_{xx} \, h_{yy} \\
  &\quad + \tfrac{3 \, B_{0}^2}{16} h_{yy}^2 + \tfrac{B_{0}^2}{4} h_{yz}^2 + \tfrac{B_{0}^2}{8} h_{tt} \, h_{zz} - \tfrac{1}{8} B_{0}^2 h_{xx} \, h_{zz} + \tfrac{B_{0}^2}{8} h_{yy} \, h_{zz} + \tfrac{3 \, B_{0}^2}{16} h_{zz}^2 - \tfrac{1}{2} \partial_z A_{t}^2 + \tfrac{1}{2} \partial_z A_{x}^2 + \tfrac{1}{2} \partial_z A_{y}^2 + \tfrac{1}{2} \partial_z A_{z}^2 \\
  &\quad - \tfrac{1}{2 \, \kappa^2} \partial_z h_{tx}^2 - \tfrac{1}{2 \, \kappa^2} \partial_z h_{ty}^2 + \tfrac{1}{2 \, \kappa^2} \partial_z h_{tt} \, \partial_z h_{xx} + \tfrac{1}{2 \, \kappa^2} \partial_z h_{xy}^2 + \tfrac{1}{2 \, \kappa^2} \partial_z h_{tt} \, \partial_z h_{yy} - \tfrac{1}{2 \, \kappa^2} \partial_z h_{xx} \, \partial_z h_{yy} - \tfrac{1}{2} \dot{A}_{t}^2 + \tfrac{1}{2} \dot{A}_{x}^2 + \tfrac{1}{2} \dot{A}_{y}^2 \\
  &\quad + \tfrac{1}{2} \dot{A}_{z}^2 + \tfrac{1}{2 \, \kappa^2} \dot{h}_{xy}^2 + \tfrac{1}{2 \, \kappa^2} \dot{h}_{xz}^2 - \tfrac{1}{2 \, \kappa^2} \dot{h}_{xx} \, \dot{h}_{yy} + \tfrac{1}{2 \, \kappa^2} \dot{h}_{yz}^2 - \tfrac{1}{2 \, \kappa^2} \dot{h}_{xx} \, \dot{h}_{zz} - \tfrac{1}{2 \, \kappa^2} \dot{h}_{yy} \, \dot{h}_{zz} - \tfrac{2}{\kappa^2} h_{xz} \, \partial_t \partial_z h_{tx} - \tfrac{2}{\kappa^2} h_{yz} \, \partial_t \partial_z h_{ty} \\
  &\quad + \tfrac{1}{\kappa^{2}} h_{tt} \, \partial_t \partial_z h_{tz} + \tfrac{1}{\kappa^{2}} h_{xx} \, \partial_t \partial_z h_{tz} + \tfrac{1}{\kappa^{2}} h_{yy} \, \partial_t \partial_z h_{tz} - \tfrac{1}{\kappa^{2}} h_{zz} \, \partial_t \partial_z h_{tz} + \tfrac{2}{\kappa^2} h_{tz} \, \partial_t \partial_z h_{xx} - \tfrac{2}{\kappa^2} h_{tx} \, \partial_t \partial_z h_{xz} \\
  &\quad + \tfrac{2}{\kappa^2} h_{tz} \, \partial_t \partial_z h_{yy} - \tfrac{2}{\kappa^2} h_{ty} \, \partial_t \partial_z h_{yz}
  \end{aligned}

\end{gathered}
$}
\medskip


% --- E.4  Propagating PGT with the curvature-squared term -----------
\subsection{$\MAGFt{^2}$--PGT}%
\label{SymbolicMechanism}
The $\Alp{1}\MAGFt{^2}$
curvature-squared term promotes the metric-perturbation sector to
a fourth-order system. We include this theory \emph{for
inspection} of the resulting equation-of-motion structure rather
than as a simulated theory. The $\Alp{1}$-weighted
contributions land only on the diagonal $h_{ii}$ kinetic equations
(see, for instance, the $2\,\Alp{1}\,\partial_t^{4}h_{xx}$ term) and are absent from the
off-diagonal $h_{xy}$ equation that mediates the Gertsenshtein
channel.

The Hamiltonian density is suppressed;
see~\cref{PerturbativeReduction}. The symbolic output comprises
\MechanismEomCount{} equations of motion (\MechanismConstraintCount{}
constraints) with \MechanismTermCount{} RHS terms.

\medskip
\noindent\adjustbox{max width=\linewidth}{$\displaystyle
\setlength{\jot}{1pt}%
\begin{gathered}
% !TEX root = ../../../main.tex
  \begin{aligned}
    \mathscr{L} &= \tfrac{1}{\kappa^2} \tilde{\mathcal{R}} + \Bet{1} \tensor{T}{_a_b_c} \tensor{g}{^a^d} \tensor{g}{^b^e} \tensor{g}{^c^f} \tensor{T}{_d_e_f} + \Bet{2} \tensor{T}{_a_b_c} \tensor{g}{^b^d} \tensor{g}{^a^e} \tensor{g}{^c^f} \tensor{T}{_d_e_f} \\
    &\quad + \Bet{3} \tensor{T}{^b_b_a} \tensor{g}{^a^d} \tensor{T}{^c_c_d} + \Alp{1} \tilde{\mathcal{R}}^2 - \tfrac{1}{4} \tensor{\mathcal{F}}{_a_b} \tensor{g}{^a^c} \tensor{g}{^b^d} \tensor{\mathcal{F}}{_c_d}
  \end{aligned} \\[1ex]
  \begin{aligned}
    % \partial_t^{2} A_{t} &= -B_{0} \partial_z h_{ty} - \partial_z^2 A_{t} \\
    \partial_t^{2} A_{x} &= B_{0} \partial_z h_{xy} + \partial_z^2 A_{x} \\
    % \partial_t^{2} A_{y} &= \tfrac{B_{0}}{2} \partial_z h_{tt} -\tfrac{1}{2} B_{0} \partial_z h_{xx} + \tfrac{B_{0}}{2} \partial_z h_{yy} + \tfrac{B_{0}}{2} \partial_z h_{zz} + \partial_z^2 A_{y} -B_{0} \partial_t h_{tz} \\
    % \partial_t^{2} A_{z} &= \partial_z^2 A_{z} + B_{0} \partial_t h_{ty}
  \end{aligned} \\[1ex]
  \begin{aligned}
    -\tfrac{1}{\kappa^{2}} \, \partial_t^{2} h_{xy} &= \tfrac{B_{0}^2}{2} h_{xy} + B_{0} \partial_z A_{x} -\tfrac{1}{\kappa^{2}} \partial_z^2 h_{xy} \\
    % -\tfrac{1}{\kappa^{2}} \, \partial_t^{2} h_{xz} &= \tfrac{B_{0}^2}{2} h_{xz} -\tfrac{1}{\kappa^{2}} \partial_z \dot{h}_{tx} \\
    % -\tfrac{1}{\kappa^{2}} \, \partial_t^{2} h_{yz} &= \tfrac{B_{0}^2}{2} h_{yz} -\tfrac{1}{\kappa^{2}} \partial_z \dot{h}_{ty} \\
    % 0 &= -\tfrac{1}{2} B_{0}^2 h_{tx} -\tfrac{1}{\kappa^{2}} \partial_z^2 h_{tx} + \tfrac{1}{\kappa^{2}} \partial_z \dot{h}_{xz} \\
    % 0 &= \tfrac{B_{0}^2}{2} h_{ty} + B_{0} \partial_z A_{t} -\tfrac{1}{\kappa^{2}} \partial_z^2 h_{ty} -B_{0} \partial_t A_{z} + \tfrac{1}{\kappa^{2}} \partial_z \dot{h}_{yz}
  \end{aligned} \\[1ex]
  \begin{aligned}
  %   0 &= \tfrac{B_{0}^2}{8} h_{tt} + \tfrac{B_{0}^2}{8} h_{xx} - \tfrac{1}{8} B_{0}^2 h_{yy} - \tfrac{1}{8} B_{0}^2 h_{zz} - \tfrac{1}{2} B_{0} \partial_z A_{y} + \tfrac{1}{2 \, \kappa^2} \partial_z^2 h_{xx} + \tfrac{1}{2 \, \kappa^2} \partial_z^2 h_{yy} + 5 \, \Alp{1} \partial_z^{3} \mathcal{T}_{xxz} \\
  % &\quad + 5 \, \Alp{1} \partial_z^{3} \mathcal{T}_{yyz} + 5 \, \Alp{1} \partial_z^{3} \mathcal{T}_{ttz} + 2 \, \Alp{1} \partial_z^4 h_{tt} - 2 \, \Alp{1} \partial_z^4 h_{xx} - 2 \, \Alp{1} \partial_z^4 h_{yy} + 5 \, \Alp{1} \partial_z^2 \dot{\mathcal{T}}_{yty} + 5 \, \Alp{1} \partial_z^2 \dot{\mathcal{T}}_{ztz} + 5 \, \Alp{1} \partial_z^2 \dot{\mathcal{T}}_{xtx} \\
  % &\quad - 4 \, \Alp{1} \partial_z^{3} \dot{h}_{tz} + 2 \, \Alp{1} \partial_t^{2} \partial_z^{2} h_{xx} + 2 \, \Alp{1} \partial_t^{2} \partial_z^{2} h_{yy} + 2 \, \Alp{1} \partial_t^{2} \partial_z^{2} h_{zz} \\
  %   0 &= \tfrac{B_{0}^2}{2} h_{tz} + B_{0} \partial_t A_{y} - \tfrac{1}{\kappa^{2}} \partial_z \dot{h}_{xx} - \tfrac{1}{\kappa^{2}} \partial_z \dot{h}_{yy} - 10 \, \Alp{1} \partial_z^2 \dot{\mathcal{T}}_{xxz} - 10 \, \Alp{1} \partial_z^2 \dot{\mathcal{T}}_{yyz} - 10 \, \Alp{1} \partial_z^2 \dot{\mathcal{T}}_{ttz} - 4 \, \Alp{1} \partial_z^{3} \dot{h}_{tt} \\
  % &\quad + 4 \, \Alp{1} \partial_z^{3} \dot{h}_{xx} + 4 \, \Alp{1} \partial_z^{3} \dot{h}_{yy} - 10 \, \Alp{1} \partial_t^{2} \partial_z \mathcal{T}_{yty} - 10 \, \Alp{1} \partial_t^{2} \partial_z \mathcal{T}_{ztz} - 10 \, \Alp{1} \partial_t^{2} \partial_z \mathcal{T}_{xtx} + 8 \, \Alp{1} \partial_t^{2} \partial_z^{2} h_{tz} - 4 \, \Alp{1} \partial_t^{3} \partial_z h_{xx} \\
  % &\quad - 4 \, \Alp{1} \partial_t^{3} \partial_z h_{yy} - 4 \, \Alp{1} \partial_t^{3} \partial_z h_{zz} \\
    2 \, \Alp{1} \, \partial_t^{4} h_{xx} &= -\tfrac{1}{8} B_{0}^2 h_{tt} - \tfrac{1}{8} B_{0}^2 h_{xx} - \tfrac{1}{8} B_{0}^2 h_{yy} - \tfrac{1}{8} B_{0}^2 h_{zz} - \tfrac{1}{2} B_{0} \partial_z A_{y} \\
      &\quad - \tfrac{1}{2 \, \kappa^2} \partial_z^2 h_{tt} + \tfrac{1}{2 \, \kappa^2} \partial_z^2 h_{yy} + 5 \, \Alp{1} \partial_z^{3} \mathcal{T}_{xxz} + 5 \, \Alp{1} \partial_z^{3} \mathcal{T}_{yyz} + 5 \, \Alp{1} \partial_z^{3} \mathcal{T}_{ttz} \\
      &\quad + 2 \, \Alp{1} \partial_z^4 h_{tt} - 2 \, \Alp{1} \partial_z^4 h_{xx} - 2 \, \Alp{1} \partial_z^4 h_{yy} + \tfrac{1}{\kappa^{2}} \partial_z \dot{h}_{tz} + 5 \, \Alp{1} \partial_z^2 \dot{\mathcal{T}}_{yty} \\
      &\quad + 5 \, \Alp{1} \partial_z^2 \dot{\mathcal{T}}_{ztz} + 5 \, \Alp{1} \partial_z^2 \dot{\mathcal{T}}_{xtx} - 4 \, \Alp{1} \partial_z^{3} \dot{h}_{tz} - \tfrac{1}{2 \, \kappa^2} \partial_t^2 h_{yy} - \tfrac{1}{2 \, \kappa^2} \partial_t^2 h_{zz} \\
      &\quad - 5 \, \Alp{1} \partial_t^{2} \partial_z \mathcal{T}_{xxz} - 5 \, \Alp{1} \partial_t^{2} \partial_z \mathcal{T}_{yyz} - 5 \, \Alp{1} \partial_t^{2} \partial_z \mathcal{T}_{ttz} - 2 \, \Alp{1} \partial_t^{2} \partial_z^{2} h_{tt} \\
      &\quad + 4 \, \Alp{1} \partial_t^{2} \partial_z^{2} h_{xx} + 4 \, \Alp{1} \partial_t^{2} \partial_z^{2} h_{yy} + 2 \, \Alp{1} \partial_t^{2} \partial_z^{2} h_{zz} - 5 \, \Alp{1} \partial_t^{3} \mathcal{T}_{yty} \\
      &\quad - 5 \, \Alp{1} \partial_t^{3} \mathcal{T}_{ztz} - 5 \, \Alp{1} \partial_t^{3} \mathcal{T}_{xtx} + 4 \, \Alp{1} \partial_t^{3} \partial_z h_{tz} - 2 \, \Alp{1} \partial_t^{4} h_{yy} - 2 \, \Alp{1} \partial_t^{4} h_{zz} \\
  %   2 \, \Alp{1} \, \partial_t^{4} h_{yy} &= \tfrac{B_{0}^2}{8} h_{tt} - \tfrac{1}{8} B_{0}^2 h_{xx} + \tfrac{3 \, B_{0}^2}{8} h_{yy} + \tfrac{B_{0}^2}{8} h_{zz} + \tfrac{B_{0}}{2} \partial_z A_{y} - \tfrac{1}{2 \, \kappa^2} \partial_z^2 h_{tt} + \tfrac{1}{2 \, \kappa^2} \partial_z^2 h_{xx} \\
  % &\quad + 5 \, \Alp{1} \partial_z^{3} \mathcal{T}_{xxz} + 5 \, \Alp{1} \partial_z^{3} \mathcal{T}_{yyz} + 5 \, \Alp{1} \partial_z^{3} \mathcal{T}_{ttz} + 2 \, \Alp{1} \partial_z^4 h_{tt} - 2 \, \Alp{1} \partial_z^4 h_{xx} - 2 \, \Alp{1} \partial_z^4 h_{yy} + \tfrac{1}{\kappa^{2}} \partial_z \dot{h}_{tz} + 5 \, \Alp{1} \partial_z^2 \dot{\mathcal{T}}_{yty} \\
  % &\quad + 5 \, \Alp{1} \partial_z^2 \dot{\mathcal{T}}_{ztz} + 5 \, \Alp{1} \partial_z^2 \dot{\mathcal{T}}_{xtx} - 4 \, \Alp{1} \partial_z^{3} \dot{h}_{tz} - \tfrac{1}{2 \, \kappa^2} \partial_t^2 h_{xx} - \tfrac{1}{2 \, \kappa^2} \partial_t^2 h_{zz} - 5 \, \Alp{1} \partial_t^{2} \partial_z \mathcal{T}_{xxz} - 5 \, \Alp{1} \partial_t^{2} \partial_z \mathcal{T}_{yyz} \\
  % &\quad - 5 \, \Alp{1} \partial_t^{2} \partial_z \mathcal{T}_{ttz} - 2 \, \Alp{1} \partial_t^{2} \partial_z^{2} h_{tt} + 4 \, \Alp{1} \partial_t^{2} \partial_z^{2} h_{xx} + 4 \, \Alp{1} \partial_t^{2} \partial_z^{2} h_{yy} + 2 \, \Alp{1} \partial_t^{2} \partial_z^{2} h_{zz} - 5 \, \Alp{1} \partial_t^{3} \mathcal{T}_{yty} - 5 \, \Alp{1} \partial_t^{3} \mathcal{T}_{ztz} \\
  % &\quad - 5 \, \Alp{1} \partial_t^{3} \mathcal{T}_{xtx} + 4 \, \Alp{1} \partial_t^{3} \partial_z h_{tz} - 2 \, \Alp{1} \partial_t^{4} h_{xx} - 2 \, \Alp{1} \partial_t^{4} h_{zz} \\
  %   2 \, \Alp{1} \, \partial_t^{4} h_{zz} &= \tfrac{B_{0}^2}{8} h_{tt} - \tfrac{1}{8} B_{0}^2 h_{xx} + \tfrac{B_{0}^2}{8} h_{yy} + \tfrac{3 \, B_{0}^2}{8} h_{zz} + \tfrac{B_{0}}{2} \partial_z A_{y} - \tfrac{1}{2 \, \kappa^2} \partial_t^2 h_{xx} - \tfrac{1}{2 \, \kappa^2} \partial_t^2 h_{yy} \\
  % &\quad - 5 \, \Alp{1} \partial_t^{2} \partial_z \mathcal{T}_{xxz} - 5 \, \Alp{1} \partial_t^{2} \partial_z \mathcal{T}_{yyz} - 5 \, \Alp{1} \partial_t^{2} \partial_z \mathcal{T}_{ttz} - 2 \, \Alp{1} \partial_t^{2} \partial_z^{2} h_{tt} + 2 \, \Alp{1} \partial_t^{2} \partial_z^{2} h_{xx} + 2 \, \Alp{1} \partial_t^{2} \partial_z^{2} h_{yy} \\
  % &\quad - 5 \, \Alp{1} \partial_t^{3} \mathcal{T}_{yty} - 5 \, \Alp{1} \partial_t^{3} \mathcal{T}_{ztz} - 5 \, \Alp{1} \partial_t^{3} \mathcal{T}_{xtx} + 4 \, \Alp{1} \partial_t^{3} \partial_z h_{tz} - 2 \, \Alp{1} \partial_t^{4} h_{xx} - 2 \, \Alp{1} \partial_t^{4} h_{yy}
  \end{aligned} \\[1ex]
  % \begin{aligned}
  %   0 &= 4 \, \Bet{1} \mathcal{T}_{ttx} + 2 \, \Bet{2} \mathcal{T}_{ttx} + 2 \, \Bet{3} \mathcal{T}_{ttx} + \tfrac{1}{2 \, \kappa^2} \mathcal{T}_{ttx} -2 \, \Bet{3} \mathcal{T}_{yxy} + \tfrac{3}{\kappa^2} \mathcal{T}_{yxy} -2 \, \Bet{3} \mathcal{T}_{zxz} + \tfrac{3}{\kappa^2} \mathcal{T}_{zxz} \\
  %   0 &= 4 \, \Bet{1} \mathcal{T}_{tty} + 2 \, \Bet{2} \mathcal{T}_{tty} + 2 \, \Bet{3} \mathcal{T}_{tty} + \tfrac{1}{2 \, \kappa^2} \mathcal{T}_{tty} -2 \, \Bet{3} \mathcal{T}_{zyz} + \tfrac{3}{\kappa^2} \mathcal{T}_{zyz} + 2 \, \Bet{3} \mathcal{T}_{xxy} -\tfrac{3}{\kappa^2} \mathcal{T}_{xxy} \\
  %   0 &= 4 \, \Bet{1} \mathcal{T}_{xyz} + \tfrac{2}{\kappa^2} \mathcal{T}_{xyz} + 2 \, \Bet{2} \mathcal{T}_{yxz} + \tfrac{3}{2 \, \kappa^2} \mathcal{T}_{yxz} -2 \, \Bet{2} \mathcal{T}_{zxy} -\tfrac{3}{2 \, \kappa^2} \mathcal{T}_{zxy} \\
  %   0 &= -4 \, \Bet{1} \mathcal{T}_{ytx} -\tfrac{2}{\kappa^2} \mathcal{T}_{ytx} -2 \, \Bet{2} \mathcal{T}_{txy} -\tfrac{3}{2 \, \kappa^2} \mathcal{T}_{txy} -2 \, \Bet{2} \mathcal{T}_{xty} -\tfrac{3}{2 \, \kappa^2} \mathcal{T}_{xty} \\
  %   0 &= -4 \, \Bet{1} \mathcal{T}_{ytz} -\tfrac{2}{\kappa^2} \mathcal{T}_{ytz} -2 \, \Bet{2} \mathcal{T}_{zty} -\tfrac{3}{2 \, \kappa^2} \mathcal{T}_{zty} + 2 \, \Bet{2} \mathcal{T}_{tyz} + \tfrac{3}{2 \, \kappa^2} \mathcal{T}_{tyz} \\
  %   0 &= -2 \, \Bet{3} \mathcal{T}_{ttx} + \tfrac{3}{\kappa^2} \mathcal{T}_{ttx} + 4 \, \Bet{1} \mathcal{T}_{yxy} + 2 \, \Bet{2} \mathcal{T}_{yxy} + 2 \, \Bet{3} \mathcal{T}_{yxy} + \tfrac{1}{2 \, \kappa^2} \mathcal{T}_{yxy} + 2 \, \Bet{3} \mathcal{T}_{zxz} -\tfrac{3}{\kappa^2} \mathcal{T}_{zxz} \\
  %   0 &= 2 \, \Bet{2} \mathcal{T}_{xyz} + \tfrac{3}{2 \, \kappa^2} \mathcal{T}_{xyz} + 4 \, \Bet{1} \mathcal{T}_{yxz} + \tfrac{2}{\kappa^2} \mathcal{T}_{yxz} + 2 \, \Bet{2} \mathcal{T}_{zxy} + \tfrac{3}{2 \, \kappa^2} \mathcal{T}_{zxy} \\
  %   0 &= -4 \, \Bet{1} \mathcal{T}_{ztx} -\tfrac{2}{\kappa^2} \mathcal{T}_{ztx} -2 \, \Bet{2} \mathcal{T}_{txz} -\tfrac{3}{2 \, \kappa^2} \mathcal{T}_{txz} -2 \, \Bet{2} \mathcal{T}_{xtz} -\tfrac{3}{2 \, \kappa^2} \mathcal{T}_{xtz} \\
  %   0 &= -2 \, \Bet{2} \mathcal{T}_{ytz} -\tfrac{3}{2 \, \kappa^2} \mathcal{T}_{ytz} -4 \, \Bet{1} \mathcal{T}_{zty} -\tfrac{2}{\kappa^2} \mathcal{T}_{zty} -2 \, \Bet{2} \mathcal{T}_{tyz} -\tfrac{3}{2 \, \kappa^2} \mathcal{T}_{tyz} \\
  %   0 &= -2 \, \Bet{2} \mathcal{T}_{xyz} -\tfrac{3}{2 \, \kappa^2} \mathcal{T}_{xyz} + 2 \, \Bet{2} \mathcal{T}_{yxz} + \tfrac{3}{2 \, \kappa^2} \mathcal{T}_{yxz} + 4 \, \Bet{1} \mathcal{T}_{zxy} + \tfrac{2}{\kappa^2} \mathcal{T}_{zxy} \\
  %   0 &= -2 \, \Bet{3} \mathcal{T}_{ttx} + \tfrac{3}{\kappa^2} \mathcal{T}_{ttx} + 2 \, \Bet{3} \mathcal{T}_{yxy} -\tfrac{3}{\kappa^2} \mathcal{T}_{yxy} + 4 \, \Bet{1} \mathcal{T}_{zxz} + 2 \, \Bet{2} \mathcal{T}_{zxz} + 2 \, \Bet{3} \mathcal{T}_{zxz} + \tfrac{1}{2 \, \kappa^2} \mathcal{T}_{zxz} \\
  %   0 &= -2 \, \Bet{3} \mathcal{T}_{tty} + \tfrac{3}{\kappa^2} \mathcal{T}_{tty} + 4 \, \Bet{1} \mathcal{T}_{zyz} + 2 \, \Bet{2} \mathcal{T}_{zyz} + 2 \, \Bet{3} \mathcal{T}_{zyz} + \tfrac{1}{2 \, \kappa^2} \mathcal{T}_{zyz} -2 \, \Bet{3} \mathcal{T}_{xxy} + \tfrac{3}{\kappa^2} \mathcal{T}_{xxy} \\
  %   0 &= -2 \, \Bet{2} \mathcal{T}_{ytx} -\tfrac{3}{2 \, \kappa^2} \mathcal{T}_{ytx} -4 \, \Bet{1} \mathcal{T}_{txy} -\tfrac{2}{\kappa^2} \mathcal{T}_{txy} + 2 \, \Bet{2} \mathcal{T}_{xty} + \tfrac{3}{2 \, \kappa^2} \mathcal{T}_{xty} \\
  %   0 &= -2 \, \Bet{2} \mathcal{T}_{ztx} -\tfrac{3}{2 \, \kappa^2} \mathcal{T}_{ztx} -4 \, \Bet{1} \mathcal{T}_{txz} -\tfrac{2}{\kappa^2} \mathcal{T}_{txz} + 2 \, \Bet{2} \mathcal{T}_{xtz} + \tfrac{3}{2 \, \kappa^2} \mathcal{T}_{xtz} \\
  %   0 &= 2 \, \Bet{2} \mathcal{T}_{ytz} + \tfrac{3}{2 \, \kappa^2} \mathcal{T}_{ytz} -2 \, \Bet{2} \mathcal{T}_{zty} -\tfrac{3}{2 \, \kappa^2} \mathcal{T}_{zty} -4 \, \Bet{1} \mathcal{T}_{tyz} -\tfrac{2}{\kappa^2} \mathcal{T}_{tyz} \\
  %   0 &= -2 \, \Bet{2} \mathcal{T}_{ytx} -\tfrac{3}{2 \, \kappa^2} \mathcal{T}_{ytx} + 2 \, \Bet{2} \mathcal{T}_{txy} + \tfrac{3}{2 \, \kappa^2} \mathcal{T}_{txy} -4 \, \Bet{1} \mathcal{T}_{xty} -\tfrac{2}{\kappa^2} \mathcal{T}_{xty} \\
  %   0 &= -2 \, \Bet{2} \mathcal{T}_{ztx} -\tfrac{3}{2 \, \kappa^2} \mathcal{T}_{ztx} + 2 \, \Bet{2} \mathcal{T}_{txz} + \tfrac{3}{2 \, \kappa^2} \mathcal{T}_{txz} -4 \, \Bet{1} \mathcal{T}_{xtz} -\tfrac{2}{\kappa^2} \mathcal{T}_{xtz} \\
  %   0 &= 2 \, \Bet{3} \mathcal{T}_{tty} -\tfrac{3}{\kappa^2} \mathcal{T}_{tty} -2 \, \Bet{3} \mathcal{T}_{zyz} + \tfrac{3}{\kappa^2} \mathcal{T}_{zyz} + 4 \, \Bet{1} \mathcal{T}_{xxy} + 2 \, \Bet{2} \mathcal{T}_{xxy} + 2 \, \Bet{3} \mathcal{T}_{xxy} + \tfrac{1}{2 \, \kappa^2} \mathcal{T}_{xxy}
  % \end{aligned} \\[1ex]
  % \begin{aligned}
  %   0 &= 4 \, \Bet{1} \mathcal{T}_{xxz} + 2 \, \Bet{2} \mathcal{T}_{xxz} + 2 \, \Bet{3} \mathcal{T}_{xxz} + \tfrac{1}{2 \, \kappa^2} \mathcal{T}_{xxz} + 2 \, \Bet{3} \mathcal{T}_{yyz} - \tfrac{3}{\kappa^2} \mathcal{T}_{yyz} + 2 \, \Bet{3} \mathcal{T}_{ttz} - \tfrac{3}{\kappa^2} \mathcal{T}_{ttz} - \tfrac{25 \, \Alp{1}}{2} \partial_z^2 \mathcal{T}_{xxz} \\
  % &\quad - \tfrac{25 \, \Alp{1}}{2} \partial_z^2 \mathcal{T}_{yyz} - \tfrac{25 \, \Alp{1}}{2} \partial_z^2 \mathcal{T}_{ttz} - 5 \, \Alp{1} \partial_z^{3} h_{tt} + 5 \, \Alp{1} \partial_z^{3} h_{xx} + 5 \, \Alp{1} \partial_z^{3} h_{yy} - \tfrac{25 \, \Alp{1}}{2} \partial_z \dot{\mathcal{T}}_{yty} - \tfrac{25 \, \Alp{1}}{2} \partial_z \dot{\mathcal{T}}_{ztz} \\
  % &\quad - \tfrac{25 \, \Alp{1}}{2} \partial_z \dot{\mathcal{T}}_{xtx} + 10 \, \Alp{1} \partial_z^2 \dot{h}_{tz} - 5 \, \Alp{1} \partial_t^{2} \partial_z h_{xx} - 5 \, \Alp{1} \partial_t^{2} \partial_z h_{yy} - 5 \, \Alp{1} \partial_t^{2} \partial_z h_{zz}
  % \end{aligned} \\[1ex]
  % \begin{aligned}
  %   -\tfrac{25 \, \Alp{1}}{2} \, \partial_t^{2} \mathcal{T}_{yty} &= 4 \, \Bet{1} \mathcal{T}_{yty} + 2 \, \Bet{2} \mathcal{T}_{yty} + 2 \, \Bet{3} \mathcal{T}_{yty} + \tfrac{1}{2 \, \kappa^2} \mathcal{T}_{yty} + 2 \, \Bet{3} \mathcal{T}_{ztz} - \tfrac{3}{\kappa^2} \mathcal{T}_{ztz} + 2 \, \Bet{3} \mathcal{T}_{xtx} - \tfrac{3}{\kappa^2} \mathcal{T}_{xtx} \\
  % &\quad + \tfrac{25 \, \Alp{1}}{2} \partial_z \dot{\mathcal{T}}_{xxz} + \tfrac{25 \, \Alp{1}}{2} \partial_z \dot{\mathcal{T}}_{yyz} + \tfrac{25 \, \Alp{1}}{2} \partial_z \dot{\mathcal{T}}_{ttz} + 5 \, \Alp{1} \partial_z^2 \dot{h}_{tt} - 5 \, \Alp{1} \partial_z^2 \dot{h}_{xx} - 5 \, \Alp{1} \partial_z^2 \dot{h}_{yy} + \tfrac{25 \, \Alp{1}}{2} \partial_t^2 \mathcal{T}_{ztz} \\
  % &\quad + \tfrac{25 \, \Alp{1}}{2} \partial_t^2 \mathcal{T}_{xtx} - 10 \, \Alp{1} \partial_t^{2} \partial_z h_{tz} + 5 \, \Alp{1} \partial_t^{3} h_{xx} + 5 \, \Alp{1} \partial_t^{3} h_{yy} + 5 \, \Alp{1} \partial_t^{3} h_{zz}
  % \end{aligned} \\[1ex]
  % \begin{aligned}
  %   0 &= 2 \, \Bet{3} \mathcal{T}_{xxz} - \tfrac{3}{\kappa^2} \mathcal{T}_{xxz} + 4 \, \Bet{1} \mathcal{T}_{yyz} + 2 \, \Bet{2} \mathcal{T}_{yyz} + 2 \, \Bet{3} \mathcal{T}_{yyz} + \tfrac{1}{2 \, \kappa^2} \mathcal{T}_{yyz} + 2 \, \Bet{3} \mathcal{T}_{ttz} - \tfrac{3}{\kappa^2} \mathcal{T}_{ttz} - \tfrac{25 \, \Alp{1}}{2} \partial_z^2 \mathcal{T}_{xxz} \\
  % &\quad - \tfrac{25 \, \Alp{1}}{2} \partial_z^2 \mathcal{T}_{yyz} - \tfrac{25 \, \Alp{1}}{2} \partial_z^2 \mathcal{T}_{ttz} - 5 \, \Alp{1} \partial_z^{3} h_{tt} + 5 \, \Alp{1} \partial_z^{3} h_{xx} + 5 \, \Alp{1} \partial_z^{3} h_{yy} - \tfrac{25 \, \Alp{1}}{2} \partial_z \dot{\mathcal{T}}_{yty} - \tfrac{25 \, \Alp{1}}{2} \partial_z \dot{\mathcal{T}}_{ztz} \\
  % &\quad - \tfrac{25 \, \Alp{1}}{2} \partial_z \dot{\mathcal{T}}_{xtx} + 10 \, \Alp{1} \partial_z^2 \dot{h}_{tz} - 5 \, \Alp{1} \partial_t^{2} \partial_z h_{xx} - 5 \, \Alp{1} \partial_t^{2} \partial_z h_{yy} - 5 \, \Alp{1} \partial_t^{2} \partial_z h_{zz}
  % \end{aligned} \\[1ex]
  % \begin{aligned}
  %   0 &= 2 \, \Bet{3} \mathcal{T}_{xxz} - \tfrac{3}{\kappa^2} \mathcal{T}_{xxz} + 2 \, \Bet{3} \mathcal{T}_{yyz} - \tfrac{3}{\kappa^2} \mathcal{T}_{yyz} + 4 \, \Bet{1} \mathcal{T}_{ttz} + 2 \, \Bet{2} \mathcal{T}_{ttz} + 2 \, \Bet{3} \mathcal{T}_{ttz} + \tfrac{1}{2 \, \kappa^2} \mathcal{T}_{ttz} - \tfrac{25 \, \Alp{1}}{2} \partial_z^2 \mathcal{T}_{xxz} \\
  % &\quad - \tfrac{25 \, \Alp{1}}{2} \partial_z^2 \mathcal{T}_{yyz} - \tfrac{25 \, \Alp{1}}{2} \partial_z^2 \mathcal{T}_{ttz} - 5 \, \Alp{1} \partial_z^{3} h_{tt} + 5 \, \Alp{1} \partial_z^{3} h_{xx} + 5 \, \Alp{1} \partial_z^{3} h_{yy} - \tfrac{25 \, \Alp{1}}{2} \partial_z \dot{\mathcal{T}}_{yty} - \tfrac{25 \, \Alp{1}}{2} \partial_z \dot{\mathcal{T}}_{ztz} \\
  % &\quad - \tfrac{25 \, \Alp{1}}{2} \partial_z \dot{\mathcal{T}}_{xtx} + 10 \, \Alp{1} \partial_z^2 \dot{h}_{tz} - 5 \, \Alp{1} \partial_t^{2} \partial_z h_{xx} - 5 \, \Alp{1} \partial_t^{2} \partial_z h_{yy} - 5 \, \Alp{1} \partial_t^{2} \partial_z h_{zz}
  % \end{aligned} \\[1ex]
  % \begin{aligned}
  %   -\tfrac{25 \, \Alp{1}}{2} \, \partial_t^{2} \mathcal{T}_{ztz} &= 2 \, \Bet{3} \mathcal{T}_{yty} - \tfrac{3}{\kappa^2} \mathcal{T}_{yty} + 4 \, \Bet{1} \mathcal{T}_{ztz} + 2 \, \Bet{2} \mathcal{T}_{ztz} + 2 \, \Bet{3} \mathcal{T}_{ztz} + \tfrac{1}{2 \, \kappa^2} \mathcal{T}_{ztz} + 2 \, \Bet{3} \mathcal{T}_{xtx} - \tfrac{3}{\kappa^2} \mathcal{T}_{xtx} \\
  % &\quad + \tfrac{25 \, \Alp{1}}{2} \partial_z \dot{\mathcal{T}}_{xxz} + \tfrac{25 \, \Alp{1}}{2} \partial_z \dot{\mathcal{T}}_{yyz} + \tfrac{25 \, \Alp{1}}{2} \partial_z \dot{\mathcal{T}}_{ttz} + 5 \, \Alp{1} \partial_z^2 \dot{h}_{tt} - 5 \, \Alp{1} \partial_z^2 \dot{h}_{xx} - 5 \, \Alp{1} \partial_z^2 \dot{h}_{yy} + \tfrac{25 \, \Alp{1}}{2} \partial_t^2 \mathcal{T}_{yty} \\
  % &\quad + \tfrac{25 \, \Alp{1}}{2} \partial_t^2 \mathcal{T}_{xtx} - 10 \, \Alp{1} \partial_t^{2} \partial_z h_{tz} + 5 \, \Alp{1} \partial_t^{3} h_{xx} + 5 \, \Alp{1} \partial_t^{3} h_{yy} + 5 \, \Alp{1} \partial_t^{3} h_{zz}
  % \end{aligned} \\[1ex]
  % \begin{aligned}
  %   -\tfrac{25 \, \Alp{1}}{2} \, \partial_t^{2} \mathcal{T}_{xtx} &= 2 \, \Bet{3} \mathcal{T}_{yty} - \tfrac{3}{\kappa^2} \mathcal{T}_{yty} + 2 \, \Bet{3} \mathcal{T}_{ztz} - \tfrac{3}{\kappa^2} \mathcal{T}_{ztz} + 4 \, \Bet{1} \mathcal{T}_{xtx} + 2 \, \Bet{2} \mathcal{T}_{xtx} + 2 \, \Bet{3} \mathcal{T}_{xtx} + \tfrac{1}{2 \, \kappa^2} \mathcal{T}_{xtx} \\
  % &\quad + \tfrac{25 \, \Alp{1}}{2} \partial_z \dot{\mathcal{T}}_{xxz} + \tfrac{25 \, \Alp{1}}{2} \partial_z \dot{\mathcal{T}}_{yyz} + \tfrac{25 \, \Alp{1}}{2} \partial_z \dot{\mathcal{T}}_{ttz} + 5 \, \Alp{1} \partial_z^2 \dot{h}_{tt} - 5 \, \Alp{1} \partial_z^2 \dot{h}_{xx} - 5 \, \Alp{1} \partial_z^2 \dot{h}_{yy} + \tfrac{25 \, \Alp{1}}{2} \partial_t^2 \mathcal{T}_{yty} \\
  % &\quad + \tfrac{25 \, \Alp{1}}{2} \partial_t^2 \mathcal{T}_{ztz} - 10 \, \Alp{1} \partial_t^{2} \partial_z h_{tz} + 5 \, \Alp{1} \partial_t^{3} h_{xx} + 5 \, \Alp{1} \partial_t^{3} h_{yy} + 5 \, \Alp{1} \partial_t^{3} h_{zz}
  % \end{aligned} \\[1ex]
  % \begin{aligned}
  %   \mathscr{H} &\supset -\tfrac{1}{16} B_{0}^2 h_{tt}^2 + \tfrac{B_{0}^2}{4} h_{tx}^2 - \tfrac{1}{4} B_{0}^2 h_{ty}^2 - \tfrac{1}{4} B_{0}^2 h_{tz}^2 - \tfrac{1}{8} B_{0}^2 h_{tt} \, h_{xx} - \tfrac{1}{16} B_{0}^2 h_{xx}^2 + \tfrac{B_{0}^2}{4} h_{xy}^2 + \tfrac{B_{0}^2}{4} h_{xz}^2 \\
  % &\quad + \tfrac{B_{0}^2}{8} h_{tt} \, h_{yy} - \tfrac{1}{8} B_{0}^2 h_{xx} \, h_{yy} + \tfrac{3 \, B_{0}^2}{16} h_{yy}^2 + \tfrac{B_{0}^2}{4} h_{yz}^2 + \tfrac{B_{0}^2}{8} h_{tt} \, h_{zz} - \tfrac{1}{8} B_{0}^2 h_{xx} \, h_{zz} + \tfrac{B_{0}^2}{8} h_{yy} \, h_{zz} \\
  % &\quad + \tfrac{3 \, B_{0}^2}{16} h_{zz}^2 - \tfrac{1}{2} \partial_z A_{t}^2 + \tfrac{1}{2} \partial_z A_{x}^2 + \tfrac{1}{2} \partial_z A_{y}^2 + \tfrac{1}{2} \partial_z A_{z}^2 - \tfrac{1}{2 \, \kappa^2} \partial_z h_{tx}^2 - \tfrac{1}{2 \, \kappa^2} \partial_z h_{ty}^2 \\
  % &\quad + \tfrac{1}{2 \, \kappa^2} \partial_z h_{tt} \, \partial_z h_{xx} + \tfrac{1}{2 \, \kappa^2} \partial_z h_{xy}^2 + \tfrac{1}{2 \, \kappa^2} \partial_z h_{tt} \, \partial_z h_{yy} - \tfrac{1}{2 \, \kappa^2} \partial_z h_{xx} \, \partial_z h_{yy} - \Alp{1} \partial_z^2 h_{tt}^2 \\
  % &\quad + 2 \, \Alp{1} \partial_z^2 h_{tt} \, \partial_z^2 h_{xx} - \Alp{1} \partial_z^2 h_{xx}^2 + 2 \, \Alp{1} \partial_z^2 h_{tt} \, \partial_z^2 h_{yy} - 2 \, \Alp{1} \partial_z^2 h_{xx} \, \partial_z^2 h_{yy} - \Alp{1} \partial_z^2 h_{yy}^2 - \tfrac{1}{2} \dot{A}_{t}^2 + \tfrac{1}{2} \dot{A}_{x}^2 \\
  % &\quad + \tfrac{1}{2} \dot{A}_{y}^2 + \tfrac{1}{2} \dot{A}_{z}^2 + \tfrac{1}{2 \, \kappa^2} \dot{h}_{xy}^2 + \tfrac{1}{2 \, \kappa^2} \dot{h}_{xz}^2 - \tfrac{1}{2 \, \kappa^2} \dot{h}_{xx} \, \dot{h}_{yy} + \tfrac{1}{2 \, \kappa^2} \dot{h}_{yz}^2 - \tfrac{1}{2 \, \kappa^2} \dot{h}_{xx} \, \dot{h}_{zz} \\
  % &\quad - \tfrac{1}{2 \, \kappa^2} \dot{h}_{yy} \, \dot{h}_{zz} - \tfrac{2}{\kappa^2} h_{xz} \, \partial_t \partial_z h_{tx} - \tfrac{2}{\kappa^2} h_{yz} \, \partial_t \partial_z h_{ty} + \tfrac{1}{\kappa^{2}} h_{tt} \, \partial_t \partial_z h_{tz} \\
  % &\quad + \tfrac{1}{\kappa^{2}} h_{xx} \, \partial_t \partial_z h_{tz} + \tfrac{1}{\kappa^{2}} h_{yy} \, \partial_t \partial_z h_{tz} - \tfrac{1}{\kappa^{2}} h_{zz} \, \partial_t \partial_z h_{tz} + 4 \, \Alp{1} \partial_z^2 h_{tt} \, \partial_t \partial_z h_{tz} \\
  % &\quad - 4 \, \Alp{1} \partial_z^2 h_{xx} \, \partial_t \partial_z h_{tz} - 4 \, \Alp{1} \partial_z^2 h_{yy} \, \partial_t \partial_z h_{tz} - 4 \, \Alp{1} \partial_t \partial_z h_{tz}^2 + \tfrac{2}{\kappa^2} h_{tz} \, \partial_t \partial_z h_{xx} \\
  % &\quad - 2 \, \Alp{1} \partial_t \partial_z h_{tt} \, \partial_t \partial_z h_{xx} + 2 \, \Alp{1} \partial_t \partial_z h_{xx}^2 - \tfrac{2}{\kappa^2} h_{tx} \, \partial_t \partial_z h_{xz} + \tfrac{2}{\kappa^2} h_{tz} \, \partial_t \partial_z h_{yy} \\
  % &\quad - 2 \, \Alp{1} \partial_t \partial_z h_{tt} \, \partial_t \partial_z h_{yy} + 4 \, \Alp{1} \partial_t \partial_z h_{xx} \, \partial_t \partial_z h_{yy} + 2 \, \Alp{1} \partial_t \partial_z h_{yy}^2 - \tfrac{2}{\kappa^2} h_{ty} \, \partial_t \partial_z h_{yz} \\
  % &\quad - 2 \, \Alp{1} \partial_t \partial_z h_{tt} \, \partial_t \partial_z h_{zz} + 2 \, \Alp{1} \partial_t \partial_z h_{xx} \, \partial_t \partial_z h_{zz} + 2 \, \Alp{1} \partial_t \partial_z h_{yy} \, \partial_t \partial_z h_{zz} - 2 \, \Alp{1} \dot{h}_{xx} \, \partial_t \partial_z^{2} h_{tt} \\
  % &\quad - 2 \, \Alp{1} \dot{h}_{yy} \, \partial_t \partial_z^{2} h_{tt} - 2 \, \Alp{1} \dot{h}_{zz} \, \partial_t \partial_z^{2} h_{tt} + 2 \, \Alp{1} \dot{h}_{xx} \, \partial_t \partial_z^{2} h_{xx} + 2 \, \Alp{1} \dot{h}_{yy} \, \partial_t \partial_z^{2} h_{xx} \\
  % &\quad + 2 \, \Alp{1} \dot{h}_{zz} \, \partial_t \partial_z^{2} h_{xx} + 2 \, \Alp{1} \dot{h}_{xx} \, \partial_t \partial_z^{2} h_{yy} + 2 \, \Alp{1} \dot{h}_{yy} \, \partial_t \partial_z^{2} h_{yy} + 2 \, \Alp{1} \dot{h}_{zz} \, \partial_t \partial_z^{2} h_{yy} - \Alp{1} \partial_t^{2} h_{xx}^2 \\
  % &\quad - \Alp{1} \partial_t^{2} h_{yy}^2 - \Alp{1} \partial_t^{2} h_{zz}^2
  % \end{aligned}

\end{gathered}
$}
\medskip


% --- E.5  Ricci--EM nonminimal coupling -----------------------------
\subsection{Riemann--Cartan electromagnetic coupling}%
\label{SymbolicNonminimal}
A single new vertex
$\delta_{1}\MAGFt{_{\mu\nu}}\FieldA{^{\mu\nu}}$ is added to the
Einstein--Cartan terms. The symbolic output
comprises \NonminimalEomCount{} equations of motion
(\NonminimalConstraintCount{} constraints) with \NonminimalTermCount{}
RHS terms.

\medskip
\noindent\adjustbox{max width=\linewidth}{$\displaystyle
\setlength{\jot}{1pt}%
\begin{gathered}
  \begin{aligned}
    \mathscr{L} &= \tfrac{1}{\kappa^2} \tilde{\mathcal{R}} + \alpha_{1} \tensor{T}{_a_b_c} \tensor{g}{^a^d} \tensor{g}{^b^e} \tensor{g}{^c^f} \tensor{T}{_d_e_f} + \alpha_{2} \tensor{T}{_a_b_c} \tensor{g}{^b^d} \tensor{g}{^a^e} \tensor{g}{^c^f} \tensor{T}{_d_e_f} + \alpha_{3} \tensor{T}{^b_b_a} \tensor{g}{^a^d} \tensor{T}{^c_c_d} \\
  &\quad + \delta_{1} \tensor{\tilde{R}}{_a_b} \tensor{g}{^a^c} \tensor{g}{^b^d} \tensor{\mathcal{F}}{_c_d} - \tfrac{1}{4} \tensor{\mathcal{F}}{_a_b} \tensor{g}{^a^c} \tensor{g}{^b^d} \tensor{\mathcal{F}}{_c_d}
  \end{aligned} \\[1ex]
  % \begin{aligned}
  %   \partial_t^{2} A_{t} &= -B_{0} \partial_z h_{ty} - \partial_z^2 A_{t} - \tfrac{3 \, \delta_{1}}{2} \partial_z^2 \mathcal{T}_{yty} + \tfrac{\delta_{1}}{2} \partial_z^2 \mathcal{T}_{ztz} - \tfrac{3 \, \delta_{1}}{2} \partial_z^2 \mathcal{T}_{xtx} - \tfrac{3 \, \delta_{1}}{2} \partial_z \dot{\mathcal{T}}_{xxz} \\
  % &\quad - \tfrac{3 \, \delta_{1}}{2} \partial_z \dot{\mathcal{T}}_{yyz} + \tfrac{\delta_{1}}{2} \partial_z \dot{\mathcal{T}}_{ttz}
  % \end{aligned} \\[1ex]
  \begin{aligned}
    \partial_t^{2} A_{x} &= B_{0} \partial_z h_{xy} + \partial_z^2 A_{x} - \tfrac{3 \, \delta_{1}}{2} \partial_z^2 \mathcal{T}_{ttx} + \tfrac{3 \, \delta_{1}}{2} \partial_z^2 \mathcal{T}_{yxy} - \tfrac{1}{2} \delta_{1} \partial_z^2 \mathcal{T}_{zxz} - 2 \, \delta_{1} \partial_z \dot{\mathcal{T}}_{ztx} - 2 \, \delta_{1} \partial_z \dot{\mathcal{T}}_{txz} \\
  &\quad - \tfrac{1}{2} \delta_{1} \partial_t^2 \mathcal{T}_{ttx} - \tfrac{3 \, \delta_{1}}{2} \partial_t^2 \mathcal{T}_{yxy} - \tfrac{3 \, \delta_{1}}{2} \partial_t^2 \mathcal{T}_{zxz}
  \end{aligned} \\[1ex]
  % \begin{aligned}
  %   \partial_t^{2} A_{y} &= \tfrac{B_{0}}{2} \partial_z h_{tt} - \tfrac{1}{2} B_{0} \partial_z h_{xx} + \tfrac{B_{0}}{2} \partial_z h_{yy} + \tfrac{B_{0}}{2} \partial_z h_{zz} + \partial_z^2 A_{y} - \tfrac{3 \, \delta_{1}}{2} \partial_z^2 \mathcal{T}_{tty} - \tfrac{1}{2} \delta_{1} \partial_z^2 \mathcal{T}_{zyz} \\
  % &\quad - \tfrac{3 \, \delta_{1}}{2} \partial_z^2 \mathcal{T}_{xxy} - B_{0} \partial_t h_{tz} - 2 \, \delta_{1} \partial_z \dot{\mathcal{T}}_{zty} - 2 \, \delta_{1} \partial_z \dot{\mathcal{T}}_{tyz} - \tfrac{1}{2} \delta_{1} \partial_t^2 \mathcal{T}_{tty} - \tfrac{3 \, \delta_{1}}{2} \partial_t^2 \mathcal{T}_{zyz} + \tfrac{3 \, \delta_{1}}{2} \partial_t^2 \mathcal{T}_{xxy}
  % \end{aligned} \\[1ex]
  % \begin{aligned}
  %   \partial_t^{2} A_{z} &= \partial_z^2 A_{z} + B_{0} \partial_t h_{ty} + \tfrac{3 \, \delta_{1}}{2} \partial_z \dot{\mathcal{T}}_{yty} - \tfrac{1}{2} \delta_{1} \partial_z \dot{\mathcal{T}}_{ztz} + \tfrac{3 \, \delta_{1}}{2} \partial_z \dot{\mathcal{T}}_{xtx} + \tfrac{3 \, \delta_{1}}{2} \partial_t^2 \mathcal{T}_{xxz} \\
  % &\quad + \tfrac{3 \, \delta_{1}}{2} \partial_t^2 \mathcal{T}_{yyz} - \tfrac{1}{2} \delta_{1} \partial_t^2 \mathcal{T}_{ttz}
  % \end{aligned} \\[1ex]
  % \begin{aligned}
  %   0 &= \tfrac{B_{0}^2}{8} h_{tt} + \tfrac{B_{0}^2}{8} h_{xx} - \tfrac{1}{8} B_{0}^2 h_{yy} - \tfrac{1}{8} B_{0}^2 h_{zz} - \tfrac{1}{2} B_{0} \partial_z A_{y} - \tfrac{1}{4} B_{0} \, \delta_{1} \partial_z \mathcal{T}_{tty} - \tfrac{3 \, B_{0} \, \delta_{1}}{4} \partial_z \mathcal{T}_{zyz} \\
  % &\quad + \tfrac{3 \, B_{0} \, \delta_{1}}{4} \partial_z \mathcal{T}_{xxy} + \tfrac{1}{2 \, \kappa^2} \partial_z^2 h_{xx} + \tfrac{1}{2 \, \kappa^2} \partial_z^2 h_{yy}
  % \end{aligned} \\[1ex]
  % \begin{aligned}
  %   0 &= -\tfrac{1}{2} B_{0}^2 h_{tx} + B_{0} \, \delta_{1} \partial_z \mathcal{T}_{txy} -B_{0} \, \delta_{1} \partial_z \mathcal{T}_{xty} -\tfrac{1}{\kappa^{2}} \partial_z^2 h_{tx} + \tfrac{1}{\kappa^{2}} \partial_z \dot{h}_{xz}
  % \end{aligned} \\[1ex]
  % \begin{aligned}
  %   0 &= \tfrac{B_{0}^2}{2} h_{ty} + B_{0} \partial_z A_{t} + \tfrac{B_{0} \, \delta_{1}}{2} \partial_z \mathcal{T}_{yty} + \tfrac{B_{0} \, \delta_{1}}{2} \partial_z \mathcal{T}_{ztz} + \tfrac{3 \, B_{0} \, \delta_{1}}{2} \partial_z \mathcal{T}_{xtx} - \tfrac{1}{\kappa^{2}} \partial_z^2 h_{ty} - B_{0} \partial_t A_{z} \\
  % &\quad + \tfrac{3 \, B_{0} \, \delta_{1}}{2} \partial_t \mathcal{T}_{xxz} + \tfrac{3 \, B_{0} \, \delta_{1}}{2} \partial_t \mathcal{T}_{yyz} - \tfrac{1}{2} B_{0} \, \delta_{1} \partial_t \mathcal{T}_{ttz} + \tfrac{1}{\kappa^{2}} \partial_z \dot{h}_{yz}
  % \end{aligned} \\[1ex]
  % \begin{aligned}
  %   0 &= \tfrac{B_{0}^2}{2} h_{tz} + B_{0} \partial_t A_{y} + \tfrac{B_{0} \, \delta_{1}}{2} \partial_t \mathcal{T}_{tty} + \tfrac{3 \, B_{0} \, \delta_{1}}{2} \partial_t \mathcal{T}_{zyz} -\tfrac{3 \, B_{0} \, \delta_{1}}{2} \partial_t \mathcal{T}_{xxy} -\tfrac{1}{\kappa^{2}} \partial_z \dot{h}_{xx} -\tfrac{1}{\kappa^{2}} \partial_z \dot{h}_{yy}
  % \end{aligned} \\[1ex]
  % \begin{aligned}
  %   0 &= \tfrac{B_{0}^2}{8} h_{tt} + \tfrac{B_{0}^2}{8} h_{xx} + \tfrac{B_{0}^2}{8} h_{yy} + \tfrac{B_{0}^2}{8} h_{zz} + \tfrac{B_{0}}{2} \partial_z A_{y} - \tfrac{3 \, B_{0} \, \delta_{1}}{4} \partial_z \mathcal{T}_{tty} + \tfrac{3 \, B_{0} \, \delta_{1}}{4} \partial_z \mathcal{T}_{zyz} + \tfrac{B_{0} \, \delta_{1}}{4} \partial_z \mathcal{T}_{xxy} \\
  % &\quad + \tfrac{1}{2 \, \kappa^2} \partial_z^2 h_{tt} - \tfrac{1}{2 \, \kappa^2} \partial_z^2 h_{yy} - \tfrac{1}{\kappa^{2}} \partial_z \dot{h}_{tz} + \tfrac{1}{2 \, \kappa^2} \partial_t^2 h_{yy} + \tfrac{1}{2 \, \kappa^2} \partial_t^2 h_{zz}
  % \end{aligned} \\[1ex]
  \begin{aligned}
    -\tfrac{1}{\kappa^{2}} \, \partial_t^{2} h_{xy} &= \tfrac{B_{0}^2}{2} h_{xy} + B_{0} \partial_z A_{x} - \tfrac{3 \, B_{0} \, \delta_{1}}{2} \partial_z \mathcal{T}_{ttx} + \tfrac{B_{0} \, \delta_{1}}{2} \partial_z \mathcal{T}_{yxy} + \tfrac{B_{0} \, \delta_{1}}{2} \partial_z \mathcal{T}_{zxz} - \tfrac{1}{\kappa^{2}} \partial_z^2 h_{xy} \\
  &\quad - B_{0} \, \delta_{1} \partial_t \mathcal{T}_{txz} + B_{0} \, \delta_{1} \partial_t \mathcal{T}_{xtz}
  \end{aligned} \\[1ex]
  % \begin{aligned}
  %   -\tfrac{1}{\kappa^{2}} \, \partial_t^{2} h_{xz} &= \tfrac{B_{0}^2}{2} h_{xz} + B_{0} \, \delta_{1} \partial_t \mathcal{T}_{txy} -B_{0} \, \delta_{1} \partial_t \mathcal{T}_{xty} -\tfrac{1}{\kappa^{2}} \partial_z \dot{h}_{tx}
  % \end{aligned} \\[1ex]
  % \begin{aligned}
  %   0 &= -\tfrac{1}{8} B_{0}^2 h_{tt} + \tfrac{B_{0}^2}{8} h_{xx} - \tfrac{3 \, B_{0}^2}{8} h_{yy} - \tfrac{1}{8} B_{0}^2 h_{zz} - \tfrac{1}{2} B_{0} \partial_z A_{y} + \tfrac{3 \, B_{0} \, \delta_{1}}{4} \partial_z \mathcal{T}_{tty} + \tfrac{B_{0} \, \delta_{1}}{4} \partial_z \mathcal{T}_{zyz} \\
  % &\quad + \tfrac{3 \, B_{0} \, \delta_{1}}{4} \partial_z \mathcal{T}_{xxy} + \tfrac{1}{2 \, \kappa^2} \partial_z^2 h_{tt} - \tfrac{1}{2 \, \kappa^2} \partial_z^2 h_{xx} - B_{0} \, \delta_{1} \partial_t \mathcal{T}_{ytz} + B_{0} \, \delta_{1} \partial_t \mathcal{T}_{tyz} - \tfrac{1}{\kappa^{2}} \partial_z \dot{h}_{tz} \\
  % &\quad + \tfrac{1}{2 \, \kappa^2} \partial_t^2 h_{xx} + \tfrac{1}{2 \, \kappa^2} \partial_t^2 h_{zz}
  % \end{aligned} \\[1ex]
  % \begin{aligned}
  %   -\tfrac{1}{\kappa^{2}} \, \partial_t^{2} h_{yz} &= \tfrac{B_{0}^2}{2} h_{yz} -B_{0} \, \delta_{1} \partial_t \mathcal{T}_{yty} + B_{0} \, \delta_{1} \partial_t \mathcal{T}_{ztz} -\tfrac{1}{\kappa^{2}} \partial_z \dot{h}_{ty}
  % \end{aligned} \\[1ex]
  % \begin{aligned}
  %   0 &= -\tfrac{1}{8} B_{0}^2 h_{tt} + \tfrac{B_{0}^2}{8} h_{xx} - \tfrac{1}{8} B_{0}^2 h_{yy} - \tfrac{3 \, B_{0}^2}{8} h_{zz} - \tfrac{1}{2} B_{0} \partial_z A_{y} + \tfrac{3 \, B_{0} \, \delta_{1}}{4} \partial_z \mathcal{T}_{tty} + \tfrac{B_{0} \, \delta_{1}}{4} \partial_z \mathcal{T}_{zyz} \\
  % &\quad + \tfrac{3 \, B_{0} \, \delta_{1}}{4} \partial_z \mathcal{T}_{xxy} + B_{0} \, \delta_{1} \partial_t \mathcal{T}_{zty} + B_{0} \, \delta_{1} \partial_t \mathcal{T}_{tyz} + \tfrac{1}{2 \, \kappa^2} \partial_t^2 h_{xx} + \tfrac{1}{2 \, \kappa^2} \partial_t^2 h_{yy}
  % \end{aligned} \\[1ex]
  % \begin{aligned}
  %   0 &= 4 \, \alpha_{1} \mathcal{T}_{ttx} + 2 \, \alpha_{2} \mathcal{T}_{ttx} + 2 \, \alpha_{3} \mathcal{T}_{ttx} + \tfrac{1}{2 \, \kappa^2} \mathcal{T}_{ttx} + \tfrac{3 \, B_{0} \, \delta_{1}}{2} \mathcal{T}_{xyz} - 2 \, \alpha_{3} \mathcal{T}_{yxy} + \tfrac{3}{\kappa^2} \mathcal{T}_{yxy} + \tfrac{3 \, B_{0} \, \delta_{1}}{2} \mathcal{T}_{yxz} \\
  % &\quad - \tfrac{3 \, B_{0} \, \delta_{1}}{2} \mathcal{T}_{zxy} - 2 \, \alpha_{3} \mathcal{T}_{zxz} + \tfrac{3}{\kappa^2} \mathcal{T}_{zxz} - \tfrac{3 \, B_{0} \, \delta_{1}}{2} \partial_z h_{xy} - \tfrac{3 \, \delta_{1}}{2} \partial_z^2 A_{x} - \tfrac{1}{2} \delta_{1} \partial_t^2 A_{x}
  % \end{aligned} \\[1ex]
  % \begin{aligned}
  %   0 &= 4 \, \alpha_{1} \mathcal{T}_{tty} + 2 \, \alpha_{2} \mathcal{T}_{tty} + 2 \, \alpha_{3} \mathcal{T}_{tty} + \tfrac{1}{2 \, \kappa^2} \mathcal{T}_{tty} + 3 \, B_{0} \, \delta_{1} \mathcal{T}_{yyz} - 2 \, \alpha_{3} \mathcal{T}_{zyz} + \tfrac{3}{\kappa^2} \mathcal{T}_{zyz} + 2 \, \alpha_{3} \mathcal{T}_{xxy} - \tfrac{3}{\kappa^2} \mathcal{T}_{xxy} \\
  % &\quad + \tfrac{B_{0} \, \delta_{1}}{4} \partial_z h_{tt} + \tfrac{3 \, B_{0} \, \delta_{1}}{4} \partial_z h_{xx} - \tfrac{3 \, B_{0} \, \delta_{1}}{4} \partial_z h_{yy} - \tfrac{3 \, B_{0} \, \delta_{1}}{4} \partial_z h_{zz} - \tfrac{3 \, \delta_{1}}{2} \partial_z^2 A_{y} - \tfrac{1}{2} B_{0} \, \delta_{1} \partial_t h_{tz} \\
  % &\quad - \tfrac{1}{2} \delta_{1} \partial_t^2 A_{y}
  % \end{aligned} \\[1ex]
  % \begin{aligned}
  %   0 &= 4 \, \alpha_{1} \mathcal{T}_{xxz} + 2 \, \alpha_{2} \mathcal{T}_{xxz} + 2 \, \alpha_{3} \mathcal{T}_{xxz} + \tfrac{1}{2 \, \kappa^2} \mathcal{T}_{xxz} + 2 \, \alpha_{3} \mathcal{T}_{yyz} - \tfrac{3}{\kappa^2} \mathcal{T}_{yyz} + 2 \, \alpha_{3} \mathcal{T}_{ttz} - \tfrac{3}{\kappa^2} \mathcal{T}_{ttz} + 3 \, B_{0} \, \delta_{1} \mathcal{T}_{zyz} \\
  % &\quad - \tfrac{3 \, B_{0} \, \delta_{1}}{2} \partial_t h_{ty} - \tfrac{3 \, \delta_{1}}{2} \partial_z \dot{A}_{t} + \tfrac{3 \, \delta_{1}}{2} \partial_t^2 A_{z}
  % \end{aligned} \\[1ex]
  % \begin{aligned}
  %   0 &= \tfrac{3 \, B_{0} \, \delta_{1}}{2} \mathcal{T}_{ttx} + 4 \, \alpha_{1} \mathcal{T}_{xyz} + \tfrac{2}{\kappa^2} \mathcal{T}_{xyz} -2 \, B_{0} \, \delta_{1} \mathcal{T}_{yxy} + 2 \, \alpha_{2} \mathcal{T}_{yxz} + \tfrac{3}{2 \, \kappa^2} \mathcal{T}_{yxz} -2 \, \alpha_{2} \mathcal{T}_{zxy} -\tfrac{3}{2 \, \kappa^2} \mathcal{T}_{zxy} -2 \, B_{0} \, \delta_{1} \mathcal{T}_{zxz}
  % \end{aligned} \\[1ex]
  % \begin{aligned}
  %   0 &= -4 \, \alpha_{1} \mathcal{T}_{ytx} -\tfrac{2}{\kappa^2} \mathcal{T}_{ytx} -2 \, \alpha_{2} \mathcal{T}_{txy} -\tfrac{3}{2 \, \kappa^2} \mathcal{T}_{txy} -\tfrac{1}{2} B_{0} \, \delta_{1} \mathcal{T}_{txz} -2 \, \alpha_{2} \mathcal{T}_{xty} -\tfrac{3}{2 \, \kappa^2} \mathcal{T}_{xty} -\tfrac{1}{2} B_{0} \, \delta_{1} \mathcal{T}_{xtz}
  % \end{aligned} \\[1ex]
  % \begin{aligned}
  %   0 &= -4 \, \alpha_{1} \mathcal{T}_{yty} - 2 \, \alpha_{2} \mathcal{T}_{yty} - 2 \, \alpha_{3} \mathcal{T}_{yty} - \tfrac{1}{2 \, \kappa^2} \mathcal{T}_{yty} + B_{0} \, \delta_{1} \mathcal{T}_{ytz} - B_{0} \, \delta_{1} \mathcal{T}_{zty} - 2 \, \alpha_{3} \mathcal{T}_{ztz} + \tfrac{3}{\kappa^2} \mathcal{T}_{ztz} - 2 \, B_{0} \, \delta_{1} \mathcal{T}_{tyz} - 2 \, \alpha_{3} \mathcal{T}_{xtx} \\
  % &\quad + \tfrac{3}{\kappa^2} \mathcal{T}_{xtx} - \tfrac{1}{2} B_{0} \, \delta_{1} \partial_z h_{ty} - \tfrac{3 \, \delta_{1}}{2} \partial_z^2 A_{t} - B_{0} \, \delta_{1} \partial_t h_{yz} + \tfrac{3 \, \delta_{1}}{2} \partial_z \dot{A}_{z}
  % \end{aligned} \\[1ex]
  % \begin{aligned}
  %   0 &= B_{0} \, \delta_{1} \mathcal{T}_{yty} - 4 \, \alpha_{1} \mathcal{T}_{ytz} - \tfrac{2}{\kappa^2} \mathcal{T}_{ytz} - 2 \, \alpha_{2} \mathcal{T}_{zty} - \tfrac{3}{2 \, \kappa^2} \mathcal{T}_{zty} + B_{0} \, \delta_{1} \mathcal{T}_{ztz} + 2 \, \alpha_{2} \mathcal{T}_{tyz} + \tfrac{3}{2 \, \kappa^2} \mathcal{T}_{tyz} + \tfrac{3 \, B_{0} \, \delta_{1}}{2} \mathcal{T}_{xtx} \\
  % &\quad + B_{0} \, \delta_{1} \partial_t h_{yy}
  % \end{aligned} \\[1ex]
  % \begin{aligned}
  %   0 &= -2 \, \alpha_{3} \mathcal{T}_{ttx} + \tfrac{3}{\kappa^2} \mathcal{T}_{ttx} - 2 \, B_{0} \, \delta_{1} \mathcal{T}_{xyz} + 4 \, \alpha_{1} \mathcal{T}_{yxy} + 2 \, \alpha_{2} \mathcal{T}_{yxy} + 2 \, \alpha_{3} \mathcal{T}_{yxy} + \tfrac{1}{2 \, \kappa^2} \mathcal{T}_{yxy} - B_{0} \, \delta_{1} \mathcal{T}_{yxz} + B_{0} \, \delta_{1} \mathcal{T}_{zxy} + 2 \, \alpha_{3} \mathcal{T}_{zxz} \\
  % &\quad - \tfrac{3}{\kappa^2} \mathcal{T}_{zxz} + \tfrac{B_{0} \, \delta_{1}}{2} \partial_z h_{xy} + \tfrac{3 \, \delta_{1}}{2} \partial_z^2 A_{x} - \tfrac{3 \, \delta_{1}}{2} \partial_t^2 A_{x}
  % \end{aligned} \\[1ex]
  % \begin{aligned}
  %   0 &= \tfrac{3 \, B_{0} \, \delta_{1}}{2} \mathcal{T}_{ttx} + 2 \, \alpha_{2} \mathcal{T}_{xyz} + \tfrac{3}{2 \, \kappa^2} \mathcal{T}_{xyz} -B_{0} \, \delta_{1} \mathcal{T}_{yxy} + 4 \, \alpha_{1} \mathcal{T}_{yxz} + \tfrac{2}{\kappa^2} \mathcal{T}_{yxz} + 2 \, \alpha_{2} \mathcal{T}_{zxy} + \tfrac{3}{2 \, \kappa^2} \mathcal{T}_{zxy} -B_{0} \, \delta_{1} \mathcal{T}_{zxz}
  % \end{aligned} \\[1ex]
  % \begin{aligned}
  %   0 &= 3 \, B_{0} \, \delta_{1} \mathcal{T}_{tty} + 2 \, \alpha_{3} \mathcal{T}_{xxz} - \tfrac{3}{\kappa^2} \mathcal{T}_{xxz} + 4 \, \alpha_{1} \mathcal{T}_{yyz} + 2 \, \alpha_{2} \mathcal{T}_{yyz} + 2 \, \alpha_{3} \mathcal{T}_{yyz} + \tfrac{1}{2 \, \kappa^2} \mathcal{T}_{yyz} + 2 \, \alpha_{3} \mathcal{T}_{ttz} - \tfrac{3}{\kappa^2} \mathcal{T}_{ttz} \\
  % &\quad + 3 \, B_{0} \, \delta_{1} \mathcal{T}_{xxy} - \tfrac{3 \, B_{0} \, \delta_{1}}{2} \partial_t h_{ty} - \tfrac{3 \, \delta_{1}}{2} \partial_z \dot{A}_{t} + \tfrac{3 \, \delta_{1}}{2} \partial_t^2 A_{z}
  % \end{aligned} \\[1ex]
  % \begin{aligned}
  %   0 &= -4 \, \alpha_{1} \mathcal{T}_{ztx} -\tfrac{2}{\kappa^2} \mathcal{T}_{ztx} + \tfrac{B_{0} \, \delta_{1}}{2} \mathcal{T}_{txy} -2 \, \alpha_{2} \mathcal{T}_{txz} -\tfrac{3}{2 \, \kappa^2} \mathcal{T}_{txz} + \tfrac{B_{0} \, \delta_{1}}{2} \mathcal{T}_{xty} -2 \, \alpha_{2} \mathcal{T}_{xtz} -\tfrac{3}{2 \, \kappa^2} \mathcal{T}_{xtz} -2 \, \delta_{1} \partial_z \dot{A}_{x}
  % \end{aligned} \\[1ex]
  % \begin{aligned}
  %   0 &= -B_{0} \, \delta_{1} \mathcal{T}_{yty} - 2 \, \alpha_{2} \mathcal{T}_{ytz} - \tfrac{3}{2 \, \kappa^2} \mathcal{T}_{ytz} - 4 \, \alpha_{1} \mathcal{T}_{zty} - \tfrac{2}{\kappa^2} \mathcal{T}_{zty} - B_{0} \, \delta_{1} \mathcal{T}_{ztz} - 2 \, \alpha_{2} \mathcal{T}_{tyz} - \tfrac{3}{2 \, \kappa^2} \mathcal{T}_{tyz} - \tfrac{3 \, B_{0} \, \delta_{1}}{2} \mathcal{T}_{xtx} \\
  % &\quad - B_{0} \, \delta_{1} \partial_t h_{zz} - 2 \, \delta_{1} \partial_z \dot{A}_{y}
  % \end{aligned} \\[1ex]
  % \begin{aligned}
  %   0 &= 2 \, \alpha_{3} \mathcal{T}_{xxz} - \tfrac{3}{\kappa^2} \mathcal{T}_{xxz} + 2 \, \alpha_{3} \mathcal{T}_{yyz} - \tfrac{3}{\kappa^2} \mathcal{T}_{yyz} + 4 \, \alpha_{1} \mathcal{T}_{ttz} + 2 \, \alpha_{2} \mathcal{T}_{ttz} + 2 \, \alpha_{3} \mathcal{T}_{ttz} + \tfrac{1}{2 \, \kappa^2} \mathcal{T}_{ttz} + 3 \, B_{0} \, \delta_{1} \mathcal{T}_{zyz} \\
  % &\quad + \tfrac{B_{0} \, \delta_{1}}{2} \partial_t h_{ty} + \tfrac{\delta_{1}}{2} \partial_z \dot{A}_{t} - \tfrac{1}{2} \delta_{1} \partial_t^2 A_{z}
  % \end{aligned} \\[1ex]
  % \begin{aligned}
  %   0 &= -2 \, \alpha_{3} \mathcal{T}_{yty} + \tfrac{3}{\kappa^2} \mathcal{T}_{yty} + B_{0} \, \delta_{1} \mathcal{T}_{ytz} - B_{0} \, \delta_{1} \mathcal{T}_{zty} - 4 \, \alpha_{1} \mathcal{T}_{ztz} - 2 \, \alpha_{2} \mathcal{T}_{ztz} - 2 \, \alpha_{3} \mathcal{T}_{ztz} - \tfrac{1}{2 \, \kappa^2} \mathcal{T}_{ztz} - 2 \, B_{0} \, \delta_{1} \mathcal{T}_{tyz} - 2 \, \alpha_{3} \mathcal{T}_{xtx} \\
  % &\quad + \tfrac{3}{\kappa^2} \mathcal{T}_{xtx} - \tfrac{1}{2} B_{0} \, \delta_{1} \partial_z h_{ty} + \tfrac{\delta_{1}}{2} \partial_z^2 A_{t} + B_{0} \, \delta_{1} \partial_t h_{yz} - \tfrac{1}{2} \delta_{1} \partial_z \dot{A}_{z}
  % \end{aligned} \\[1ex]
  % \begin{aligned}
  %   0 &= -\tfrac{3 \, B_{0} \, \delta_{1}}{2} \mathcal{T}_{ttx} -2 \, \alpha_{2} \mathcal{T}_{xyz} -\tfrac{3}{2 \, \kappa^2} \mathcal{T}_{xyz} + B_{0} \, \delta_{1} \mathcal{T}_{yxy} + 2 \, \alpha_{2} \mathcal{T}_{yxz} + \tfrac{3}{2 \, \kappa^2} \mathcal{T}_{yxz} + 4 \, \alpha_{1} \mathcal{T}_{zxy} + \tfrac{2}{\kappa^2} \mathcal{T}_{zxy} + B_{0} \, \delta_{1} \mathcal{T}_{zxz}
  % \end{aligned} \\[1ex]
  % \begin{aligned}
  %   0 &= -2 \, \alpha_{3} \mathcal{T}_{ttx} + \tfrac{3}{\kappa^2} \mathcal{T}_{ttx} - 2 \, B_{0} \, \delta_{1} \mathcal{T}_{xyz} + 2 \, \alpha_{3} \mathcal{T}_{yxy} - \tfrac{3}{\kappa^2} \mathcal{T}_{yxy} - B_{0} \, \delta_{1} \mathcal{T}_{yxz} + B_{0} \, \delta_{1} \mathcal{T}_{zxy} + 4 \, \alpha_{1} \mathcal{T}_{zxz} + 2 \, \alpha_{2} \mathcal{T}_{zxz} + 2 \, \alpha_{3} \mathcal{T}_{zxz} \\
  % &\quad + \tfrac{1}{2 \, \kappa^2} \mathcal{T}_{zxz} + \tfrac{B_{0} \, \delta_{1}}{2} \partial_z h_{xy} - \tfrac{1}{2} \delta_{1} \partial_z^2 A_{x} - \tfrac{3 \, \delta_{1}}{2} \partial_t^2 A_{x}
  % \end{aligned} \\[1ex]
  % \begin{aligned}
  %   0 &= -2 \, \alpha_{3} \mathcal{T}_{tty} + \tfrac{3}{\kappa^2} \mathcal{T}_{tty} + 3 \, B_{0} \, \delta_{1} \mathcal{T}_{xxz} + 3 \, B_{0} \, \delta_{1} \mathcal{T}_{ttz} + 4 \, \alpha_{1} \mathcal{T}_{zyz} + 2 \, \alpha_{2} \mathcal{T}_{zyz} + 2 \, \alpha_{3} \mathcal{T}_{zyz} + \tfrac{1}{2 \, \kappa^2} \mathcal{T}_{zyz} - 2 \, \alpha_{3} \mathcal{T}_{xxy} \\
  % &\quad + \tfrac{3}{\kappa^2} \mathcal{T}_{xxy} + \tfrac{3 \, B_{0} \, \delta_{1}}{4} \partial_z h_{tt} - \tfrac{3 \, B_{0} \, \delta_{1}}{4} \partial_z h_{xx} - \tfrac{1}{4} B_{0} \, \delta_{1} \partial_z h_{yy} - \tfrac{1}{4} B_{0} \, \delta_{1} \partial_z h_{zz} - \tfrac{1}{2} \delta_{1} \partial_z^2 A_{y} \\
  % &\quad - \tfrac{3 \, B_{0} \, \delta_{1}}{2} \partial_t h_{tz} - \tfrac{3 \, \delta_{1}}{2} \partial_t^2 A_{y}
  % \end{aligned} \\[1ex]
  % \begin{aligned}
  %   0 &= -2 \, \alpha_{2} \mathcal{T}_{ytx} -\tfrac{3}{2 \, \kappa^2} \mathcal{T}_{ytx} + \tfrac{B_{0} \, \delta_{1}}{2} \mathcal{T}_{ztx} -4 \, \alpha_{1} \mathcal{T}_{txy} -\tfrac{2}{\kappa^2} \mathcal{T}_{txy} + 2 \, \alpha_{2} \mathcal{T}_{xty} + \tfrac{3}{2 \, \kappa^2} \mathcal{T}_{xty} -B_{0} \, \delta_{1} \partial_z h_{tx} + B_{0} \, \delta_{1} \partial_t h_{xz}
  % \end{aligned} \\[1ex]
  % \begin{aligned}
  %   0 &= -\tfrac{1}{2} B_{0} \, \delta_{1} \mathcal{T}_{ytx} -2 \, \alpha_{2} \mathcal{T}_{ztx} -\tfrac{3}{2 \, \kappa^2} \mathcal{T}_{ztx} -4 \, \alpha_{1} \mathcal{T}_{txz} -\tfrac{2}{\kappa^2} \mathcal{T}_{txz} + 2 \, \alpha_{2} \mathcal{T}_{xtz} + \tfrac{3}{2 \, \kappa^2} \mathcal{T}_{xtz} -B_{0} \, \delta_{1} \partial_t h_{xy} -2 \, \delta_{1} \partial_z \dot{A}_{x}
  % \end{aligned} \\[1ex]
  % \begin{aligned}
  %   0 &= -2 \, B_{0} \, \delta_{1} \mathcal{T}_{yty} + 2 \, \alpha_{2} \mathcal{T}_{ytz} + \tfrac{3}{2 \, \kappa^2} \mathcal{T}_{ytz} - 2 \, \alpha_{2} \mathcal{T}_{zty} - \tfrac{3}{2 \, \kappa^2} \mathcal{T}_{zty} - 2 \, B_{0} \, \delta_{1} \mathcal{T}_{ztz} - 4 \, \alpha_{1} \mathcal{T}_{tyz} - \tfrac{2}{\kappa^2} \mathcal{T}_{tyz} - \tfrac{3 \, B_{0} \, \delta_{1}}{2} \mathcal{T}_{xtx} \\
  % &\quad - B_{0} \, \delta_{1} \partial_t h_{yy} - B_{0} \, \delta_{1} \partial_t h_{zz} - 2 \, \delta_{1} \partial_z \dot{A}_{y}
  % \end{aligned} \\[1ex]
  % \begin{aligned}
  %   0 &= -2 \, \alpha_{3} \mathcal{T}_{yty} + \tfrac{3}{\kappa^2} \mathcal{T}_{yty} + \tfrac{3 \, B_{0} \, \delta_{1}}{2} \mathcal{T}_{ytz} - \tfrac{3 \, B_{0} \, \delta_{1}}{2} \mathcal{T}_{zty} - 2 \, \alpha_{3} \mathcal{T}_{ztz} + \tfrac{3}{\kappa^2} \mathcal{T}_{ztz} - \tfrac{3 \, B_{0} \, \delta_{1}}{2} \mathcal{T}_{tyz} - 4 \, \alpha_{1} \mathcal{T}_{xtx} - 2 \, \alpha_{2} \mathcal{T}_{xtx} \\
  % &\quad - 2 \, \alpha_{3} \mathcal{T}_{xtx} - \tfrac{1}{2 \, \kappa^2} \mathcal{T}_{xtx} - \tfrac{3 \, B_{0} \, \delta_{1}}{2} \partial_z h_{ty} - \tfrac{3 \, \delta_{1}}{2} \partial_z^2 A_{t} + \tfrac{3 \, \delta_{1}}{2} \partial_z \dot{A}_{z}
  % \end{aligned} \\[1ex]
  % \begin{aligned}
  %   0 &= -2 \, \alpha_{2} \mathcal{T}_{ytx} -\tfrac{3}{2 \, \kappa^2} \mathcal{T}_{ytx} + \tfrac{B_{0} \, \delta_{1}}{2} \mathcal{T}_{ztx} + 2 \, \alpha_{2} \mathcal{T}_{txy} + \tfrac{3}{2 \, \kappa^2} \mathcal{T}_{txy} -4 \, \alpha_{1} \mathcal{T}_{xty} -\tfrac{2}{\kappa^2} \mathcal{T}_{xty} + B_{0} \, \delta_{1} \partial_z h_{tx} -B_{0} \, \delta_{1} \partial_t h_{xz}
  % \end{aligned} \\[1ex]
  % \begin{aligned}
  %   0 &= -\tfrac{1}{2} B_{0} \, \delta_{1} \mathcal{T}_{ytx} -2 \, \alpha_{2} \mathcal{T}_{ztx} -\tfrac{3}{2 \, \kappa^2} \mathcal{T}_{ztx} + 2 \, \alpha_{2} \mathcal{T}_{txz} + \tfrac{3}{2 \, \kappa^2} \mathcal{T}_{txz} -4 \, \alpha_{1} \mathcal{T}_{xtz} -\tfrac{2}{\kappa^2} \mathcal{T}_{xtz} + B_{0} \, \delta_{1} \partial_t h_{xy}
  % \end{aligned} \\[1ex]
  % \begin{aligned}
  %   0 &= 2 \, \alpha_{3} \mathcal{T}_{tty} - \tfrac{3}{\kappa^2} \mathcal{T}_{tty} + 3 \, B_{0} \, \delta_{1} \mathcal{T}_{yyz} - 2 \, \alpha_{3} \mathcal{T}_{zyz} + \tfrac{3}{\kappa^2} \mathcal{T}_{zyz} + 4 \, \alpha_{1} \mathcal{T}_{xxy} + 2 \, \alpha_{2} \mathcal{T}_{xxy} + 2 \, \alpha_{3} \mathcal{T}_{xxy} + \tfrac{1}{2 \, \kappa^2} \mathcal{T}_{xxy} \\
  % &\quad - \tfrac{3 \, B_{0} \, \delta_{1}}{4} \partial_z h_{tt} - \tfrac{1}{4} B_{0} \, \delta_{1} \partial_z h_{xx} - \tfrac{3 \, B_{0} \, \delta_{1}}{4} \partial_z h_{yy} - \tfrac{3 \, B_{0} \, \delta_{1}}{4} \partial_z h_{zz} - \tfrac{3 \, \delta_{1}}{2} \partial_z^2 A_{y} + \tfrac{3 \, B_{0} \, \delta_{1}}{2} \partial_t h_{tz} \\
  % &\quad + \tfrac{3 \, \delta_{1}}{2} \partial_t^2 A_{y}
  % \end{aligned} \\[1ex]
  \begin{aligned}
    \mathscr{H} &\supset -\tfrac{1}{16} B_{0}^2 h_{tt}^2 + \tfrac{B_{0}^2}{4} h_{tx}^2 - \tfrac{1}{4} B_{0}^2 h_{ty}^2 - \tfrac{1}{4} B_{0}^2 h_{tz}^2 - \tfrac{1}{8} B_{0}^2 h_{tt} \, h_{xx} - \tfrac{1}{16} B_{0}^2 h_{xx}^2 + \tfrac{B_{0}^2}{4} h_{xy}^2 + \tfrac{B_{0}^2}{4} h_{xz}^2 \\
  &\quad + \tfrac{B_{0}^2}{8} h_{tt} \, h_{yy} - \tfrac{1}{8} B_{0}^2 h_{xx} \, h_{yy} + \tfrac{3 \, B_{0}^2}{16} h_{yy}^2 + \tfrac{B_{0}^2}{4} h_{yz}^2 + \tfrac{B_{0}^2}{8} h_{tt} \, h_{zz} - \tfrac{1}{8} B_{0}^2 h_{xx} \, h_{zz} + \tfrac{B_{0}^2}{8} h_{yy} \, h_{zz} \\
  &\quad + \tfrac{3 \, B_{0}^2}{16} h_{zz}^2 - \tfrac{1}{2} \partial_z A_{t}^2 + \tfrac{1}{2} \partial_z A_{x}^2 + \tfrac{1}{2} \partial_z A_{y}^2 + \tfrac{1}{2} \partial_z A_{z}^2 - \tfrac{1}{2 \, \kappa^2} \partial_z h_{tx}^2 - \tfrac{1}{2 \, \kappa^2} \partial_z h_{ty}^2 \\
  &\quad + \tfrac{1}{2 \, \kappa^2} \partial_z h_{tt} \, \partial_z h_{xx} + \tfrac{1}{2 \, \kappa^2} \partial_z h_{xy}^2 + \tfrac{1}{2 \, \kappa^2} \partial_z h_{tt} \, \partial_z h_{yy} - \tfrac{1}{2 \, \kappa^2} \partial_z h_{xx} \, \partial_z h_{yy} - \tfrac{1}{2} \dot{A}_{t}^2 \\
  &\quad + \tfrac{1}{2} \dot{A}_{x}^2 + \tfrac{1}{2} \dot{A}_{y}^2 + \tfrac{1}{2} \dot{A}_{z}^2 + \tfrac{1}{2 \, \kappa^2} \dot{h}_{xy}^2 + \tfrac{1}{2 \, \kappa^2} \dot{h}_{xz}^2 - \tfrac{1}{2 \, \kappa^2} \dot{h}_{xx} \, \dot{h}_{yy} + \tfrac{1}{2 \, \kappa^2} \dot{h}_{yz}^2 \\
  &\quad - \tfrac{1}{2 \, \kappa^2} \dot{h}_{xx} \, \dot{h}_{zz} - \tfrac{1}{2 \, \kappa^2} \dot{h}_{yy} \, \dot{h}_{zz} - \tfrac{2}{\kappa^2} h_{xz} \, \partial_t \partial_z h_{tx} - \tfrac{2}{\kappa^2} h_{yz} \, \partial_t \partial_z h_{ty} \\
  &\quad + \tfrac{1}{\kappa^{2}} h_{tt} \, \partial_t \partial_z h_{tz} + \tfrac{1}{\kappa^{2}} h_{xx} \, \partial_t \partial_z h_{tz} + \tfrac{1}{\kappa^{2}} h_{yy} \, \partial_t \partial_z h_{tz} - \tfrac{1}{\kappa^{2}} h_{zz} \, \partial_t \partial_z h_{tz} \\
  &\quad + \tfrac{2}{\kappa^2} h_{tz} \, \partial_t \partial_z h_{xx} - \tfrac{2}{\kappa^2} h_{tx} \, \partial_t \partial_z h_{xz} + \tfrac{2}{\kappa^2} h_{tz} \, \partial_t \partial_z h_{yy} - \tfrac{2}{\kappa^2} h_{ty} \, \partial_t \partial_z h_{yz}
  \end{aligned}

\end{gathered}
$}
\medskip


% --- E.6  Full YM--PGT nonminimal -----------------------------------
\subsection{Yang--Mills--PGT}%
\label{SymbolicGeneralNonminimal}
 The full Yang--Mills-form Poincar\'e
gauge Lagrangian turns on $\beta_{1,2,3}$, $\xi$, $\delta_{1}$,
$\chi$, $\zeta_{1,2,3}$ simultaneously. The Bahamonde et al.~\cite{bahamonde2025torsion},
Barker~\cite{barker2024poincare}, and Shapiro~\cite{shapiro2002torsion}
variants of~\cref{SectorSurvey} are parameter restrictions of
\cref{NonMinimalLagrangian} (\cref{SectorSurvey}). The symbolic output
comprises \GeneralNonminimalEomCount{} equations of motion
(\GeneralNonminimalConstraintCount{} constraints) with
\GeneralNonminimalTermCount{} RHS terms.

\medskip
\noindent\adjustbox{max width=\linewidth}{$\displaystyle
\setlength{\jot}{1pt}%
\begin{gathered}
\mathscr{L} &= \frac{1}{\kappa^2} \tilde{\mathcal{R}} + \beta_{1} \tensor{T}{_a_b_c} \tensor{g}{^a^d} \tensor{g}{^b^e} \tensor{g}{^c^f} \tensor{T}{_d_e_f} + \beta_{2} \tensor{T}{_a_b_c} \tensor{g}{^b^d} \tensor{g}{^a^e} \tensor{g}{^c^f} \tensor{T}{_d_e_f} + \beta_{3} \tensor{T}{^b_b_a} \tensor{g}{^a^d} \tensor{T}{^c_c_d} - \frac{1}{4} \xi \tensor{\tilde{F}}{_a_b} \tensor{g}{^a^c} \tensor{g}{^b^d} \tensor{\tilde{F}}{_c_d} + \delta_{1} \tensor{\tilde{R}}{_a_b} \tensor{g}{^a^c} \tensor{g}{^b^d} \tensor{F}{_c_d} + \chi \tensor{\tilde{R}}{_a_b} \tensor{g}{^a^c} \tensor{g}{^b^d} \tensor{\tilde{F}}{_c_d} + \zeta_{1} \nabla_{a} \tensor{T}{^b_b_c} \tensor{g}{^a^d} \tensor{g}{^c^e} \tensor{F}{_d_e} + \zeta_{2} \nabla_{a} \tensor{T}{^a_b_c} \tensor{g}{^b^d} \tensor{g}{^c^e} \tensor{F}{_d_e} + \zeta_{3} \nabla_{a} \tensor{T}{_b^a_c} \tensor{g}{^b^d} \tensor{g}{^c^e} \tensor{F}{_d_e} - \frac{1}{4} \tensor{F}{_a_b} \tensor{g}{^a^c} \tensor{g}{^b^d} \tensor{F}{_c_d} \\
\partial_t^{2} \mathcal{A}_{t} &= -B_{0} \partial_z \mathcal{H}_{ty} - \partial_z^2 \mathcal{A}_{t} -\frac{3 \, \delta_{1}}{2} \partial_z^2 \mathcal{T}_{yty} -\zeta_{1} \partial_z^2 \mathcal{T}_{yty} + \frac{\delta_{1}}{2} \partial_z^2 \mathcal{T}_{ztz} -\zeta_{1} \partial_z^2 \mathcal{T}_{ztz} -2 \, \zeta_{2} \partial_z^2 \mathcal{T}_{ztz} -\zeta_{3} \partial_z^2 \mathcal{T}_{ztz} -\frac{3 \, \delta_{1}}{2} \partial_z^2 \mathcal{T}_{xtx} -\zeta_{1} \partial_z^2 \mathcal{T}_{xtx} -\frac{3 \, \delta_{1}}{2} \partial_z \dot{\mathcal{T}}_{xxz} -\zeta_{1} \partial_z \dot{\mathcal{T}}_{xxz} -\frac{3 \, \delta_{1}}{2} \partial_z \dot{\mathcal{T}}_{yyz} -\zeta_{1} \partial_z \dot{\mathcal{T}}_{yyz} + \frac{\delta_{1}}{2} \partial_z \dot{\mathcal{T}}_{ttz} -\zeta_{1} \partial_z \dot{\mathcal{T}}_{ttz} -2 \, \zeta_{2} \partial_z \dot{\mathcal{T}}_{ttz} -\zeta_{3} \partial_z \dot{\mathcal{T}}_{ttz} \\
\partial_t^{2} \mathcal{A}_{x} &= B_{0} \partial_z \mathcal{H}_{xy} + \partial_z^2 \mathcal{A}_{x} -\frac{3 \, \delta_{1}}{2} \partial_z^2 \mathcal{T}_{ttx} -\zeta_{1} \partial_z^2 \mathcal{T}_{ttx} + \frac{3 \, \delta_{1}}{2} \partial_z^2 \mathcal{T}_{yxy} + \zeta_{1} \partial_z^2 \mathcal{T}_{yxy} -\frac{1}{2} \delta_{1} \partial_z^2 \mathcal{T}_{zxz} + \zeta_{1} \partial_z^2 \mathcal{T}_{zxz} + 2 \, \zeta_{2} \partial_z^2 \mathcal{T}_{zxz} + \zeta_{3} \partial_z^2 \mathcal{T}_{zxz} -2 \, \delta_{1} \partial_z \dot{\mathcal{T}}_{ztx} + 2 \, \zeta_{2} \partial_z \dot{\mathcal{T}}_{ztx} + \zeta_{3} \partial_z \dot{\mathcal{T}}_{ztx} -2 \, \delta_{1} \partial_z \dot{\mathcal{T}}_{txz} + 2 \, \zeta_{2} \partial_z \dot{\mathcal{T}}_{txz} + \zeta_{3} \partial_z \dot{\mathcal{T}}_{txz} -\frac{1}{2} \delta_{1} \partial_t^2 \mathcal{T}_{ttx} + \zeta_{1} \partial_t^2 \mathcal{T}_{ttx} + 2 \, \zeta_{2} \partial_t^2 \mathcal{T}_{ttx} + \zeta_{3} \partial_t^2 \mathcal{T}_{ttx} -\frac{3 \, \delta_{1}}{2} \partial_t^2 \mathcal{T}_{yxy} -\zeta_{1} \partial_t^2 \mathcal{T}_{yxy} -\frac{3 \, \delta_{1}}{2} \partial_t^2 \mathcal{T}_{zxz} -\zeta_{1} \partial_t^2 \mathcal{T}_{zxz} \\
\partial_t^{2} \mathcal{A}_{y} &= \frac{B_{0}}{2} \partial_z \mathcal{H}_{tt} -\frac{1}{2} B_{0} \partial_z \mathcal{H}_{xx} + \frac{B_{0}}{2} \partial_z \mathcal{H}_{yy} + \frac{B_{0}}{2} \partial_z \mathcal{H}_{zz} + \partial_z^2 \mathcal{A}_{y} -\frac{3 \, \delta_{1}}{2} \partial_z^2 \mathcal{T}_{tty} -\zeta_{1} \partial_z^2 \mathcal{T}_{tty} -\frac{1}{2} \delta_{1} \partial_z^2 \mathcal{T}_{zyz} + \zeta_{1} \partial_z^2 \mathcal{T}_{zyz} + 2 \, \zeta_{2} \partial_z^2 \mathcal{T}_{zyz} + \zeta_{3} \partial_z^2 \mathcal{T}_{zyz} -\frac{3 \, \delta_{1}}{2} \partial_z^2 \mathcal{T}_{xxy} -\zeta_{1} \partial_z^2 \mathcal{T}_{xxy} -B_{0} \partial_t \mathcal{H}_{tz} -2 \, \delta_{1} \partial_z \dot{\mathcal{T}}_{zty} + 2 \, \zeta_{2} \partial_z \dot{\mathcal{T}}_{zty} + \zeta_{3} \partial_z \dot{\mathcal{T}}_{zty} -2 \, \delta_{1} \partial_z \dot{\mathcal{T}}_{tyz} + 2 \, \zeta_{2} \partial_z \dot{\mathcal{T}}_{tyz} + \zeta_{3} \partial_z \dot{\mathcal{T}}_{tyz} -\frac{1}{2} \delta_{1} \partial_t^2 \mathcal{T}_{tty} + \zeta_{1} \partial_t^2 \mathcal{T}_{tty} + 2 \, \zeta_{2} \partial_t^2 \mathcal{T}_{tty} + \zeta_{3} \partial_t^2 \mathcal{T}_{tty} -\frac{3 \, \delta_{1}}{2} \partial_t^2 \mathcal{T}_{zyz} -\zeta_{1} \partial_t^2 \mathcal{T}_{zyz} + \frac{3 \, \delta_{1}}{2} \partial_t^2 \mathcal{T}_{xxy} + \zeta_{1} \partial_t^2 \mathcal{T}_{xxy} \\
\partial_t^{2} \mathcal{A}_{z} &= \partial_z^2 \mathcal{A}_{z} + B_{0} \partial_t \mathcal{H}_{ty} + \frac{3 \, \delta_{1}}{2} \partial_z \dot{\mathcal{T}}_{yty} + \zeta_{1} \partial_z \dot{\mathcal{T}}_{yty} -\frac{1}{2} \delta_{1} \partial_z \dot{\mathcal{T}}_{ztz} + \zeta_{1} \partial_z \dot{\mathcal{T}}_{ztz} + 2 \, \zeta_{2} \partial_z \dot{\mathcal{T}}_{ztz} + \zeta_{3} \partial_z \dot{\mathcal{T}}_{ztz} + \frac{3 \, \delta_{1}}{2} \partial_z \dot{\mathcal{T}}_{xtx} + \zeta_{1} \partial_z \dot{\mathcal{T}}_{xtx} + \frac{3 \, \delta_{1}}{2} \partial_t^2 \mathcal{T}_{xxz} + \zeta_{1} \partial_t^2 \mathcal{T}_{xxz} + \frac{3 \, \delta_{1}}{2} \partial_t^2 \mathcal{T}_{yyz} + \zeta_{1} \partial_t^2 \mathcal{T}_{yyz} -\frac{1}{2} \delta_{1} \partial_t^2 \mathcal{T}_{ttz} + \zeta_{1} \partial_t^2 \mathcal{T}_{ttz} + 2 \, \zeta_{2} \partial_t^2 \mathcal{T}_{ttz} + \zeta_{3} \partial_t^2 \mathcal{T}_{ttz} \\
0 &= \frac{B_{0}^2}{8} \mathcal{H}_{tt} + \frac{B_{0}^2}{8} \mathcal{H}_{xx} -\frac{1}{8} B_{0}^2 \mathcal{H}_{yy} -\frac{1}{8} B_{0}^2 \mathcal{H}_{zz} -\frac{1}{2} B_{0} \partial_z \mathcal{A}_{y} -\frac{1}{4} B_{0} \, \delta_{1} \partial_z \mathcal{T}_{tty} + \frac{B_{0} \, \zeta_{1}}{2} \partial_z \mathcal{T}_{tty} + B_{0} \, \zeta_{2} \partial_z \mathcal{T}_{tty} + \frac{B_{0} \, \zeta_{3}}{2} \partial_z \mathcal{T}_{tty} -\frac{3 \, B_{0} \, \delta_{1}}{4} \partial_z \mathcal{T}_{zyz} -\frac{1}{2} B_{0} \, \zeta_{1} \partial_z \mathcal{T}_{zyz} + \frac{3 \, B_{0} \, \delta_{1}}{4} \partial_z \mathcal{T}_{xxy} + \frac{B_{0} \, \zeta_{1}}{2} \partial_z \mathcal{T}_{xxy} + \frac{1}{2 \, \kappa^2} \partial_z^2 \mathcal{H}_{xx} + \frac{1}{2 \, \kappa^2} \partial_z^2 \mathcal{H}_{yy} \\
0 &= -\frac{1}{2} B_{0}^2 \mathcal{H}_{tx} + B_{0} \, \delta_{1} \partial_z \mathcal{T}_{txy} -B_{0} \, \zeta_{2} \partial_z \mathcal{T}_{txy} -\frac{1}{2} B_{0} \, \zeta_{3} \partial_z \mathcal{T}_{txy} -B_{0} \, \delta_{1} \partial_z \mathcal{T}_{xty} + B_{0} \, \zeta_{2} \partial_z \mathcal{T}_{xty} + \frac{B_{0} \, \zeta_{3}}{2} \partial_z \mathcal{T}_{xty} -\frac{1}{\kappa^{2}} \partial_z^2 \mathcal{H}_{tx} + \frac{1}{\kappa^{2}} \partial_z \dot{\mathcal{H}}_{xz} \\
0 &= \frac{B_{0}^2}{2} \mathcal{H}_{ty} + B_{0} \partial_z \mathcal{A}_{t} + \frac{B_{0} \, \delta_{1}}{2} \partial_z \mathcal{T}_{yty} + B_{0} \, \zeta_{1} \partial_z \mathcal{T}_{yty} + B_{0} \, \zeta_{2} \partial_z \mathcal{T}_{yty} + \frac{B_{0} \, \zeta_{3}}{2} \partial_z \mathcal{T}_{yty} + \frac{B_{0} \, \delta_{1}}{2} \partial_z \mathcal{T}_{ztz} + B_{0} \, \zeta_{1} \partial_z \mathcal{T}_{ztz} + B_{0} \, \zeta_{2} \partial_z \mathcal{T}_{ztz} + \frac{B_{0} \, \zeta_{3}}{2} \partial_z \mathcal{T}_{ztz} + \frac{3 \, B_{0} \, \delta_{1}}{2} \partial_z \mathcal{T}_{xtx} + B_{0} \, \zeta_{1} \partial_z \mathcal{T}_{xtx} -\frac{1}{\kappa^{2}} \partial_z^2 \mathcal{H}_{ty} -B_{0} \partial_t \mathcal{A}_{z} + \frac{3 \, B_{0} \, \delta_{1}}{2} \partial_t \mathcal{T}_{xxz} + B_{0} \, \zeta_{1} \partial_t \mathcal{T}_{xxz} + \frac{3 \, B_{0} \, \delta_{1}}{2} \partial_t \mathcal{T}_{yyz} + B_{0} \, \zeta_{1} \partial_t \mathcal{T}_{yyz} -\frac{1}{2} B_{0} \, \delta_{1} \partial_t \mathcal{T}_{ttz} + B_{0} \, \zeta_{1} \partial_t \mathcal{T}_{ttz} + 2 \, B_{0} \, \zeta_{2} \partial_t \mathcal{T}_{ttz} + B_{0} \, \zeta_{3} \partial_t \mathcal{T}_{ttz} + \frac{1}{\kappa^{2}} \partial_z \dot{\mathcal{H}}_{yz} \\
0 &= \frac{B_{0}^2}{2} \mathcal{H}_{tz} + B_{0} \partial_t \mathcal{A}_{y} + \frac{B_{0} \, \delta_{1}}{2} \partial_t \mathcal{T}_{tty} -B_{0} \, \zeta_{1} \partial_t \mathcal{T}_{tty} -2 \, B_{0} \, \zeta_{2} \partial_t \mathcal{T}_{tty} -B_{0} \, \zeta_{3} \partial_t \mathcal{T}_{tty} + \frac{3 \, B_{0} \, \delta_{1}}{2} \partial_t \mathcal{T}_{zyz} + B_{0} \, \zeta_{1} \partial_t \mathcal{T}_{zyz} -\frac{3 \, B_{0} \, \delta_{1}}{2} \partial_t \mathcal{T}_{xxy} -B_{0} \, \zeta_{1} \partial_t \mathcal{T}_{xxy} -\frac{1}{\kappa^{2}} \partial_z \dot{\mathcal{H}}_{xx} -\frac{1}{\kappa^{2}} \partial_z \dot{\mathcal{H}}_{yy} \\
0 &= \frac{B_{0}^2}{8} \mathcal{H}_{tt} + \frac{B_{0}^2}{8} \mathcal{H}_{xx} + \frac{B_{0}^2}{8} \mathcal{H}_{yy} + \frac{B_{0}^2}{8} \mathcal{H}_{zz} + \frac{B_{0}}{2} \partial_z \mathcal{A}_{y} -\frac{3 \, B_{0} \, \delta_{1}}{4} \partial_z \mathcal{T}_{tty} -\frac{1}{2} B_{0} \, \zeta_{1} \partial_z \mathcal{T}_{tty} + \frac{3 \, B_{0} \, \delta_{1}}{4} \partial_z \mathcal{T}_{zyz} + \frac{B_{0} \, \zeta_{1}}{2} \partial_z \mathcal{T}_{zyz} + \frac{B_{0} \, \delta_{1}}{4} \partial_z \mathcal{T}_{xxy} -\frac{1}{2} B_{0} \, \zeta_{1} \partial_z \mathcal{T}_{xxy} -B_{0} \, \zeta_{2} \partial_z \mathcal{T}_{xxy} -\frac{1}{2} B_{0} \, \zeta_{3} \partial_z \mathcal{T}_{xxy} + \frac{1}{2 \, \kappa^2} \partial_z^2 \mathcal{H}_{tt} -\frac{1}{2 \, \kappa^2} \partial_z^2 \mathcal{H}_{yy} -\frac{1}{\kappa^{2}} \partial_z \dot{\mathcal{H}}_{tz} + \frac{1}{2 \, \kappa^2} \partial_t^2 \mathcal{H}_{yy} + \frac{1}{2 \, \kappa^2} \partial_t^2 \mathcal{H}_{zz} \\
\partial_t^{2} \mathcal{H}_{xy} &= \frac{B_{0}^2}{2} \mathcal{H}_{xy} + B_{0} \partial_z \mathcal{A}_{x} -\frac{3 \, B_{0} \, \delta_{1}}{2} \partial_z \mathcal{T}_{ttx} -B_{0} \, \zeta_{1} \partial_z \mathcal{T}_{ttx} + \frac{B_{0} \, \delta_{1}}{2} \partial_z \mathcal{T}_{yxy} + B_{0} \, \zeta_{1} \partial_z \mathcal{T}_{yxy} + B_{0} \, \zeta_{2} \partial_z \mathcal{T}_{yxy} + \frac{B_{0} \, \zeta_{3}}{2} \partial_z \mathcal{T}_{yxy} + \frac{B_{0} \, \delta_{1}}{2} \partial_z \mathcal{T}_{zxz} + B_{0} \, \zeta_{1} \partial_z \mathcal{T}_{zxz} + B_{0} \, \zeta_{2} \partial_z \mathcal{T}_{zxz} + \frac{B_{0} \, \zeta_{3}}{2} \partial_z \mathcal{T}_{zxz} -\frac{1}{\kappa^{2}} \partial_z^2 \mathcal{H}_{xy} -B_{0} \, \delta_{1} \partial_t \mathcal{T}_{txz} + B_{0} \, \zeta_{2} \partial_t \mathcal{T}_{txz} + \frac{B_{0} \, \zeta_{3}}{2} \partial_t \mathcal{T}_{txz} + B_{0} \, \delta_{1} \partial_t \mathcal{T}_{xtz} -B_{0} \, \zeta_{2} \partial_t \mathcal{T}_{xtz} -\frac{1}{2} B_{0} \, \zeta_{3} \partial_t \mathcal{T}_{xtz} \\
\partial_t^{2} \mathcal{H}_{xz} &= \frac{B_{0}^2}{2} \mathcal{H}_{xz} + B_{0} \, \delta_{1} \partial_t \mathcal{T}_{txy} -B_{0} \, \zeta_{2} \partial_t \mathcal{T}_{txy} -\frac{1}{2} B_{0} \, \zeta_{3} \partial_t \mathcal{T}_{txy} -B_{0} \, \delta_{1} \partial_t \mathcal{T}_{xty} + B_{0} \, \zeta_{2} \partial_t \mathcal{T}_{xty} + \frac{B_{0} \, \zeta_{3}}{2} \partial_t \mathcal{T}_{xty} -\frac{1}{\kappa^{2}} \partial_z \dot{\mathcal{H}}_{tx} \\
0 &= -\frac{1}{8} B_{0}^2 \mathcal{H}_{tt} + \frac{B_{0}^2}{8} \mathcal{H}_{xx} -\frac{3 \, B_{0}^2}{8} \mathcal{H}_{yy} -\frac{1}{8} B_{0}^2 \mathcal{H}_{zz} -\frac{1}{2} B_{0} \partial_z \mathcal{A}_{y} + \frac{3 \, B_{0} \, \delta_{1}}{4} \partial_z \mathcal{T}_{tty} + \frac{B_{0} \, \zeta_{1}}{2} \partial_z \mathcal{T}_{tty} + \frac{B_{0} \, \delta_{1}}{4} \partial_z \mathcal{T}_{zyz} -\frac{1}{2} B_{0} \, \zeta_{1} \partial_z \mathcal{T}_{zyz} -B_{0} \, \zeta_{2} \partial_z \mathcal{T}_{zyz} -\frac{1}{2} B_{0} \, \zeta_{3} \partial_z \mathcal{T}_{zyz} + \frac{3 \, B_{0} \, \delta_{1}}{4} \partial_z \mathcal{T}_{xxy} + \frac{B_{0} \, \zeta_{1}}{2} \partial_z \mathcal{T}_{xxy} + \frac{1}{2 \, \kappa^2} \partial_z^2 \mathcal{H}_{tt} -\frac{1}{2 \, \kappa^2} \partial_z^2 \mathcal{H}_{xx} -B_{0} \, \delta_{1} \partial_t \mathcal{T}_{ytz} + B_{0} \, \zeta_{2} \partial_t \mathcal{T}_{ytz} + \frac{B_{0} \, \zeta_{3}}{2} \partial_t \mathcal{T}_{ytz} + B_{0} \, \delta_{1} \partial_t \mathcal{T}_{tyz} -B_{0} \, \zeta_{2} \partial_t \mathcal{T}_{tyz} -\frac{1}{2} B_{0} \, \zeta_{3} \partial_t \mathcal{T}_{tyz} -\frac{1}{\kappa^{2}} \partial_z \dot{\mathcal{H}}_{tz} + \frac{1}{2 \, \kappa^2} \partial_t^2 \mathcal{H}_{xx} + \frac{1}{2 \, \kappa^2} \partial_t^2 \mathcal{H}_{zz} \\
\partial_t^{2} \mathcal{H}_{yz} &= \frac{B_{0}^2}{2} \mathcal{H}_{yz} -B_{0} \, \delta_{1} \partial_t \mathcal{T}_{yty} + B_{0} \, \zeta_{2} \partial_t \mathcal{T}_{yty} + \frac{B_{0} \, \zeta_{3}}{2} \partial_t \mathcal{T}_{yty} + B_{0} \, \delta_{1} \partial_t \mathcal{T}_{ztz} -B_{0} \, \zeta_{2} \partial_t \mathcal{T}_{ztz} -\frac{1}{2} B_{0} \, \zeta_{3} \partial_t \mathcal{T}_{ztz} -\frac{1}{\kappa^{2}} \partial_z \dot{\mathcal{H}}_{ty} \\
0 &= -\frac{1}{8} B_{0}^2 \mathcal{H}_{tt} + \frac{B_{0}^2}{8} \mathcal{H}_{xx} -\frac{1}{8} B_{0}^2 \mathcal{H}_{yy} -\frac{3 \, B_{0}^2}{8} \mathcal{H}_{zz} -\frac{1}{2} B_{0} \partial_z \mathcal{A}_{y} + \frac{3 \, B_{0} \, \delta_{1}}{4} \partial_z \mathcal{T}_{tty} + \frac{B_{0} \, \zeta_{1}}{2} \partial_z \mathcal{T}_{tty} + \frac{B_{0} \, \delta_{1}}{4} \partial_z \mathcal{T}_{zyz} -\frac{1}{2} B_{0} \, \zeta_{1} \partial_z \mathcal{T}_{zyz} -B_{0} \, \zeta_{2} \partial_z \mathcal{T}_{zyz} -\frac{1}{2} B_{0} \, \zeta_{3} \partial_z \mathcal{T}_{zyz} + \frac{3 \, B_{0} \, \delta_{1}}{4} \partial_z \mathcal{T}_{xxy} + \frac{B_{0} \, \zeta_{1}}{2} \partial_z \mathcal{T}_{xxy} + B_{0} \, \delta_{1} \partial_t \mathcal{T}_{zty} -B_{0} \, \zeta_{2} \partial_t \mathcal{T}_{zty} -\frac{1}{2} B_{0} \, \zeta_{3} \partial_t \mathcal{T}_{zty} + B_{0} \, \delta_{1} \partial_t \mathcal{T}_{tyz} -B_{0} \, \zeta_{2} \partial_t \mathcal{T}_{tyz} -\frac{1}{2} B_{0} \, \zeta_{3} \partial_t \mathcal{T}_{tyz} + \frac{1}{2 \, \kappa^2} \partial_t^2 \mathcal{H}_{xx} + \frac{1}{2 \, \kappa^2} \partial_t^2 \mathcal{H}_{yy} \\
\partial_t^{2} \mathcal{T}_{ttx} &= -4 \, \beta_{1} \mathcal{T}_{ttx} -2 \, \beta_{2} \mathcal{T}_{ttx} -2 \, \beta_{3} \mathcal{T}_{ttx} -\frac{1}{2 \, \kappa^2} \mathcal{T}_{ttx} -\frac{3 \, B_{0} \, \delta_{1}}{2} \mathcal{T}_{xyz} + 2 \, \beta_{3} \mathcal{T}_{yxy} -\frac{3}{\kappa^2} \mathcal{T}_{yxy} -\frac{3 \, B_{0} \, \delta_{1}}{2} \mathcal{T}_{yxz} + \frac{3 \, B_{0} \, \delta_{1}}{2} \mathcal{T}_{zxy} + 2 \, \beta_{3} \mathcal{T}_{zxz} -\frac{3}{\kappa^2} \mathcal{T}_{zxz} + \frac{3 \, B_{0} \, \delta_{1}}{2} \partial_z \mathcal{H}_{xy} + B_{0} \, \zeta_{1} \partial_z \mathcal{H}_{xy} + \frac{3 \, \delta_{1}}{2} \partial_z^2 \mathcal{A}_{x} + \zeta_{1} \partial_z^2 \mathcal{A}_{x} + 3 \, \chi \partial_z^2 \mathcal{T}_{ttx} -\xi \partial_z^2 \mathcal{T}_{ttx} -3 \, \chi \partial_z^2 \mathcal{T}_{yxy} + \xi \partial_z^2 \mathcal{T}_{yxy} -\chi \partial_z^2 \mathcal{T}_{zxz} + \xi \partial_z^2 \mathcal{T}_{zxz} + 2 \, \chi \partial_z \dot{\mathcal{T}}_{ztx} + 2 \, \chi \partial_z \dot{\mathcal{T}}_{txz} + \frac{\delta_{1}}{2} \partial_t^2 \mathcal{A}_{x} -\zeta_{1} \partial_t^2 \mathcal{A}_{x} -2 \, \zeta_{2} \partial_t^2 \mathcal{A}_{x} -\zeta_{3} \partial_t^2 \mathcal{A}_{x} + \chi \partial_t^2 \mathcal{T}_{yxy} -\xi \partial_t^2 \mathcal{T}_{yxy} + \chi \partial_t^2 \mathcal{T}_{zxz} -\xi \partial_t^2 \mathcal{T}_{zxz} \\
\partial_t^{2} \mathcal{T}_{tty} &= -4 \, \beta_{1} \mathcal{T}_{tty} -2 \, \beta_{2} \mathcal{T}_{tty} -2 \, \beta_{3} \mathcal{T}_{tty} -\frac{1}{2 \, \kappa^2} \mathcal{T}_{tty} -3 \, B_{0} \, \delta_{1} \mathcal{T}_{yyz} + 2 \, \beta_{3} \mathcal{T}_{zyz} -\frac{3}{\kappa^2} \mathcal{T}_{zyz} -2 \, \beta_{3} \mathcal{T}_{xxy} + \frac{3}{\kappa^2} \mathcal{T}_{xxy} -\frac{1}{4} B_{0} \, \delta_{1} \partial_z \mathcal{H}_{tt} + \frac{B_{0} \, \zeta_{1}}{2} \partial_z \mathcal{H}_{tt} + B_{0} \, \zeta_{2} \partial_z \mathcal{H}_{tt} + \frac{B_{0} \, \zeta_{3}}{2} \partial_z \mathcal{H}_{tt} -\frac{3 \, B_{0} \, \delta_{1}}{4} \partial_z \mathcal{H}_{xx} -\frac{1}{2} B_{0} \, \zeta_{1} \partial_z \mathcal{H}_{xx} + \frac{3 \, B_{0} \, \delta_{1}}{4} \partial_z \mathcal{H}_{yy} + \frac{B_{0} \, \zeta_{1}}{2} \partial_z \mathcal{H}_{yy} + \frac{3 \, B_{0} \, \delta_{1}}{4} \partial_z \mathcal{H}_{zz} + \frac{B_{0} \, \zeta_{1}}{2} \partial_z \mathcal{H}_{zz} + \frac{3 \, \delta_{1}}{2} \partial_z^2 \mathcal{A}_{y} + \zeta_{1} \partial_z^2 \mathcal{A}_{y} + 3 \, \chi \partial_z^2 \mathcal{T}_{tty} -\xi \partial_z^2 \mathcal{T}_{tty} -\chi \partial_z^2 \mathcal{T}_{zyz} + \xi \partial_z^2 \mathcal{T}_{zyz} + 3 \, \chi \partial_z^2 \mathcal{T}_{xxy} -\xi \partial_z^2 \mathcal{T}_{xxy} + \frac{B_{0} \, \delta_{1}}{2} \partial_t \mathcal{H}_{tz} -B_{0} \, \zeta_{1} \partial_t \mathcal{H}_{tz} -2 \, B_{0} \, \zeta_{2} \partial_t \mathcal{H}_{tz} -B_{0} \, \zeta_{3} \partial_t \mathcal{H}_{tz} + 2 \, \chi \partial_z \dot{\mathcal{T}}_{zty} + 2 \, \chi \partial_z \dot{\mathcal{T}}_{tyz} + \frac{\delta_{1}}{2} \partial_t^2 \mathcal{A}_{y} -\zeta_{1} \partial_t^2 \mathcal{A}_{y} -2 \, \zeta_{2} \partial_t^2 \mathcal{A}_{y} -\zeta_{3} \partial_t^2 \mathcal{A}_{y} + \chi \partial_t^2 \mathcal{T}_{zyz} -\xi \partial_t^2 \mathcal{T}_{zyz} -\chi \partial_t^2 \mathcal{T}_{xxy} + \xi \partial_t^2 \mathcal{T}_{xxy} \\
\partial_t^{2} \mathcal{T}_{xxz} &= -4 \, \beta_{1} \mathcal{T}_{xxz} -2 \, \beta_{2} \mathcal{T}_{xxz} -2 \, \beta_{3} \mathcal{T}_{xxz} -\frac{1}{2 \, \kappa^2} \mathcal{T}_{xxz} -2 \, \beta_{3} \mathcal{T}_{yyz} + \frac{3}{\kappa^2} \mathcal{T}_{yyz} -2 \, \beta_{3} \mathcal{T}_{ttz} + \frac{3}{\kappa^2} \mathcal{T}_{ttz} -3 \, B_{0} \, \delta_{1} \mathcal{T}_{zyz} + \frac{3 \, B_{0} \, \delta_{1}}{2} \partial_t \mathcal{H}_{ty} + B_{0} \, \zeta_{1} \partial_t \mathcal{H}_{ty} + \frac{3 \, \delta_{1}}{2} \partial_z \dot{\mathcal{A}}_{t} + \zeta_{1} \partial_z \dot{\mathcal{A}}_{t} -3 \, \chi \partial_z \dot{\mathcal{T}}_{yty} + \xi \partial_z \dot{\mathcal{T}}_{yty} -\chi \partial_z \dot{\mathcal{T}}_{ztz} + \xi \partial_z \dot{\mathcal{T}}_{ztz} -3 \, \chi \partial_z \dot{\mathcal{T}}_{xtx} + \xi \partial_z \dot{\mathcal{T}}_{xtx} -\frac{3 \, \delta_{1}}{2} \partial_t^2 \mathcal{A}_{z} -\zeta_{1} \partial_t^2 \mathcal{A}_{z} -3 \, \chi \partial_t^2 \mathcal{T}_{yyz} + \xi \partial_t^2 \mathcal{T}_{yyz} -\chi \partial_t^2 \mathcal{T}_{ttz} + \xi \partial_t^2 \mathcal{T}_{ttz} \\
0 &= \frac{3 \, B_{0} \, \delta_{1}}{2} \mathcal{T}_{ttx} + 4 \, \beta_{1} \mathcal{T}_{xyz} + \frac{2}{\kappa^2} \mathcal{T}_{xyz} -2 \, B_{0} \, \delta_{1} \mathcal{T}_{yxy} + 2 \, \beta_{2} \mathcal{T}_{yxz} + \frac{3}{2 \, \kappa^2} \mathcal{T}_{yxz} -2 \, \beta_{2} \mathcal{T}_{zxy} -\frac{3}{2 \, \kappa^2} \mathcal{T}_{zxy} -2 \, B_{0} \, \delta_{1} \mathcal{T}_{zxz} \\
0 &= -4 \, \beta_{1} \mathcal{T}_{ytx} -\frac{2}{\kappa^2} \mathcal{T}_{ytx} -2 \, \beta_{2} \mathcal{T}_{txy} -\frac{3}{2 \, \kappa^2} \mathcal{T}_{txy} -\frac{1}{2} B_{0} \, \delta_{1} \mathcal{T}_{txz} -2 \, \beta_{2} \mathcal{T}_{xty} -\frac{3}{2 \, \kappa^2} \mathcal{T}_{xty} -\frac{1}{2} B_{0} \, \delta_{1} \mathcal{T}_{xtz} \\
0 &= -4 \, \beta_{1} \mathcal{T}_{yty} -2 \, \beta_{2} \mathcal{T}_{yty} -2 \, \beta_{3} \mathcal{T}_{yty} -\frac{1}{2 \, \kappa^2} \mathcal{T}_{yty} + B_{0} \, \delta_{1} \mathcal{T}_{ytz} -B_{0} \, \delta_{1} \mathcal{T}_{zty} -2 \, \beta_{3} \mathcal{T}_{ztz} + \frac{3}{\kappa^2} \mathcal{T}_{ztz} -2 \, B_{0} \, \delta_{1} \mathcal{T}_{tyz} -2 \, \beta_{3} \mathcal{T}_{xtx} + \frac{3}{\kappa^2} \mathcal{T}_{xtx} -\frac{1}{2} B_{0} \, \delta_{1} \partial_z \mathcal{H}_{ty} -B_{0} \, \zeta_{1} \partial_z \mathcal{H}_{ty} -B_{0} \, \zeta_{2} \partial_z \mathcal{H}_{ty} -\frac{1}{2} B_{0} \, \zeta_{3} \partial_z \mathcal{H}_{ty} -\frac{3 \, \delta_{1}}{2} \partial_z^2 \mathcal{A}_{t} -\zeta_{1} \partial_z^2 \mathcal{A}_{t} + 3 \, \chi \partial_z^2 \mathcal{T}_{yty} -\xi \partial_z^2 \mathcal{T}_{yty} + \chi \partial_z^2 \mathcal{T}_{ztz} -\xi \partial_z^2 \mathcal{T}_{ztz} + 3 \, \chi \partial_z^2 \mathcal{T}_{xtx} -\xi \partial_z^2 \mathcal{T}_{xtx} -B_{0} \, \delta_{1} \partial_t \mathcal{H}_{yz} + B_{0} \, \zeta_{2} \partial_t \mathcal{H}_{yz} + \frac{B_{0} \, \zeta_{3}}{2} \partial_t \mathcal{H}_{yz} + \frac{3 \, \delta_{1}}{2} \partial_z \dot{\mathcal{A}}_{z} + \zeta_{1} \partial_z \dot{\mathcal{A}}_{z} + 3 \, \chi \partial_z \dot{\mathcal{T}}_{xxz} -\xi \partial_z \dot{\mathcal{T}}_{xxz} + 3 \, \chi \partial_z \dot{\mathcal{T}}_{yyz} -\xi \partial_z \dot{\mathcal{T}}_{yyz} + \chi \partial_z \dot{\mathcal{T}}_{ttz} -\xi \partial_z \dot{\mathcal{T}}_{ttz} \\
0 &= B_{0} \, \delta_{1} \mathcal{T}_{yty} -4 \, \beta_{1} \mathcal{T}_{ytz} -\frac{2}{\kappa^2} \mathcal{T}_{ytz} -2 \, \beta_{2} \mathcal{T}_{zty} -\frac{3}{2 \, \kappa^2} \mathcal{T}_{zty} + B_{0} \, \delta_{1} \mathcal{T}_{ztz} + 2 \, \beta_{2} \mathcal{T}_{tyz} + \frac{3}{2 \, \kappa^2} \mathcal{T}_{tyz} + \frac{3 \, B_{0} \, \delta_{1}}{2} \mathcal{T}_{xtx} + B_{0} \, \delta_{1} \partial_t \mathcal{H}_{yy} -B_{0} \, \zeta_{2} \partial_t \mathcal{H}_{yy} -\frac{1}{2} B_{0} \, \zeta_{3} \partial_t \mathcal{H}_{yy} \\
\partial_t^{2} \mathcal{T}_{yxy} &= 2 \, \beta_{3} \mathcal{T}_{ttx} -\frac{3}{\kappa^2} \mathcal{T}_{ttx} + 2 \, B_{0} \, \delta_{1} \mathcal{T}_{xyz} -4 \, \beta_{1} \mathcal{T}_{yxy} -2 \, \beta_{2} \mathcal{T}_{yxy} -2 \, \beta_{3} \mathcal{T}_{yxy} -\frac{1}{2 \, \kappa^2} \mathcal{T}_{yxy} + B_{0} \, \delta_{1} \mathcal{T}_{yxz} -B_{0} \, \delta_{1} \mathcal{T}_{zxy} -2 \, \beta_{3} \mathcal{T}_{zxz} + \frac{3}{\kappa^2} \mathcal{T}_{zxz} -\frac{1}{2} B_{0} \, \delta_{1} \partial_z \mathcal{H}_{xy} -B_{0} \, \zeta_{1} \partial_z \mathcal{H}_{xy} -B_{0} \, \zeta_{2} \partial_z \mathcal{H}_{xy} -\frac{1}{2} B_{0} \, \zeta_{3} \partial_z \mathcal{H}_{xy} -\frac{3 \, \delta_{1}}{2} \partial_z^2 \mathcal{A}_{x} -\zeta_{1} \partial_z^2 \mathcal{A}_{x} -3 \, \chi \partial_z^2 \mathcal{T}_{ttx} + \xi \partial_z^2 \mathcal{T}_{ttx} + 3 \, \chi \partial_z^2 \mathcal{T}_{yxy} -\xi \partial_z^2 \mathcal{T}_{yxy} + \chi \partial_z^2 \mathcal{T}_{zxz} -\xi \partial_z^2 \mathcal{T}_{zxz} -2 \, \chi \partial_z \dot{\mathcal{T}}_{ztx} -2 \, \chi \partial_z \dot{\mathcal{T}}_{txz} + \frac{3 \, \delta_{1}}{2} \partial_t^2 \mathcal{A}_{x} + \zeta_{1} \partial_t^2 \mathcal{A}_{x} + \chi \partial_t^2 \mathcal{T}_{ttx} -\xi \partial_t^2 \mathcal{T}_{ttx} -3 \, \chi \partial_t^2 \mathcal{T}_{zxz} + \xi \partial_t^2 \mathcal{T}_{zxz} \\
0 &= \frac{3 \, B_{0} \, \delta_{1}}{2} \mathcal{T}_{ttx} + 2 \, \beta_{2} \mathcal{T}_{xyz} + \frac{3}{2 \, \kappa^2} \mathcal{T}_{xyz} -B_{0} \, \delta_{1} \mathcal{T}_{yxy} + 4 \, \beta_{1} \mathcal{T}_{yxz} + \frac{2}{\kappa^2} \mathcal{T}_{yxz} + 2 \, \beta_{2} \mathcal{T}_{zxy} + \frac{3}{2 \, \kappa^2} \mathcal{T}_{zxy} -B_{0} \, \delta_{1} \mathcal{T}_{zxz} \\
\partial_t^{2} \mathcal{T}_{yyz} &= -3 \, B_{0} \, \delta_{1} \mathcal{T}_{tty} -2 \, \beta_{3} \mathcal{T}_{xxz} + \frac{3}{\kappa^2} \mathcal{T}_{xxz} -4 \, \beta_{1} \mathcal{T}_{yyz} -2 \, \beta_{2} \mathcal{T}_{yyz} -2 \, \beta_{3} \mathcal{T}_{yyz} -\frac{1}{2 \, \kappa^2} \mathcal{T}_{yyz} -2 \, \beta_{3} \mathcal{T}_{ttz} + \frac{3}{\kappa^2} \mathcal{T}_{ttz} -3 \, B_{0} \, \delta_{1} \mathcal{T}_{xxy} + \frac{3 \, B_{0} \, \delta_{1}}{2} \partial_t \mathcal{H}_{ty} + B_{0} \, \zeta_{1} \partial_t \mathcal{H}_{ty} + \frac{3 \, \delta_{1}}{2} \partial_z \dot{\mathcal{A}}_{t} + \zeta_{1} \partial_z \dot{\mathcal{A}}_{t} -3 \, \chi \partial_z \dot{\mathcal{T}}_{yty} + \xi \partial_z \dot{\mathcal{T}}_{yty} -\chi \partial_z \dot{\mathcal{T}}_{ztz} + \xi \partial_z \dot{\mathcal{T}}_{ztz} -3 \, \chi \partial_z \dot{\mathcal{T}}_{xtx} + \xi \partial_z \dot{\mathcal{T}}_{xtx} -\frac{3 \, \delta_{1}}{2} \partial_t^2 \mathcal{A}_{z} -\zeta_{1} \partial_t^2 \mathcal{A}_{z} -3 \, \chi \partial_t^2 \mathcal{T}_{xxz} + \xi \partial_t^2 \mathcal{T}_{xxz} -\chi \partial_t^2 \mathcal{T}_{ttz} + \xi \partial_t^2 \mathcal{T}_{ttz} \\
0 &= -4 \, \beta_{1} \mathcal{T}_{ztx} -\frac{2}{\kappa^2} \mathcal{T}_{ztx} + \frac{B_{0} \, \delta_{1}}{2} \mathcal{T}_{txy} -2 \, \beta_{2} \mathcal{T}_{txz} -\frac{3}{2 \, \kappa^2} \mathcal{T}_{txz} + \frac{B_{0} \, \delta_{1}}{2} \mathcal{T}_{xty} -2 \, \beta_{2} \mathcal{T}_{xtz} -\frac{3}{2 \, \kappa^2} \mathcal{T}_{xtz} -2 \, \delta_{1} \partial_z \dot{\mathcal{A}}_{x} + 2 \, \zeta_{2} \partial_z \dot{\mathcal{A}}_{x} + \zeta_{3} \partial_z \dot{\mathcal{A}}_{x} -2 \, \chi \partial_z \dot{\mathcal{T}}_{ttx} + 2 \, \chi \partial_z \dot{\mathcal{T}}_{yxy} + 2 \, \chi \partial_z \dot{\mathcal{T}}_{zxz} \\
0 &= -B_{0} \, \delta_{1} \mathcal{T}_{yty} -2 \, \beta_{2} \mathcal{T}_{ytz} -\frac{3}{2 \, \kappa^2} \mathcal{T}_{ytz} -4 \, \beta_{1} \mathcal{T}_{zty} -\frac{2}{\kappa^2} \mathcal{T}_{zty} -B_{0} \, \delta_{1} \mathcal{T}_{ztz} -2 \, \beta_{2} \mathcal{T}_{tyz} -\frac{3}{2 \, \kappa^2} \mathcal{T}_{tyz} -\frac{3 \, B_{0} \, \delta_{1}}{2} \mathcal{T}_{xtx} -B_{0} \, \delta_{1} \partial_t \mathcal{H}_{zz} + B_{0} \, \zeta_{2} \partial_t \mathcal{H}_{zz} + \frac{B_{0} \, \zeta_{3}}{2} \partial_t \mathcal{H}_{zz} -2 \, \delta_{1} \partial_z \dot{\mathcal{A}}_{y} + 2 \, \zeta_{2} \partial_z \dot{\mathcal{A}}_{y} + \zeta_{3} \partial_z \dot{\mathcal{A}}_{y} -2 \, \chi \partial_z \dot{\mathcal{T}}_{tty} + 2 \, \chi \partial_z \dot{\mathcal{T}}_{zyz} -2 \, \chi \partial_z \dot{\mathcal{T}}_{xxy} \\
\partial_t^{2} \mathcal{T}_{ttz} &= -2 \, \beta_{3} \mathcal{T}_{xxz} + \frac{3}{\kappa^2} \mathcal{T}_{xxz} -2 \, \beta_{3} \mathcal{T}_{yyz} + \frac{3}{\kappa^2} \mathcal{T}_{yyz} -4 \, \beta_{1} \mathcal{T}_{ttz} -2 \, \beta_{2} \mathcal{T}_{ttz} -2 \, \beta_{3} \mathcal{T}_{ttz} -\frac{1}{2 \, \kappa^2} \mathcal{T}_{ttz} -3 \, B_{0} \, \delta_{1} \mathcal{T}_{zyz} -\frac{1}{2} B_{0} \, \delta_{1} \partial_t \mathcal{H}_{ty} + B_{0} \, \zeta_{1} \partial_t \mathcal{H}_{ty} + 2 \, B_{0} \, \zeta_{2} \partial_t \mathcal{H}_{ty} + B_{0} \, \zeta_{3} \partial_t \mathcal{H}_{ty} -\frac{1}{2} \delta_{1} \partial_z \dot{\mathcal{A}}_{t} + \zeta_{1} \partial_z \dot{\mathcal{A}}_{t} + 2 \, \zeta_{2} \partial_z \dot{\mathcal{A}}_{t} + \zeta_{3} \partial_z \dot{\mathcal{A}}_{t} -\chi \partial_z \dot{\mathcal{T}}_{yty} + \xi \partial_z \dot{\mathcal{T}}_{yty} + \chi \partial_z \dot{\mathcal{T}}_{ztz} + \xi \partial_z \dot{\mathcal{T}}_{ztz} -\chi \partial_z \dot{\mathcal{T}}_{xtx} + \xi \partial_z \dot{\mathcal{T}}_{xtx} + \frac{\delta_{1}}{2} \partial_t^2 \mathcal{A}_{z} -\zeta_{1} \partial_t^2 \mathcal{A}_{z} -2 \, \zeta_{2} \partial_t^2 \mathcal{A}_{z} -\zeta_{3} \partial_t^2 \mathcal{A}_{z} -\chi \partial_t^2 \mathcal{T}_{xxz} + \xi \partial_t^2 \mathcal{T}_{xxz} -\chi \partial_t^2 \mathcal{T}_{yyz} + \xi \partial_t^2 \mathcal{T}_{yyz} \\
0 &= -2 \, \beta_{3} \mathcal{T}_{yty} + \frac{3}{\kappa^2} \mathcal{T}_{yty} + B_{0} \, \delta_{1} \mathcal{T}_{ytz} -B_{0} \, \delta_{1} \mathcal{T}_{zty} -4 \, \beta_{1} \mathcal{T}_{ztz} -2 \, \beta_{2} \mathcal{T}_{ztz} -2 \, \beta_{3} \mathcal{T}_{ztz} -\frac{1}{2 \, \kappa^2} \mathcal{T}_{ztz} -2 \, B_{0} \, \delta_{1} \mathcal{T}_{tyz} -2 \, \beta_{3} \mathcal{T}_{xtx} + \frac{3}{\kappa^2} \mathcal{T}_{xtx} -\frac{1}{2} B_{0} \, \delta_{1} \partial_z \mathcal{H}_{ty} -B_{0} \, \zeta_{1} \partial_z \mathcal{H}_{ty} -B_{0} \, \zeta_{2} \partial_z \mathcal{H}_{ty} -\frac{1}{2} B_{0} \, \zeta_{3} \partial_z \mathcal{H}_{ty} + \frac{\delta_{1}}{2} \partial_z^2 \mathcal{A}_{t} -\zeta_{1} \partial_z^2 \mathcal{A}_{t} -2 \, \zeta_{2} \partial_z^2 \mathcal{A}_{t} -\zeta_{3} \partial_z^2 \mathcal{A}_{t} + \chi \partial_z^2 \mathcal{T}_{yty} -\xi \partial_z^2 \mathcal{T}_{yty} -\chi \partial_z^2 \mathcal{T}_{ztz} -\xi \partial_z^2 \mathcal{T}_{ztz} + \chi \partial_z^2 \mathcal{T}_{xtx} -\xi \partial_z^2 \mathcal{T}_{xtx} + B_{0} \, \delta_{1} \partial_t \mathcal{H}_{yz} -B_{0} \, \zeta_{2} \partial_t \mathcal{H}_{yz} -\frac{1}{2} B_{0} \, \zeta_{3} \partial_t \mathcal{H}_{yz} -\frac{1}{2} \delta_{1} \partial_z \dot{\mathcal{A}}_{z} + \zeta_{1} \partial_z \dot{\mathcal{A}}_{z} + 2 \, \zeta_{2} \partial_z \dot{\mathcal{A}}_{z} + \zeta_{3} \partial_z \dot{\mathcal{A}}_{z} + \chi \partial_z \dot{\mathcal{T}}_{xxz} -\xi \partial_z \dot{\mathcal{T}}_{xxz} + \chi \partial_z \dot{\mathcal{T}}_{yyz} -\xi \partial_z \dot{\mathcal{T}}_{yyz} -\chi \partial_z \dot{\mathcal{T}}_{ttz} -\xi \partial_z \dot{\mathcal{T}}_{ttz} \\
0 &= -\frac{3 \, B_{0} \, \delta_{1}}{2} \mathcal{T}_{ttx} -2 \, \beta_{2} \mathcal{T}_{xyz} -\frac{3}{2 \, \kappa^2} \mathcal{T}_{xyz} + B_{0} \, \delta_{1} \mathcal{T}_{yxy} + 2 \, \beta_{2} \mathcal{T}_{yxz} + \frac{3}{2 \, \kappa^2} \mathcal{T}_{yxz} + 4 \, \beta_{1} \mathcal{T}_{zxy} + \frac{2}{\kappa^2} \mathcal{T}_{zxy} + B_{0} \, \delta_{1} \mathcal{T}_{zxz} \\
\partial_t^{2} \mathcal{T}_{zxz} &= 2 \, \beta_{3} \mathcal{T}_{ttx} -\frac{3}{\kappa^2} \mathcal{T}_{ttx} + 2 \, B_{0} \, \delta_{1} \mathcal{T}_{xyz} -2 \, \beta_{3} \mathcal{T}_{yxy} + \frac{3}{\kappa^2} \mathcal{T}_{yxy} + B_{0} \, \delta_{1} \mathcal{T}_{yxz} -B_{0} \, \delta_{1} \mathcal{T}_{zxy} -4 \, \beta_{1} \mathcal{T}_{zxz} -2 \, \beta_{2} \mathcal{T}_{zxz} -2 \, \beta_{3} \mathcal{T}_{zxz} -\frac{1}{2 \, \kappa^2} \mathcal{T}_{zxz} -\frac{1}{2} B_{0} \, \delta_{1} \partial_z \mathcal{H}_{xy} -B_{0} \, \zeta_{1} \partial_z \mathcal{H}_{xy} -B_{0} \, \zeta_{2} \partial_z \mathcal{H}_{xy} -\frac{1}{2} B_{0} \, \zeta_{3} \partial_z \mathcal{H}_{xy} + \frac{\delta_{1}}{2} \partial_z^2 \mathcal{A}_{x} -\zeta_{1} \partial_z^2 \mathcal{A}_{x} -2 \, \zeta_{2} \partial_z^2 \mathcal{A}_{x} -\zeta_{3} \partial_z^2 \mathcal{A}_{x} -\chi \partial_z^2 \mathcal{T}_{ttx} + \xi \partial_z^2 \mathcal{T}_{ttx} + \chi \partial_z^2 \mathcal{T}_{yxy} -\xi \partial_z^2 \mathcal{T}_{yxy} -\chi \partial_z^2 \mathcal{T}_{zxz} -\xi \partial_z^2 \mathcal{T}_{zxz} -2 \, \chi \partial_z \dot{\mathcal{T}}_{ztx} -2 \, \chi \partial_z \dot{\mathcal{T}}_{txz} + \frac{3 \, \delta_{1}}{2} \partial_t^2 \mathcal{A}_{x} + \zeta_{1} \partial_t^2 \mathcal{A}_{x} + \chi \partial_t^2 \mathcal{T}_{ttx} -\xi \partial_t^2 \mathcal{T}_{ttx} -3 \, \chi \partial_t^2 \mathcal{T}_{yxy} + \xi \partial_t^2 \mathcal{T}_{yxy} \\
\partial_t^{2} \mathcal{T}_{zyz} &= 2 \, \beta_{3} \mathcal{T}_{tty} -\frac{3}{\kappa^2} \mathcal{T}_{tty} -3 \, B_{0} \, \delta_{1} \mathcal{T}_{xxz} -3 \, B_{0} \, \delta_{1} \mathcal{T}_{ttz} -4 \, \beta_{1} \mathcal{T}_{zyz} -2 \, \beta_{2} \mathcal{T}_{zyz} -2 \, \beta_{3} \mathcal{T}_{zyz} -\frac{1}{2 \, \kappa^2} \mathcal{T}_{zyz} + 2 \, \beta_{3} \mathcal{T}_{xxy} -\frac{3}{\kappa^2} \mathcal{T}_{xxy} -\frac{3 \, B_{0} \, \delta_{1}}{4} \partial_z \mathcal{H}_{tt} -\frac{1}{2} B_{0} \, \zeta_{1} \partial_z \mathcal{H}_{tt} + \frac{3 \, B_{0} \, \delta_{1}}{4} \partial_z \mathcal{H}_{xx} + \frac{B_{0} \, \zeta_{1}}{2} \partial_z \mathcal{H}_{xx} + \frac{B_{0} \, \delta_{1}}{4} \partial_z \mathcal{H}_{yy} -\frac{1}{2} B_{0} \, \zeta_{1} \partial_z \mathcal{H}_{yy} -B_{0} \, \zeta_{2} \partial_z \mathcal{H}_{yy} -\frac{1}{2} B_{0} \, \zeta_{3} \partial_z \mathcal{H}_{yy} + \frac{B_{0} \, \delta_{1}}{4} \partial_z \mathcal{H}_{zz} -\frac{1}{2} B_{0} \, \zeta_{1} \partial_z \mathcal{H}_{zz} -B_{0} \, \zeta_{2} \partial_z \mathcal{H}_{zz} -\frac{1}{2} B_{0} \, \zeta_{3} \partial_z \mathcal{H}_{zz} + \frac{\delta_{1}}{2} \partial_z^2 \mathcal{A}_{y} -\zeta_{1} \partial_z^2 \mathcal{A}_{y} -2 \, \zeta_{2} \partial_z^2 \mathcal{A}_{y} -\zeta_{3} \partial_z^2 \mathcal{A}_{y} -\chi \partial_z^2 \mathcal{T}_{tty} + \xi \partial_z^2 \mathcal{T}_{tty} -\chi \partial_z^2 \mathcal{T}_{zyz} -\xi \partial_z^2 \mathcal{T}_{zyz} -\chi \partial_z^2 \mathcal{T}_{xxy} + \xi \partial_z^2 \mathcal{T}_{xxy} + \frac{3 \, B_{0} \, \delta_{1}}{2} \partial_t \mathcal{H}_{tz} + B_{0} \, \zeta_{1} \partial_t \mathcal{H}_{tz} -2 \, \chi \partial_z \dot{\mathcal{T}}_{zty} -2 \, \chi \partial_z \dot{\mathcal{T}}_{tyz} + \frac{3 \, \delta_{1}}{2} \partial_t^2 \mathcal{A}_{y} + \zeta_{1} \partial_t^2 \mathcal{A}_{y} + \chi \partial_t^2 \mathcal{T}_{tty} -\xi \partial_t^2 \mathcal{T}_{tty} + 3 \, \chi \partial_t^2 \mathcal{T}_{xxy} -\xi \partial_t^2 \mathcal{T}_{xxy} \\
0 &= -2 \, \beta_{2} \mathcal{T}_{ytx} -\frac{3}{2 \, \kappa^2} \mathcal{T}_{ytx} + \frac{B_{0} \, \delta_{1}}{2} \mathcal{T}_{ztx} -4 \, \beta_{1} \mathcal{T}_{txy} -\frac{2}{\kappa^2} \mathcal{T}_{txy} + 2 \, \beta_{2} \mathcal{T}_{xty} + \frac{3}{2 \, \kappa^2} \mathcal{T}_{xty} -B_{0} \, \delta_{1} \partial_z \mathcal{H}_{tx} + B_{0} \, \zeta_{2} \partial_z \mathcal{H}_{tx} + \frac{B_{0} \, \zeta_{3}}{2} \partial_z \mathcal{H}_{tx} + B_{0} \, \delta_{1} \partial_t \mathcal{H}_{xz} -B_{0} \, \zeta_{2} \partial_t \mathcal{H}_{xz} -\frac{1}{2} B_{0} \, \zeta_{3} \partial_t \mathcal{H}_{xz} \\
0 &= -\frac{1}{2} B_{0} \, \delta_{1} \mathcal{T}_{ytx} -2 \, \beta_{2} \mathcal{T}_{ztx} -\frac{3}{2 \, \kappa^2} \mathcal{T}_{ztx} -4 \, \beta_{1} \mathcal{T}_{txz} -\frac{2}{\kappa^2} \mathcal{T}_{txz} + 2 \, \beta_{2} \mathcal{T}_{xtz} + \frac{3}{2 \, \kappa^2} \mathcal{T}_{xtz} -B_{0} \, \delta_{1} \partial_t \mathcal{H}_{xy} + B_{0} \, \zeta_{2} \partial_t \mathcal{H}_{xy} + \frac{B_{0} \, \zeta_{3}}{2} \partial_t \mathcal{H}_{xy} -2 \, \delta_{1} \partial_z \dot{\mathcal{A}}_{x} + 2 \, \zeta_{2} \partial_z \dot{\mathcal{A}}_{x} + \zeta_{3} \partial_z \dot{\mathcal{A}}_{x} -2 \, \chi \partial_z \dot{\mathcal{T}}_{ttx} + 2 \, \chi \partial_z \dot{\mathcal{T}}_{yxy} + 2 \, \chi \partial_z \dot{\mathcal{T}}_{zxz} \\
0 &= -2 \, B_{0} \, \delta_{1} \mathcal{T}_{yty} + 2 \, \beta_{2} \mathcal{T}_{ytz} + \frac{3}{2 \, \kappa^2} \mathcal{T}_{ytz} -2 \, \beta_{2} \mathcal{T}_{zty} -\frac{3}{2 \, \kappa^2} \mathcal{T}_{zty} -2 \, B_{0} \, \delta_{1} \mathcal{T}_{ztz} -4 \, \beta_{1} \mathcal{T}_{tyz} -\frac{2}{\kappa^2} \mathcal{T}_{tyz} -\frac{3 \, B_{0} \, \delta_{1}}{2} \mathcal{T}_{xtx} -B_{0} \, \delta_{1} \partial_t \mathcal{H}_{yy} + B_{0} \, \zeta_{2} \partial_t \mathcal{H}_{yy} + \frac{B_{0} \, \zeta_{3}}{2} \partial_t \mathcal{H}_{yy} -B_{0} \, \delta_{1} \partial_t \mathcal{H}_{zz} + B_{0} \, \zeta_{2} \partial_t \mathcal{H}_{zz} + \frac{B_{0} \, \zeta_{3}}{2} \partial_t \mathcal{H}_{zz} -2 \, \delta_{1} \partial_z \dot{\mathcal{A}}_{y} + 2 \, \zeta_{2} \partial_z \dot{\mathcal{A}}_{y} + \zeta_{3} \partial_z \dot{\mathcal{A}}_{y} -2 \, \chi \partial_z \dot{\mathcal{T}}_{tty} + 2 \, \chi \partial_z \dot{\mathcal{T}}_{zyz} -2 \, \chi \partial_z \dot{\mathcal{T}}_{xxy} \\
0 &= -2 \, \beta_{3} \mathcal{T}_{yty} + \frac{3}{\kappa^2} \mathcal{T}_{yty} + \frac{3 \, B_{0} \, \delta_{1}}{2} \mathcal{T}_{ytz} -\frac{3 \, B_{0} \, \delta_{1}}{2} \mathcal{T}_{zty} -2 \, \beta_{3} \mathcal{T}_{ztz} + \frac{3}{\kappa^2} \mathcal{T}_{ztz} -\frac{3 \, B_{0} \, \delta_{1}}{2} \mathcal{T}_{tyz} -4 \, \beta_{1} \mathcal{T}_{xtx} -2 \, \beta_{2} \mathcal{T}_{xtx} -2 \, \beta_{3} \mathcal{T}_{xtx} -\frac{1}{2 \, \kappa^2} \mathcal{T}_{xtx} -\frac{3 \, B_{0} \, \delta_{1}}{2} \partial_z \mathcal{H}_{ty} -B_{0} \, \zeta_{1} \partial_z \mathcal{H}_{ty} -\frac{3 \, \delta_{1}}{2} \partial_z^2 \mathcal{A}_{t} -\zeta_{1} \partial_z^2 \mathcal{A}_{t} + 3 \, \chi \partial_z^2 \mathcal{T}_{yty} -\xi \partial_z^2 \mathcal{T}_{yty} + \chi \partial_z^2 \mathcal{T}_{ztz} -\xi \partial_z^2 \mathcal{T}_{ztz} + 3 \, \chi \partial_z^2 \mathcal{T}_{xtx} -\xi \partial_z^2 \mathcal{T}_{xtx} + \frac{3 \, \delta_{1}}{2} \partial_z \dot{\mathcal{A}}_{z} + \zeta_{1} \partial_z \dot{\mathcal{A}}_{z} + 3 \, \chi \partial_z \dot{\mathcal{T}}_{xxz} -\xi \partial_z \dot{\mathcal{T}}_{xxz} + 3 \, \chi \partial_z \dot{\mathcal{T}}_{yyz} -\xi \partial_z \dot{\mathcal{T}}_{yyz} + \chi \partial_z \dot{\mathcal{T}}_{ttz} -\xi \partial_z \dot{\mathcal{T}}_{ttz} \\
0 &= -2 \, \beta_{2} \mathcal{T}_{ytx} -\frac{3}{2 \, \kappa^2} \mathcal{T}_{ytx} + \frac{B_{0} \, \delta_{1}}{2} \mathcal{T}_{ztx} + 2 \, \beta_{2} \mathcal{T}_{txy} + \frac{3}{2 \, \kappa^2} \mathcal{T}_{txy} -4 \, \beta_{1} \mathcal{T}_{xty} -\frac{2}{\kappa^2} \mathcal{T}_{xty} + B_{0} \, \delta_{1} \partial_z \mathcal{H}_{tx} -B_{0} \, \zeta_{2} \partial_z \mathcal{H}_{tx} -\frac{1}{2} B_{0} \, \zeta_{3} \partial_z \mathcal{H}_{tx} -B_{0} \, \delta_{1} \partial_t \mathcal{H}_{xz} + B_{0} \, \zeta_{2} \partial_t \mathcal{H}_{xz} + \frac{B_{0} \, \zeta_{3}}{2} \partial_t \mathcal{H}_{xz} \\
0 &= -\frac{1}{2} B_{0} \, \delta_{1} \mathcal{T}_{ytx} -2 \, \beta_{2} \mathcal{T}_{ztx} -\frac{3}{2 \, \kappa^2} \mathcal{T}_{ztx} + 2 \, \beta_{2} \mathcal{T}_{txz} + \frac{3}{2 \, \kappa^2} \mathcal{T}_{txz} -4 \, \beta_{1} \mathcal{T}_{xtz} -\frac{2}{\kappa^2} \mathcal{T}_{xtz} + B_{0} \, \delta_{1} \partial_t \mathcal{H}_{xy} -B_{0} \, \zeta_{2} \partial_t \mathcal{H}_{xy} -\frac{1}{2} B_{0} \, \zeta_{3} \partial_t \mathcal{H}_{xy} \\
\partial_t^{2} \mathcal{T}_{xxy} &= -2 \, \beta_{3} \mathcal{T}_{tty} + \frac{3}{\kappa^2} \mathcal{T}_{tty} -3 \, B_{0} \, \delta_{1} \mathcal{T}_{yyz} + 2 \, \beta_{3} \mathcal{T}_{zyz} -\frac{3}{\kappa^2} \mathcal{T}_{zyz} -4 \, \beta_{1} \mathcal{T}_{xxy} -2 \, \beta_{2} \mathcal{T}_{xxy} -2 \, \beta_{3} \mathcal{T}_{xxy} -\frac{1}{2 \, \kappa^2} \mathcal{T}_{xxy} + \frac{3 \, B_{0} \, \delta_{1}}{4} \partial_z \mathcal{H}_{tt} + \frac{B_{0} \, \zeta_{1}}{2} \partial_z \mathcal{H}_{tt} + \frac{B_{0} \, \delta_{1}}{4} \partial_z \mathcal{H}_{xx} -\frac{1}{2} B_{0} \, \zeta_{1} \partial_z \mathcal{H}_{xx} -B_{0} \, \zeta_{2} \partial_z \mathcal{H}_{xx} -\frac{1}{2} B_{0} \, \zeta_{3} \partial_z \mathcal{H}_{xx} + \frac{3 \, B_{0} \, \delta_{1}}{4} \partial_z \mathcal{H}_{yy} + \frac{B_{0} \, \zeta_{1}}{2} \partial_z \mathcal{H}_{yy} + \frac{3 \, B_{0} \, \delta_{1}}{4} \partial_z \mathcal{H}_{zz} + \frac{B_{0} \, \zeta_{1}}{2} \partial_z \mathcal{H}_{zz} + \frac{3 \, \delta_{1}}{2} \partial_z^2 \mathcal{A}_{y} + \zeta_{1} \partial_z^2 \mathcal{A}_{y} + 3 \, \chi \partial_z^2 \mathcal{T}_{tty} -\xi \partial_z^2 \mathcal{T}_{tty} -\chi \partial_z^2 \mathcal{T}_{zyz} + \xi \partial_z^2 \mathcal{T}_{zyz} + 3 \, \chi \partial_z^2 \mathcal{T}_{xxy} -\xi \partial_z^2 \mathcal{T}_{xxy} -\frac{3 \, B_{0} \, \delta_{1}}{2} \partial_t \mathcal{H}_{tz} -B_{0} \, \zeta_{1} \partial_t \mathcal{H}_{tz} + 2 \, \chi \partial_z \dot{\mathcal{T}}_{zty} + 2 \, \chi \partial_z \dot{\mathcal{T}}_{tyz} -\frac{3 \, \delta_{1}}{2} \partial_t^2 \mathcal{A}_{y} -\zeta_{1} \partial_t^2 \mathcal{A}_{y} -\chi \partial_t^2 \mathcal{T}_{tty} + \xi \partial_t^2 \mathcal{T}_{tty} + 3 \, \chi \partial_t^2 \mathcal{T}_{zyz} -\xi \partial_t^2 \mathcal{T}_{zyz} \\
\mathscr{H} &= -\frac{1}{16} B_{0}^2 \mathcal{H}_{tt}^2 + \frac{B_{0}^2}{4} \mathcal{H}_{tx}^2 -\frac{1}{4} B_{0}^2 \mathcal{H}_{ty}^2 -\frac{1}{4} B_{0}^2 \mathcal{H}_{tz}^2 -\frac{1}{8} B_{0}^2 \mathcal{H}_{tt} \, \mathcal{H}_{xx} -\frac{1}{16} B_{0}^2 \mathcal{H}_{xx}^2 + \frac{B_{0}^2}{4} \mathcal{H}_{xy}^2 + \frac{B_{0}^2}{4} \mathcal{H}_{xz}^2 + \frac{B_{0}^2}{8} \mathcal{H}_{tt} \, \mathcal{H}_{yy} -\frac{1}{8} B_{0}^2 \mathcal{H}_{xx} \, \mathcal{H}_{yy} + \frac{3 \, B_{0}^2}{16} \mathcal{H}_{yy}^2 + \frac{B_{0}^2}{4} \mathcal{H}_{yz}^2 + \frac{B_{0}^2}{8} \mathcal{H}_{tt} \, \mathcal{H}_{zz} -\frac{1}{8} B_{0}^2 \mathcal{H}_{xx} \, \mathcal{H}_{zz} + \frac{B_{0}^2}{8} \mathcal{H}_{yy} \, \mathcal{H}_{zz} + \frac{3 \, B_{0}^2}{16} \mathcal{H}_{zz}^2 -\frac{1}{2} \partial_z \mathcal{A}_{t}^2 + \frac{1}{2} \partial_z \mathcal{A}_{x}^2 + \frac{1}{2} \partial_z \mathcal{A}_{y}^2 + \frac{1}{2} \partial_z \mathcal{A}_{z}^2 -\frac{1}{2 \, \kappa^2} \partial_z \mathcal{H}_{tx}^2 -\frac{1}{2 \, \kappa^2} \partial_z \mathcal{H}_{ty}^2 + \frac{1}{2 \, \kappa^2} \partial_z \mathcal{H}_{tt} \, \partial_z \mathcal{H}_{xx} + \frac{1}{2 \, \kappa^2} \partial_z \mathcal{H}_{xy}^2 + \frac{1}{2 \, \kappa^2} \partial_z \mathcal{H}_{tt} \, \partial_z \mathcal{H}_{yy} -\frac{1}{2 \, \kappa^2} \partial_z \mathcal{H}_{xx} \, \partial_z \mathcal{H}_{yy} -\frac{1}{2} \dot{\mathcal{A}}_{t}^2 + \frac{1}{2} \dot{\mathcal{A}}_{x}^2 + B_{0} \mathcal{H}_{tz} \, \dot{\mathcal{A}}_{y} + \frac{1}{2} \dot{\mathcal{A}}_{y}^2 -B_{0} \mathcal{H}_{ty} \, \dot{\mathcal{A}}_{z} + \frac{1}{2} \dot{\mathcal{A}}_{z}^2 + \frac{1}{2 \, \kappa^2} \dot{\mathcal{H}}_{xy}^2 + \frac{1}{2 \, \kappa^2} \dot{\mathcal{H}}_{xz}^2 -\frac{1}{2 \, \kappa^2} \dot{\mathcal{H}}_{xx} \, \dot{\mathcal{H}}_{yy} + \frac{1}{2 \, \kappa^2} \dot{\mathcal{H}}_{yz}^2 -\frac{1}{2 \, \kappa^2} \dot{\mathcal{H}}_{xx} \, \dot{\mathcal{H}}_{zz} -\frac{1}{2 \, \kappa^2} \dot{\mathcal{H}}_{yy} \, \dot{\mathcal{H}}_{zz} -\frac{2}{\kappa^2} \mathcal{H}_{xz} \, \partial_t \partial_z \mathcal{H}_{tx} -\frac{2}{\kappa^2} \mathcal{H}_{yz} \, \partial_t \partial_z \mathcal{H}_{ty} + \frac{1}{\kappa^{2}} \mathcal{H}_{tt} \, \partial_t \partial_z \mathcal{H}_{tz} + \frac{1}{\kappa^{2}} \mathcal{H}_{xx} \, \partial_t \partial_z \mathcal{H}_{tz} + \frac{1}{\kappa^{2}} \mathcal{H}_{yy} \, \partial_t \partial_z \mathcal{H}_{tz} -\frac{1}{\kappa^{2}} \mathcal{H}_{zz} \, \partial_t \partial_z \mathcal{H}_{tz} + \frac{2}{\kappa^2} \mathcal{H}_{tz} \, \partial_t \partial_z \mathcal{H}_{xx} -\frac{2}{\kappa^2} \mathcal{H}_{tx} \, \partial_t \partial_z \mathcal{H}_{xz} + \frac{2}{\kappa^2} \mathcal{H}_{tz} \, \partial_t \partial_z \mathcal{H}_{yy} -\frac{2}{\kappa^2} \mathcal{H}_{ty} \, \partial_t \partial_z \mathcal{H}_{yz}

\end{gathered}
$}
\medskip


% --- E.7  Parity-odd extension --------------------------------------
\subsection{Parity-odd extension}%
\label{SymbolicParityOdd}
\Cref{ParityOdd20D} introduces a twenty-parameter parity-odd non-minimal sector on top of the
Yang--Mills--PGT base.

Under the self-sector restriction of~\cref{DiracBergmann}, the
self-sector Hamiltonian for this theory is identical in form to
the YM--PGT Hamiltonian of~\cref{SymbolicGeneralNonminimal}. This
is expected: the parity-odd invariants contribute only to
cross-sector interaction terms and to an additional constraint equation.
The self-energy diagonal sector itself is unchanged.
The symbolic output comprises \ParityOddEomCount{} equations of
motion (\ParityOddConstraintCount{} constraints) with
\ParityOddTermCount{} RHS terms.

\medskip
\noindent\adjustbox{max width=\linewidth}{$\displaystyle
\setlength{\jot}{1pt}%
\begin{gathered}
  \begin{aligned}
    \mathscr{L} &= \tfrac{1}{\kappa^2} \tilde{\mathcal{R}} + \beta_{1} \tensor{T}{_a_b_c} \tensor{g}{^a^d} \tensor{g}{^b^e} \tensor{g}{^c^f} \tensor{T}{_d_e_f} + \beta_{2} \tensor{T}{_a_b_c} \tensor{g}{^b^d} \tensor{g}{^a^e} \tensor{g}{^c^f} \tensor{T}{_d_e_f} + \beta_{3} \tensor{T}{^b_b_a} \tensor{g}{^a^d} \tensor{T}{^c_c_d} \\
  &\quad - \tfrac{1}{4} \xi \tensor{\tilde{F}}{_a_b} \tensor{g}{^a^c} \tensor{g}{^b^d} \tensor{\tilde{F}}{_c_d} + \delta_{1} \tensor{\tilde{R}}{_a_b} \tensor{g}{^a^c} \tensor{g}{^b^d} \tensor{\mathcal{F}}{_c_d} + \chi \tensor{\tilde{R}}{_a_b} \tensor{g}{^a^c} \tensor{g}{^b^d} \tensor{\tilde{F}}{_c_d} \\
  &\quad + \zeta_{1} \nabla_{a} \tensor{T}{^b_b_c} \tensor{g}{^a^d} \tensor{g}{^c^e} \tensor{\mathcal{F}}{_d_e} + \zeta_{2} \nabla_{a} \tensor{T}{^a_b_c} \tensor{g}{^b^d} \tensor{g}{^c^e} \tensor{\mathcal{F}}{_d_e} + \zeta_{3} \nabla_{a} \tensor{T}{_b^a_c} \tensor{g}{^b^d} \tensor{g}{^c^e} \tensor{\mathcal{F}}{_d_e} + d_{14} \tensor{\epsilon}{^a^b^c^d} \tensor{T}{^e_a_b} \tensor{T}{_e_c_d} \\
  &\quad + d_{15} \tensor{\epsilon}{^a^b^c^d} \tensor{T}{^e_a_c} \tensor{T}{_e_b_d} + d_{16} \tensor{\epsilon}{^a^b^c^d} \tensor{g}{^e^f} \tensor{T}{_a_e_f} \tensor{T}{_b_c_d} + d_{17} \tensor{\epsilon}{^a^b^c^d} \tensor{T}{^e_a_e} \tensor{T}{_b_c_d} + d_{19} \tensor{\epsilon}{^a^b^c^d} \tensor{\tilde{\mathcal{R}}}{_a_b_c^e} \tensor{\mathcal{F}}{_d_e} \\
  &\quad + d_{20} \tensor{\epsilon}{^a^b^c^d} \tensor{\tilde{\mathcal{R}}}{_a_c_b^e} \tensor{\mathcal{F}}{_d_e} + d_{21} \tensor{\epsilon}{^a^b^c^d} \tensor{\tilde{R}}{_a_b} \tensor{\mathcal{F}}{_c_d} + \zeta_{1} \nabla_{a} \tensor{T}{_b_c_d} \tensor{\epsilon}{^b^c^d^e} \tensor{g}{^a^f} \tensor{\mathcal{F}}{_f_e} \\
  &\quad + \zeta_{2} \nabla_{a} \tensor{T}{_b_c_d} \tensor{\epsilon}{^a^c^d^e} \tensor{g}{^b^f} \tensor{\mathcal{F}}{_f_e} + \zeta_{3} \nabla_{a} \tensor{T}{_b_c_d} \tensor{\epsilon}{^a^b^d^e} \tensor{g}{^c^f} \tensor{\mathcal{F}}{_f_e} + \zeta_{4} \nabla_{a} \tensor{T}{^b_b_c} \tensor{\epsilon}{^a^c^d^e} \tensor{\mathcal{F}}{_d_e} \\
  &\quad + \zeta_{5} \nabla_{a} \tensor{T}{^a_b_c} \tensor{\epsilon}{^b^c^d^e} \tensor{\mathcal{F}}{_d_e} + \zeta_{6} \nabla_{a} \tensor{T}{_b^a_c} \tensor{\epsilon}{^b^c^d^e} \tensor{\mathcal{F}}{_d_e} - \tfrac{1}{4} \tensor{\mathcal{F}}{_a_b} \tensor{g}{^a^c} \tensor{g}{^b^d} \tensor{\mathcal{F}}{_c_d}
  \end{aligned} \\[1ex]
  % \begin{aligned}
  %   \partial_t^{2} \mathcal{A}_{t} &= -B_{0} \partial_z \mathcal{H}_{ty} - \partial_z^2 \mathcal{A}_{t} - d_{19} \partial_z^2 \mathcal{T}_{xyz} + d_{20} \partial_z^2 \mathcal{T}_{xyz} - 2 \, d_{21} \partial_z^2 \mathcal{T}_{xyz} - 2 \, {\zeta}_{1} \partial_z^2 \mathcal{T}_{xyz} + {\zeta}_{3} \partial_z^2 \mathcal{T}_{xyz} + 2 \, {\zeta}_{6} \partial_z^2 \mathcal{T}_{xyz} \\
  % &\quad - \tfrac{3 \, \delta_{1}}{2} \partial_z^2 \mathcal{T}_{yty} - \zeta_{1} \partial_z^2 \mathcal{T}_{yty} + d_{19} \partial_z^2 \mathcal{T}_{yxz} - d_{20} \partial_z^2 \mathcal{T}_{yxz} + 2 \, d_{21} \partial_z^2 \mathcal{T}_{yxz} + 2 \, {\zeta}_{1} \partial_z^2 \mathcal{T}_{yxz} - {\zeta}_{3} \partial_z^2 \mathcal{T}_{yxz} - 2 \, {\zeta}_{6} \partial_z^2 \mathcal{T}_{yxz} \\
  % &\quad + \tfrac{\delta_{1}}{2} \partial_z^2 \mathcal{T}_{ztz} - \zeta_{1} \partial_z^2 \mathcal{T}_{ztz} - 2 \, \zeta_{2} \partial_z^2 \mathcal{T}_{ztz} - \zeta_{3} \partial_z^2 \mathcal{T}_{ztz} + 2 \, d_{21} \partial_z^2 \mathcal{T}_{zxy} - 2 \, {\zeta}_{1} \partial_z^2 \mathcal{T}_{zxy} - 2 \, {\zeta}_{2} \partial_z^2 \mathcal{T}_{zxy} - 4 \, {\zeta}_{5} \partial_z^2 \mathcal{T}_{zxy} \\
  % &\quad - \tfrac{3 \, \delta_{1}}{2} \partial_z^2 \mathcal{T}_{xtx} - \zeta_{1} \partial_z^2 \mathcal{T}_{xtx} - \tfrac{3 \, \delta_{1}}{2} \partial_z \dot{\mathcal{T}}_{xxz} - \zeta_{1} \partial_z \dot{\mathcal{T}}_{xxz} + d_{19} \partial_z \dot{\mathcal{T}}_{ytx} - d_{20} \partial_z \dot{\mathcal{T}}_{ytx} + 2 \, d_{21} \partial_z \dot{\mathcal{T}}_{ytx} \\
  % &\quad + 2 \, {\zeta}_{1} \partial_z \dot{\mathcal{T}}_{ytx} - {\zeta}_{3} \partial_z \dot{\mathcal{T}}_{ytx} - 2 \, {\zeta}_{6} \partial_z \dot{\mathcal{T}}_{ytx} - \tfrac{3 \, \delta_{1}}{2} \partial_z \dot{\mathcal{T}}_{yyz} - \zeta_{1} \partial_z \dot{\mathcal{T}}_{yyz} + \tfrac{\delta_{1}}{2} \partial_z \dot{\mathcal{T}}_{ttz} - \zeta_{1} \partial_z \dot{\mathcal{T}}_{ttz} \\
  % &\quad - 2 \, \zeta_{2} \partial_z \dot{\mathcal{T}}_{ttz} - \zeta_{3} \partial_z \dot{\mathcal{T}}_{ttz} + 2 \, d_{21} \partial_z \dot{\mathcal{T}}_{txy} - 2 \, {\zeta}_{1} \partial_z \dot{\mathcal{T}}_{txy} - 2 \, {\zeta}_{2} \partial_z \dot{\mathcal{T}}_{txy} - 4 \, {\zeta}_{5} \partial_z \dot{\mathcal{T}}_{txy} - d_{19} \partial_z \dot{\mathcal{T}}_{xty} \\
  % &\quad + d_{20} \partial_z \dot{\mathcal{T}}_{xty} - 2 \, d_{21} \partial_z \dot{\mathcal{T}}_{xty} - 2 \, {\zeta}_{1} \partial_z \dot{\mathcal{T}}_{xty} + {\zeta}_{3} \partial_z \dot{\mathcal{T}}_{xty} + 2 \, {\zeta}_{6} \partial_z \dot{\mathcal{T}}_{xty}
  % \end{aligned} \\[1ex]
  \begin{aligned}
    \partial_t^{2} \mathcal{A}_{x} &= B_{0} \partial_z \mathcal{H}_{xy} + \partial_z^2 \mathcal{A}_{x} - \tfrac{3 \, \delta_{1}}{2} \partial_z^2 \mathcal{T}_{ttx} - \zeta_{1} \partial_z^2 \mathcal{T}_{ttx} - d_{19} \partial_z^2 \mathcal{T}_{ytz} + d_{20} \partial_z^2 \mathcal{T}_{ytz} - 2 \, d_{21} \partial_z^2 \mathcal{T}_{ytz} - 2 \, {\zeta}_{1} \partial_z^2 \mathcal{T}_{ytz} \\
  &\quad + {\zeta}_{3} \partial_z^2 \mathcal{T}_{ytz} + 2 \, {\zeta}_{6} \partial_z^2 \mathcal{T}_{ytz} + \tfrac{3 \, \delta_{1}}{2} \partial_z^2 \mathcal{T}_{yxy} + \zeta_{1} \partial_z^2 \mathcal{T}_{yxy} - 2 \, d_{21} \partial_z^2 \mathcal{T}_{zty} + 2 \, {\zeta}_{1} \partial_z^2 \mathcal{T}_{zty} + 2 \, {\zeta}_{2} \partial_z^2 \mathcal{T}_{zty} \\
  &\quad + 4 \, {\zeta}_{5} \partial_z^2 \mathcal{T}_{zty} - \tfrac{1}{2} \delta_{1} \partial_z^2 \mathcal{T}_{zxz} + \zeta_{1} \partial_z^2 \mathcal{T}_{zxz} + 2 \, \zeta_{2} \partial_z^2 \mathcal{T}_{zxz} + \zeta_{3} \partial_z^2 \mathcal{T}_{zxz} - d_{19} \partial_z^2 \mathcal{T}_{tyz} + d_{20} \partial_z^2 \mathcal{T}_{tyz} - 2 \, d_{21} \partial_z^2 \mathcal{T}_{tyz} \\
  &\quad - 2 \, {\zeta}_{1} \partial_z^2 \mathcal{T}_{tyz} + {\zeta}_{3} \partial_z^2 \mathcal{T}_{tyz} + 2 \, {\zeta}_{6} \partial_z^2 \mathcal{T}_{tyz} - d_{19} \partial_z \dot{\mathcal{T}}_{tty} + d_{20} \partial_z \dot{\mathcal{T}}_{tty} - 4 \, d_{21} \partial_z \dot{\mathcal{T}}_{tty} + 2 \, {\zeta}_{2} \partial_z \dot{\mathcal{T}}_{tty} + {\zeta}_{3} \partial_z \dot{\mathcal{T}}_{tty} \\
  &\quad + 4 \, {\zeta}_{5} \partial_z \dot{\mathcal{T}}_{tty} + 2 \, {\zeta}_{6} \partial_z \dot{\mathcal{T}}_{tty} - 2 \, \delta_{1} \partial_z \dot{\mathcal{T}}_{ztx} + 2 \, \zeta_{2} \partial_z \dot{\mathcal{T}}_{ztx} + \zeta_{3} \partial_z \dot{\mathcal{T}}_{ztx} - d_{19} \partial_z \dot{\mathcal{T}}_{zyz} + d_{20} \partial_z \dot{\mathcal{T}}_{zyz} \\
  &\quad - 4 \, d_{21} \partial_z \dot{\mathcal{T}}_{zyz} + 2 \, {\zeta}_{2} \partial_z \dot{\mathcal{T}}_{zyz} + {\zeta}_{3} \partial_z \dot{\mathcal{T}}_{zyz} + 4 \, {\zeta}_{5} \partial_z \dot{\mathcal{T}}_{zyz} + 2 \, {\zeta}_{6} \partial_z \dot{\mathcal{T}}_{zyz} - 2 \, \delta_{1} \partial_z \dot{\mathcal{T}}_{txz} + 2 \, \zeta_{2} \partial_z \dot{\mathcal{T}}_{txz} \\
  &\quad + \zeta_{3} \partial_z \dot{\mathcal{T}}_{txz} - \tfrac{1}{2} \delta_{1} \partial_t^2 \mathcal{T}_{ttx} + \zeta_{1} \partial_t^2 \mathcal{T}_{ttx} + 2 \, \zeta_{2} \partial_t^2 \mathcal{T}_{ttx} + \zeta_{3} \partial_t^2 \mathcal{T}_{ttx} + d_{19} \partial_t^2 \mathcal{T}_{ytz} - d_{20} \partial_t^2 \mathcal{T}_{ytz} + 2 \, d_{21} \partial_t^2 \mathcal{T}_{ytz} \\
  &\quad + 2 \, {\zeta}_{1} \partial_t^2 \mathcal{T}_{ytz} - {\zeta}_{3} \partial_t^2 \mathcal{T}_{ytz} - 2 \, {\zeta}_{6} \partial_t^2 \mathcal{T}_{ytz} - \tfrac{3 \, \delta_{1}}{2} \partial_t^2 \mathcal{T}_{yxy} - \zeta_{1} \partial_t^2 \mathcal{T}_{yxy} - d_{19} \partial_t^2 \mathcal{T}_{zty} + d_{20} \partial_t^2 \mathcal{T}_{zty} - 2 \, d_{21} \partial_t^2 \mathcal{T}_{zty} \\
  &\quad - 2 \, {\zeta}_{1} \partial_t^2 \mathcal{T}_{zty} + {\zeta}_{3} \partial_t^2 \mathcal{T}_{zty} + 2 \, {\zeta}_{6} \partial_t^2 \mathcal{T}_{zty} - \tfrac{3 \, \delta_{1}}{2} \partial_t^2 \mathcal{T}_{zxz} - \zeta_{1} \partial_t^2 \mathcal{T}_{zxz} - 2 \, d_{21} \partial_t^2 \mathcal{T}_{tyz} + 2 \, {\zeta}_{1} \partial_t^2 \mathcal{T}_{tyz} \\
  &\quad + 2 \, {\zeta}_{2} \partial_t^2 \mathcal{T}_{tyz} + 4 \, {\zeta}_{5} \partial_t^2 \mathcal{T}_{tyz}
  \end{aligned} \\[1ex]
  % \begin{aligned}
  %   \partial_t^{2} \mathcal{A}_{y} &= \tfrac{B_{0}}{2} \partial_z \mathcal{H}_{tt} - \tfrac{1}{2} B_{0} \partial_z \mathcal{H}_{xx} + \tfrac{B_{0}}{2} \partial_z \mathcal{H}_{yy} + \tfrac{B_{0}}{2} \partial_z \mathcal{H}_{zz} + \partial_z^2 \mathcal{A}_{y} - \tfrac{3 \, \delta_{1}}{2} \partial_z^2 \mathcal{T}_{tty} - \zeta_{1} \partial_z^2 \mathcal{T}_{tty} \\
  % &\quad + 2 \, d_{21} \partial_z^2 \mathcal{T}_{ztx} - 2 \, {\zeta}_{1} \partial_z^2 \mathcal{T}_{ztx} - 2 \, {\zeta}_{2} \partial_z^2 \mathcal{T}_{ztx} - 4 \, {\zeta}_{5} \partial_z^2 \mathcal{T}_{ztx} - \tfrac{1}{2} \delta_{1} \partial_z^2 \mathcal{T}_{zyz} + \zeta_{1} \partial_z^2 \mathcal{T}_{zyz} + 2 \, \zeta_{2} \partial_z^2 \mathcal{T}_{zyz} + \zeta_{3} \partial_z^2 \mathcal{T}_{zyz} \\
  % &\quad + d_{19} \partial_z^2 \mathcal{T}_{txz} - d_{20} \partial_z^2 \mathcal{T}_{txz} + 2 \, d_{21} \partial_z^2 \mathcal{T}_{txz} + 2 \, {\zeta}_{1} \partial_z^2 \mathcal{T}_{txz} - {\zeta}_{3} \partial_z^2 \mathcal{T}_{txz} - 2 \, {\zeta}_{6} \partial_z^2 \mathcal{T}_{txz} + d_{19} \partial_z^2 \mathcal{T}_{xtz} - d_{20} \partial_z^2 \mathcal{T}_{xtz} \\
  % &\quad + 2 \, d_{21} \partial_z^2 \mathcal{T}_{xtz} + 2 \, {\zeta}_{1} \partial_z^2 \mathcal{T}_{xtz} - {\zeta}_{3} \partial_z^2 \mathcal{T}_{xtz} - 2 \, {\zeta}_{6} \partial_z^2 \mathcal{T}_{xtz} - \tfrac{3 \, \delta_{1}}{2} \partial_z^2 \mathcal{T}_{xxy} - \zeta_{1} \partial_z^2 \mathcal{T}_{xxy} - B_{0} \partial_t \mathcal{H}_{tz} + d_{19} \partial_z \dot{\mathcal{T}}_{ttx} \\
  % &\quad - d_{20} \partial_z \dot{\mathcal{T}}_{ttx} + 4 \, d_{21} \partial_z \dot{\mathcal{T}}_{ttx} - 2 \, {\zeta}_{2} \partial_z \dot{\mathcal{T}}_{ttx} - {\zeta}_{3} \partial_z \dot{\mathcal{T}}_{ttx} - 4 \, {\zeta}_{5} \partial_z \dot{\mathcal{T}}_{ttx} - 2 \, {\zeta}_{6} \partial_z \dot{\mathcal{T}}_{ttx} - 2 \, \delta_{1} \partial_z \dot{\mathcal{T}}_{zty} \\
  % &\quad + 2 \, \zeta_{2} \partial_z \dot{\mathcal{T}}_{zty} + \zeta_{3} \partial_z \dot{\mathcal{T}}_{zty} + d_{19} \partial_z \dot{\mathcal{T}}_{zxz} - d_{20} \partial_z \dot{\mathcal{T}}_{zxz} + 4 \, d_{21} \partial_z \dot{\mathcal{T}}_{zxz} - 2 \, {\zeta}_{2} \partial_z \dot{\mathcal{T}}_{zxz} - {\zeta}_{3} \partial_z \dot{\mathcal{T}}_{zxz} \\
  % &\quad - 4 \, {\zeta}_{5} \partial_z \dot{\mathcal{T}}_{zxz} - 2 \, {\zeta}_{6} \partial_z \dot{\mathcal{T}}_{zxz} - 2 \, \delta_{1} \partial_z \dot{\mathcal{T}}_{tyz} + 2 \, \zeta_{2} \partial_z \dot{\mathcal{T}}_{tyz} + \zeta_{3} \partial_z \dot{\mathcal{T}}_{tyz} - \tfrac{1}{2} \delta_{1} \partial_t^2 \mathcal{T}_{tty} + \zeta_{1} \partial_t^2 \mathcal{T}_{tty} \\
  % &\quad + 2 \, \zeta_{2} \partial_t^2 \mathcal{T}_{tty} + \zeta_{3} \partial_t^2 \mathcal{T}_{tty} + d_{19} \partial_t^2 \mathcal{T}_{ztx} - d_{20} \partial_t^2 \mathcal{T}_{ztx} + 2 \, d_{21} \partial_t^2 \mathcal{T}_{ztx} + 2 \, {\zeta}_{1} \partial_t^2 \mathcal{T}_{ztx} - {\zeta}_{3} \partial_t^2 \mathcal{T}_{ztx} - 2 \, {\zeta}_{6} \partial_t^2 \mathcal{T}_{ztx} \\
  % &\quad - \tfrac{3 \, \delta_{1}}{2} \partial_t^2 \mathcal{T}_{zyz} - \zeta_{1} \partial_t^2 \mathcal{T}_{zyz} + 2 \, d_{21} \partial_t^2 \mathcal{T}_{txz} - 2 \, {\zeta}_{1} \partial_t^2 \mathcal{T}_{txz} - 2 \, {\zeta}_{2} \partial_t^2 \mathcal{T}_{txz} - 4 \, {\zeta}_{5} \partial_t^2 \mathcal{T}_{txz} - d_{19} \partial_t^2 \mathcal{T}_{xtz} + d_{20} \partial_t^2 \mathcal{T}_{xtz} \\
  % &\quad - 2 \, d_{21} \partial_t^2 \mathcal{T}_{xtz} - 2 \, {\zeta}_{1} \partial_t^2 \mathcal{T}_{xtz} + {\zeta}_{3} \partial_t^2 \mathcal{T}_{xtz} + 2 \, {\zeta}_{6} \partial_t^2 \mathcal{T}_{xtz} + \tfrac{3 \, \delta_{1}}{2} \partial_t^2 \mathcal{T}_{xxy} + \zeta_{1} \partial_t^2 \mathcal{T}_{xxy}
  % \end{aligned} \\[1ex]
  % \begin{aligned}
  %   \partial_t^{2} \mathcal{A}_{z} &= \partial_z^2 \mathcal{A}_{z} + B_{0} \partial_t \mathcal{H}_{ty} + d_{19} \partial_z \dot{\mathcal{T}}_{xyz} - d_{20} \partial_z \dot{\mathcal{T}}_{xyz} + 2 \, d_{21} \partial_z \dot{\mathcal{T}}_{xyz} + 2 \, {\zeta}_{1} \partial_z \dot{\mathcal{T}}_{xyz} - {\zeta}_{3} \partial_z \dot{\mathcal{T}}_{xyz} - 2 \, {\zeta}_{6} \partial_z \dot{\mathcal{T}}_{xyz} \\
  % &\quad + \tfrac{3 \, \delta_{1}}{2} \partial_z \dot{\mathcal{T}}_{yty} + \zeta_{1} \partial_z \dot{\mathcal{T}}_{yty} - d_{19} \partial_z \dot{\mathcal{T}}_{yxz} + d_{20} \partial_z \dot{\mathcal{T}}_{yxz} - 2 \, d_{21} \partial_z \dot{\mathcal{T}}_{yxz} - 2 \, {\zeta}_{1} \partial_z \dot{\mathcal{T}}_{yxz} + {\zeta}_{3} \partial_z \dot{\mathcal{T}}_{yxz} \\
  % &\quad + 2 \, {\zeta}_{6} \partial_z \dot{\mathcal{T}}_{yxz} - \tfrac{1}{2} \delta_{1} \partial_z \dot{\mathcal{T}}_{ztz} + \zeta_{1} \partial_z \dot{\mathcal{T}}_{ztz} + 2 \, \zeta_{2} \partial_z \dot{\mathcal{T}}_{ztz} + \zeta_{3} \partial_z \dot{\mathcal{T}}_{ztz} - 2 \, d_{21} \partial_z \dot{\mathcal{T}}_{zxy} + 2 \, {\zeta}_{1} \partial_z \dot{\mathcal{T}}_{zxy} \\
  % &\quad + 2 \, {\zeta}_{2} \partial_z \dot{\mathcal{T}}_{zxy} + 4 \, {\zeta}_{5} \partial_z \dot{\mathcal{T}}_{zxy} + \tfrac{3 \, \delta_{1}}{2} \partial_z \dot{\mathcal{T}}_{xtx} + \zeta_{1} \partial_z \dot{\mathcal{T}}_{xtx} + \tfrac{3 \, \delta_{1}}{2} \partial_t^2 \mathcal{T}_{xxz} + \zeta_{1} \partial_t^2 \mathcal{T}_{xxz} - d_{19} \partial_t^2 \mathcal{T}_{ytx} \\
  % &\quad + d_{20} \partial_t^2 \mathcal{T}_{ytx} - 2 \, d_{21} \partial_t^2 \mathcal{T}_{ytx} - 2 \, {\zeta}_{1} \partial_t^2 \mathcal{T}_{ytx} + {\zeta}_{3} \partial_t^2 \mathcal{T}_{ytx} + 2 \, {\zeta}_{6} \partial_t^2 \mathcal{T}_{ytx} + \tfrac{3 \, \delta_{1}}{2} \partial_t^2 \mathcal{T}_{yyz} + \zeta_{1} \partial_t^2 \mathcal{T}_{yyz} - \tfrac{1}{2} \delta_{1} \partial_t^2 \mathcal{T}_{ttz} \\
  % &\quad + \zeta_{1} \partial_t^2 \mathcal{T}_{ttz} + 2 \, \zeta_{2} \partial_t^2 \mathcal{T}_{ttz} + \zeta_{3} \partial_t^2 \mathcal{T}_{ttz} - 2 \, d_{21} \partial_t^2 \mathcal{T}_{txy} + 2 \, {\zeta}_{1} \partial_t^2 \mathcal{T}_{txy} + 2 \, {\zeta}_{2} \partial_t^2 \mathcal{T}_{txy} + 4 \, {\zeta}_{5} \partial_t^2 \mathcal{T}_{txy} + d_{19} \partial_t^2 \mathcal{T}_{xty} \\
  % &\quad - d_{20} \partial_t^2 \mathcal{T}_{xty} + 2 \, d_{21} \partial_t^2 \mathcal{T}_{xty} + 2 \, {\zeta}_{1} \partial_t^2 \mathcal{T}_{xty} - {\zeta}_{3} \partial_t^2 \mathcal{T}_{xty} - 2 \, {\zeta}_{6} \partial_t^2 \mathcal{T}_{xty}
  % \end{aligned} \\[1ex]
  % \begin{aligned}
  %   0 &= \tfrac{B_{0}^2}{8} \mathcal{H}_{tt} + \tfrac{B_{0}^2}{8} \mathcal{H}_{xx} - \tfrac{1}{8} B_{0}^2 \mathcal{H}_{yy} - \tfrac{1}{8} B_{0}^2 \mathcal{H}_{zz} - \tfrac{1}{2} B_{0} \partial_z \mathcal{A}_{y} - \tfrac{1}{4} B_{0} \, \delta_{1} \partial_z \mathcal{T}_{tty} + \tfrac{B_{0} \, \zeta_{1}}{2} \partial_z \mathcal{T}_{tty} + B_{0} \, \zeta_{2} \partial_z \mathcal{T}_{tty} \\
  % &\quad + \tfrac{B_{0} \, \zeta_{3}}{2} \partial_z \mathcal{T}_{tty} + \tfrac{B_{0} \, d_{19}}{2} \partial_z \mathcal{T}_{ztx} - \tfrac{1}{2} B_{0} \, d_{20} \partial_z \mathcal{T}_{ztx} + B_{0} \, d_{21} \partial_z \mathcal{T}_{ztx} + B_{0} \, {\zeta}_{1} \partial_z \mathcal{T}_{ztx} - \tfrac{1}{2} B_{0} \, {\zeta}_{3} \partial_z \mathcal{T}_{ztx} \\
  % &\quad - B_{0} \, {\zeta}_{6} \partial_z \mathcal{T}_{ztx} - \tfrac{3 \, B_{0} \, \delta_{1}}{4} \partial_z \mathcal{T}_{zyz} - \tfrac{1}{2} B_{0} \, \zeta_{1} \partial_z \mathcal{T}_{zyz} + B_{0} \, d_{21} \partial_z \mathcal{T}_{txz} - B_{0} \, {\zeta}_{1} \partial_z \mathcal{T}_{txz} - B_{0} \, {\zeta}_{2} \partial_z \mathcal{T}_{txz} - 2 \, B_{0} \, {\zeta}_{5} \partial_z \mathcal{T}_{txz} \\
  % &\quad - \tfrac{1}{2} B_{0} \, d_{19} \partial_z \mathcal{T}_{xtz} + \tfrac{B_{0} \, d_{20}}{2} \partial_z \mathcal{T}_{xtz} - B_{0} \, d_{21} \partial_z \mathcal{T}_{xtz} - B_{0} \, {\zeta}_{1} \partial_z \mathcal{T}_{xtz} + \tfrac{B_{0} \, {\zeta}_{3}}{2} \partial_z \mathcal{T}_{xtz} + B_{0} \, {\zeta}_{6} \partial_z \mathcal{T}_{xtz} \\
  % &\quad + \tfrac{3 \, B_{0} \, \delta_{1}}{4} \partial_z \mathcal{T}_{xxy} + \tfrac{B_{0} \, \zeta_{1}}{2} \partial_z \mathcal{T}_{xxy} + \tfrac{1}{2 \, \kappa^2} \partial_z^2 \mathcal{H}_{xx} + \tfrac{1}{2 \, \kappa^2} \partial_z^2 \mathcal{H}_{yy}
  % \end{aligned} \\[1ex]
  % \begin{aligned}
  %   0 &= -\tfrac{1}{2} B_{0}^2 \mathcal{H}_{tx} + \tfrac{B_{0} \, d_{19}}{2} \partial_z \mathcal{T}_{xxz} - \tfrac{1}{2} B_{0} \, d_{20} \partial_z \mathcal{T}_{xxz} + 2 \, B_{0} \, d_{21} \partial_z \mathcal{T}_{xxz} - B_{0} \, {\zeta}_{2} \partial_z \mathcal{T}_{xxz} - \tfrac{1}{2} B_{0} \, {\zeta}_{3} \partial_z \mathcal{T}_{xxz} - 2 \, B_{0} \, {\zeta}_{5} \partial_z \mathcal{T}_{xxz} \\
  % &\quad - B_{0} \, {\zeta}_{6} \partial_z \mathcal{T}_{xxz} - \tfrac{1}{2} B_{0} \, d_{19} \partial_z \mathcal{T}_{ttz} + \tfrac{B_{0} \, d_{20}}{2} \partial_z \mathcal{T}_{ttz} - 2 \, B_{0} \, d_{21} \partial_z \mathcal{T}_{ttz} + B_{0} \, {\zeta}_{2} \partial_z \mathcal{T}_{ttz} + \tfrac{B_{0} \, {\zeta}_{3}}{2} \partial_z \mathcal{T}_{ttz} \\
  % &\quad + 2 \, B_{0} \, {\zeta}_{5} \partial_z \mathcal{T}_{ttz} + B_{0} \, {\zeta}_{6} \partial_z \mathcal{T}_{ttz} + B_{0} \, \delta_{1} \partial_z \mathcal{T}_{txy} - B_{0} \, \zeta_{2} \partial_z \mathcal{T}_{txy} - \tfrac{1}{2} B_{0} \, \zeta_{3} \partial_z \mathcal{T}_{txy} - B_{0} \, \delta_{1} \partial_z \mathcal{T}_{xty} + B_{0} \, \zeta_{2} \partial_z \mathcal{T}_{xty} \\
  % &\quad + \tfrac{B_{0} \, \zeta_{3}}{2} \partial_z \mathcal{T}_{xty} - \tfrac{1}{\kappa^{2}} \partial_z^2 \mathcal{H}_{tx} + \tfrac{1}{\kappa^{2}} \partial_z \dot{\mathcal{H}}_{xz}
  % \end{aligned} \\[1ex]
  \begin{aligned}
    0 &= \tfrac{B_{0}^2}{2} \mathcal{H}_{ty} + B_{0} \partial_z \mathcal{A}_{t} + B_{0} \, d_{19} \partial_z \mathcal{T}_{xyz} - B_{0} \, d_{20} \partial_z \mathcal{T}_{xyz} + 2 \, B_{0} \, d_{21} \partial_z \mathcal{T}_{xyz} + 2 \, B_{0} \, {\zeta}_{1} \partial_z \mathcal{T}_{xyz} - B_{0} \, {\zeta}_{3} \partial_z \mathcal{T}_{xyz} - 2 \, B_{0} \, {\zeta}_{6} \partial_z \mathcal{T}_{xyz} \\
  &\quad + \tfrac{B_{0} \, \delta_{1}}{2} \partial_z \mathcal{T}_{yty} + B_{0} \, \zeta_{1} \partial_z \mathcal{T}_{yty} + B_{0} \, \zeta_{2} \partial_z \mathcal{T}_{yty} + \tfrac{B_{0} \, \zeta_{3}}{2} \partial_z \mathcal{T}_{yty} - \tfrac{1}{2} B_{0} \, d_{19} \partial_z \mathcal{T}_{yxz} + \tfrac{B_{0} \, d_{20}}{2} \partial_z \mathcal{T}_{yxz} \\
  &\quad - 2 \, B_{0} \, {\zeta}_{1} \partial_z \mathcal{T}_{yxz} - B_{0} \, {\zeta}_{2} \partial_z \mathcal{T}_{yxz} + \tfrac{B_{0} \, {\zeta}_{3}}{2} \partial_z \mathcal{T}_{yxz} - 2 \, B_{0} \, {\zeta}_{5} \partial_z \mathcal{T}_{yxz} + B_{0} \, {\zeta}_{6} \partial_z \mathcal{T}_{yxz} + \tfrac{B_{0} \, \delta_{1}}{2} \partial_z \mathcal{T}_{ztz} \\
  &\quad + B_{0} \, \zeta_{1} \partial_z \mathcal{T}_{ztz} + B_{0} \, \zeta_{2} \partial_z \mathcal{T}_{ztz} + \tfrac{B_{0} \, \zeta_{3}}{2} \partial_z \mathcal{T}_{ztz} + \tfrac{B_{0} \, d_{19}}{2} \partial_z \mathcal{T}_{zxy} - \tfrac{1}{2} B_{0} \, d_{20} \partial_z \mathcal{T}_{zxy} + 2 \, B_{0} \, {\zeta}_{1} \partial_z \mathcal{T}_{zxy} + B_{0} \, {\zeta}_{2} \partial_z \mathcal{T}_{zxy} \\
  &\quad - \tfrac{1}{2} B_{0} \, {\zeta}_{3} \partial_z \mathcal{T}_{zxy} + 2 \, B_{0} \, {\zeta}_{5} \partial_z \mathcal{T}_{zxy} - B_{0} \, {\zeta}_{6} \partial_z \mathcal{T}_{zxy} + \tfrac{3 \, B_{0} \, \delta_{1}}{2} \partial_z \mathcal{T}_{xtx} + B_{0} \, \zeta_{1} \partial_z \mathcal{T}_{xtx} - \tfrac{1}{\kappa^{2}} \partial_z^2 \mathcal{H}_{ty} - B_{0} \partial_t \mathcal{A}_{z} \\
  &\quad + \tfrac{3 \, B_{0} \, \delta_{1}}{2} \partial_t \mathcal{T}_{xxz} + B_{0} \, \zeta_{1} \partial_t \mathcal{T}_{xxz} - B_{0} \, d_{19} \partial_t \mathcal{T}_{ytx} + B_{0} \, d_{20} \partial_t \mathcal{T}_{ytx} - 2 \, B_{0} \, d_{21} \partial_t \mathcal{T}_{ytx} - 2 \, B_{0} \, {\zeta}_{1} \partial_t \mathcal{T}_{ytx} + B_{0} \, {\zeta}_{3} \partial_t \mathcal{T}_{ytx} \\
  &\quad + 2 \, B_{0} \, {\zeta}_{6} \partial_t \mathcal{T}_{ytx} + \tfrac{3 \, B_{0} \, \delta_{1}}{2} \partial_t \mathcal{T}_{yyz} + B_{0} \, \zeta_{1} \partial_t \mathcal{T}_{yyz} - \tfrac{1}{2} B_{0} \, \delta_{1} \partial_t \mathcal{T}_{ttz} + B_{0} \, \zeta_{1} \partial_t \mathcal{T}_{ttz} + 2 \, B_{0} \, \zeta_{2} \partial_t \mathcal{T}_{ttz} \\
  &\quad + B_{0} \, \zeta_{3} \partial_t \mathcal{T}_{ttz} - 2 \, B_{0} \, d_{21} \partial_t \mathcal{T}_{txy} + 2 \, B_{0} \, {\zeta}_{1} \partial_t \mathcal{T}_{txy} + 2 \, B_{0} \, {\zeta}_{2} \partial_t \mathcal{T}_{txy} + 4 \, B_{0} \, {\zeta}_{5} \partial_t \mathcal{T}_{txy} + B_{0} \, d_{19} \partial_t \mathcal{T}_{xty} - B_{0} \, d_{20} \partial_t \mathcal{T}_{xty} \\
  &\quad + 2 \, B_{0} \, d_{21} \partial_t \mathcal{T}_{xty} + 2 \, B_{0} \, {\zeta}_{1} \partial_t \mathcal{T}_{xty} - B_{0} \, {\zeta}_{3} \partial_t \mathcal{T}_{xty} - 2 \, B_{0} \, {\zeta}_{6} \partial_t \mathcal{T}_{xty} + \tfrac{1}{\kappa^{2}} \partial_z \dot{\mathcal{H}}_{yz}
  \end{aligned} \\[1ex]
  % \begin{aligned}
  %   0 &= \tfrac{B_{0}^2}{2} \mathcal{H}_{tz} + B_{0} \partial_t \mathcal{A}_{y} + \tfrac{B_{0} \, \delta_{1}}{2} \partial_t \mathcal{T}_{tty} - B_{0} \, \zeta_{1} \partial_t \mathcal{T}_{tty} - 2 \, B_{0} \, \zeta_{2} \partial_t \mathcal{T}_{tty} - B_{0} \, \zeta_{3} \partial_t \mathcal{T}_{tty} - B_{0} \, d_{19} \partial_t \mathcal{T}_{ztx} + B_{0} \, d_{20} \partial_t \mathcal{T}_{ztx} \\
  % &\quad - 2 \, B_{0} \, d_{21} \partial_t \mathcal{T}_{ztx} - 2 \, B_{0} \, {\zeta}_{1} \partial_t \mathcal{T}_{ztx} + B_{0} \, {\zeta}_{3} \partial_t \mathcal{T}_{ztx} + 2 \, B_{0} \, {\zeta}_{6} \partial_t \mathcal{T}_{ztx} + \tfrac{3 \, B_{0} \, \delta_{1}}{2} \partial_t \mathcal{T}_{zyz} + B_{0} \, \zeta_{1} \partial_t \mathcal{T}_{zyz} - 2 \, B_{0} \, d_{21} \partial_t \mathcal{T}_{txz} \\
  % &\quad + 2 \, B_{0} \, {\zeta}_{1} \partial_t \mathcal{T}_{txz} + 2 \, B_{0} \, {\zeta}_{2} \partial_t \mathcal{T}_{txz} + 4 \, B_{0} \, {\zeta}_{5} \partial_t \mathcal{T}_{txz} + B_{0} \, d_{19} \partial_t \mathcal{T}_{xtz} - B_{0} \, d_{20} \partial_t \mathcal{T}_{xtz} + 2 \, B_{0} \, d_{21} \partial_t \mathcal{T}_{xtz} + 2 \, B_{0} \, {\zeta}_{1} \partial_t \mathcal{T}_{xtz} \\
  % &\quad - B_{0} \, {\zeta}_{3} \partial_t \mathcal{T}_{xtz} - 2 \, B_{0} \, {\zeta}_{6} \partial_t \mathcal{T}_{xtz} - \tfrac{3 \, B_{0} \, \delta_{1}}{2} \partial_t \mathcal{T}_{xxy} - B_{0} \, \zeta_{1} \partial_t \mathcal{T}_{xxy} - \tfrac{1}{\kappa^{2}} \partial_z \dot{\mathcal{H}}_{xx} - \tfrac{1}{\kappa^{2}} \partial_z \dot{\mathcal{H}}_{yy}
  % \end{aligned} \\[1ex]
  % \begin{aligned}
  %   0 &= \tfrac{B_{0}^2}{8} \mathcal{H}_{tt} + \tfrac{B_{0}^2}{8} \mathcal{H}_{xx} + \tfrac{B_{0}^2}{8} \mathcal{H}_{yy} + \tfrac{B_{0}^2}{8} \mathcal{H}_{zz} + \tfrac{B_{0}}{2} \partial_z \mathcal{A}_{y} - \tfrac{3 \, B_{0} \, \delta_{1}}{4} \partial_z \mathcal{T}_{tty} - \tfrac{1}{2} B_{0} \, \zeta_{1} \partial_z \mathcal{T}_{tty} - \tfrac{1}{2} B_{0} \, d_{19} \partial_z \mathcal{T}_{ztx} \\
  % &\quad + \tfrac{B_{0} \, d_{20}}{2} \partial_z \mathcal{T}_{ztx} - B_{0} \, d_{21} \partial_z \mathcal{T}_{ztx} - B_{0} \, {\zeta}_{1} \partial_z \mathcal{T}_{ztx} + \tfrac{B_{0} \, {\zeta}_{3}}{2} \partial_z \mathcal{T}_{ztx} + B_{0} \, {\zeta}_{6} \partial_z \mathcal{T}_{ztx} + \tfrac{3 \, B_{0} \, \delta_{1}}{4} \partial_z \mathcal{T}_{zyz} \\
  % &\quad + \tfrac{B_{0} \, \zeta_{1}}{2} \partial_z \mathcal{T}_{zyz} + \tfrac{B_{0} \, d_{19}}{2} \partial_z \mathcal{T}_{txz} - \tfrac{1}{2} B_{0} \, d_{20} \partial_z \mathcal{T}_{txz} + B_{0} \, d_{21} \partial_z \mathcal{T}_{txz} + B_{0} \, {\zeta}_{1} \partial_z \mathcal{T}_{txz} - \tfrac{1}{2} B_{0} \, {\zeta}_{3} \partial_z \mathcal{T}_{txz} \\
  % &\quad - B_{0} \, {\zeta}_{6} \partial_z \mathcal{T}_{txz} - B_{0} \, d_{21} \partial_z \mathcal{T}_{xtz} + B_{0} \, {\zeta}_{1} \partial_z \mathcal{T}_{xtz} + B_{0} \, {\zeta}_{2} \partial_z \mathcal{T}_{xtz} + 2 \, B_{0} \, {\zeta}_{5} \partial_z \mathcal{T}_{xtz} + \tfrac{B_{0} \, \delta_{1}}{4} \partial_z \mathcal{T}_{xxy} - \tfrac{1}{2} B_{0} \, \zeta_{1} \partial_z \mathcal{T}_{xxy} \\
  % &\quad - B_{0} \, \zeta_{2} \partial_z \mathcal{T}_{xxy} - \tfrac{1}{2} B_{0} \, \zeta_{3} \partial_z \mathcal{T}_{xxy} + \tfrac{1}{2 \, \kappa^2} \partial_z^2 \mathcal{H}_{tt} - \tfrac{1}{2 \, \kappa^2} \partial_z^2 \mathcal{H}_{yy} - \tfrac{1}{\kappa^{2}} \partial_z \dot{\mathcal{H}}_{tz} + \tfrac{1}{2 \, \kappa^2} \partial_t^2 \mathcal{H}_{yy} \\
  % &\quad + \tfrac{1}{2 \, \kappa^2} \partial_t^2 \mathcal{H}_{zz}
  % \end{aligned} \\[1ex]
  \begin{aligned}
    -\tfrac{1}{\kappa^{2}} \, \partial_t^{2} \mathcal{H}_{xy} &= \tfrac{B_{0}^2}{2} \mathcal{H}_{xy} + B_{0} \partial_z \mathcal{A}_{x} - \tfrac{3 \, B_{0} \, \delta_{1}}{2} \partial_z \mathcal{T}_{ttx} - B_{0} \, \zeta_{1} \partial_z \mathcal{T}_{ttx} - \tfrac{1}{2} B_{0} \, d_{19} \partial_z \mathcal{T}_{ytz} + \tfrac{B_{0} \, d_{20}}{2} \partial_z \mathcal{T}_{ytz} \\
  &\quad - 2 \, B_{0} \, {\zeta}_{1} \partial_z \mathcal{T}_{ytz} - B_{0} \, {\zeta}_{2} \partial_z \mathcal{T}_{ytz} + \tfrac{B_{0} \, {\zeta}_{3}}{2} \partial_z \mathcal{T}_{ytz} - 2 \, B_{0} \, {\zeta}_{5} \partial_z \mathcal{T}_{ytz} + B_{0} \, {\zeta}_{6} \partial_z \mathcal{T}_{ytz} + \tfrac{B_{0} \, \delta_{1}}{2} \partial_z \mathcal{T}_{yxy} \\
  &\quad + B_{0} \, \zeta_{1} \partial_z \mathcal{T}_{yxy} + B_{0} \, \zeta_{2} \partial_z \mathcal{T}_{yxy} + \tfrac{B_{0} \, \zeta_{3}}{2} \partial_z \mathcal{T}_{yxy} + \tfrac{B_{0} \, d_{19}}{2} \partial_z \mathcal{T}_{zty} - \tfrac{1}{2} B_{0} \, d_{20} \partial_z \mathcal{T}_{zty} + 2 \, B_{0} \, {\zeta}_{1} \partial_z \mathcal{T}_{zty} + B_{0} \, {\zeta}_{2} \partial_z \mathcal{T}_{zty} \\
  &\quad - \tfrac{1}{2} B_{0} \, {\zeta}_{3} \partial_z \mathcal{T}_{zty} + 2 \, B_{0} \, {\zeta}_{5} \partial_z \mathcal{T}_{zty} - B_{0} \, {\zeta}_{6} \partial_z \mathcal{T}_{zty} + \tfrac{B_{0} \, \delta_{1}}{2} \partial_z \mathcal{T}_{zxz} + B_{0} \, \zeta_{1} \partial_z \mathcal{T}_{zxz} + B_{0} \, \zeta_{2} \partial_z \mathcal{T}_{zxz} \\
  &\quad + \tfrac{B_{0} \, \zeta_{3}}{2} \partial_z \mathcal{T}_{zxz} - B_{0} \, d_{19} \partial_z \mathcal{T}_{tyz} + B_{0} \, d_{20} \partial_z \mathcal{T}_{tyz} - 2 \, B_{0} \, d_{21} \partial_z \mathcal{T}_{tyz} - 2 \, B_{0} \, {\zeta}_{1} \partial_z \mathcal{T}_{tyz} + B_{0} \, {\zeta}_{3} \partial_z \mathcal{T}_{tyz} + 2 \, B_{0} \, {\zeta}_{6} \partial_z \mathcal{T}_{tyz} \\
  &\quad - \tfrac{1}{\kappa^{2}} \partial_z^2 \mathcal{H}_{xy} - \tfrac{1}{2} B_{0} \, d_{19} \partial_t \mathcal{T}_{tty} + \tfrac{B_{0} \, d_{20}}{2} \partial_t \mathcal{T}_{tty} - 2 \, B_{0} \, d_{21} \partial_t \mathcal{T}_{tty} + B_{0} \, {\zeta}_{2} \partial_t \mathcal{T}_{tty} + \tfrac{B_{0} \, {\zeta}_{3}}{2} \partial_t \mathcal{T}_{tty} \\
  &\quad + 2 \, B_{0} \, {\zeta}_{5} \partial_t \mathcal{T}_{tty} + B_{0} \, {\zeta}_{6} \partial_t \mathcal{T}_{tty} - B_{0} \, \delta_{1} \partial_t \mathcal{T}_{txz} + B_{0} \, \zeta_{2} \partial_t \mathcal{T}_{txz} + \tfrac{B_{0} \, \zeta_{3}}{2} \partial_t \mathcal{T}_{txz} + B_{0} \, \delta_{1} \partial_t \mathcal{T}_{xtz} - B_{0} \, \zeta_{2} \partial_t \mathcal{T}_{xtz} \\
  &\quad - \tfrac{1}{2} B_{0} \, \zeta_{3} \partial_t \mathcal{T}_{xtz} + \tfrac{B_{0} \, d_{19}}{2} \partial_t \mathcal{T}_{xxy} - \tfrac{1}{2} B_{0} \, d_{20} \partial_t \mathcal{T}_{xxy} + 2 \, B_{0} \, d_{21} \partial_t \mathcal{T}_{xxy} - B_{0} \, {\zeta}_{2} \partial_t \mathcal{T}_{xxy} - \tfrac{1}{2} B_{0} \, {\zeta}_{3} \partial_t \mathcal{T}_{xxy} \\
  &\quad - 2 \, B_{0} \, {\zeta}_{5} \partial_t \mathcal{T}_{xxy} - B_{0} \, {\zeta}_{6} \partial_t \mathcal{T}_{xxy}
  \end{aligned} \\[1ex]
  % \begin{aligned}
  %   -\tfrac{1}{\kappa^{2}} \, \partial_t^{2} \mathcal{H}_{xz} &= \tfrac{B_{0}^2}{2} \mathcal{H}_{xz} + \tfrac{B_{0} \, d_{19}}{2} \partial_t \mathcal{T}_{xxz} - \tfrac{1}{2} B_{0} \, d_{20} \partial_t \mathcal{T}_{xxz} + 2 \, B_{0} \, d_{21} \partial_t \mathcal{T}_{xxz} - B_{0} \, {\zeta}_{2} \partial_t \mathcal{T}_{xxz} - \tfrac{1}{2} B_{0} \, {\zeta}_{3} \partial_t \mathcal{T}_{xxz} \\
  % &\quad - 2 \, B_{0} \, {\zeta}_{5} \partial_t \mathcal{T}_{xxz} - B_{0} \, {\zeta}_{6} \partial_t \mathcal{T}_{xxz} - \tfrac{1}{2} B_{0} \, d_{19} \partial_t \mathcal{T}_{ttz} + \tfrac{B_{0} \, d_{20}}{2} \partial_t \mathcal{T}_{ttz} - 2 \, B_{0} \, d_{21} \partial_t \mathcal{T}_{ttz} + B_{0} \, {\zeta}_{2} \partial_t \mathcal{T}_{ttz} \\
  % &\quad + \tfrac{B_{0} \, {\zeta}_{3}}{2} \partial_t \mathcal{T}_{ttz} + 2 \, B_{0} \, {\zeta}_{5} \partial_t \mathcal{T}_{ttz} + B_{0} \, {\zeta}_{6} \partial_t \mathcal{T}_{ttz} + B_{0} \, \delta_{1} \partial_t \mathcal{T}_{txy} - B_{0} \, \zeta_{2} \partial_t \mathcal{T}_{txy} - \tfrac{1}{2} B_{0} \, \zeta_{3} \partial_t \mathcal{T}_{txy} - B_{0} \, \delta_{1} \partial_t \mathcal{T}_{xty} \\
  % &\quad + B_{0} \, \zeta_{2} \partial_t \mathcal{T}_{xty} + \tfrac{B_{0} \, \zeta_{3}}{2} \partial_t \mathcal{T}_{xty} - \tfrac{1}{\kappa^{2}} \partial_z \dot{\mathcal{H}}_{tx}
  % \end{aligned} \\[1ex]
  % \begin{aligned}
  %   0 &= -\tfrac{1}{8} B_{0}^2 \mathcal{H}_{tt} + \tfrac{B_{0}^2}{8} \mathcal{H}_{xx} - \tfrac{3 \, B_{0}^2}{8} \mathcal{H}_{yy} - \tfrac{1}{8} B_{0}^2 \mathcal{H}_{zz} - \tfrac{1}{2} B_{0} \partial_z \mathcal{A}_{y} + \tfrac{3 \, B_{0} \, \delta_{1}}{4} \partial_z \mathcal{T}_{tty} + \tfrac{B_{0} \, \zeta_{1}}{2} \partial_z \mathcal{T}_{tty} - B_{0} \, d_{21} \partial_z \mathcal{T}_{ztx} \\
  % &\quad + B_{0} \, {\zeta}_{1} \partial_z \mathcal{T}_{ztx} + B_{0} \, {\zeta}_{2} \partial_z \mathcal{T}_{ztx} + 2 \, B_{0} \, {\zeta}_{5} \partial_z \mathcal{T}_{ztx} + \tfrac{B_{0} \, \delta_{1}}{4} \partial_z \mathcal{T}_{zyz} - \tfrac{1}{2} B_{0} \, \zeta_{1} \partial_z \mathcal{T}_{zyz} - B_{0} \, \zeta_{2} \partial_z \mathcal{T}_{zyz} \\
  % &\quad - \tfrac{1}{2} B_{0} \, \zeta_{3} \partial_z \mathcal{T}_{zyz} - \tfrac{1}{2} B_{0} \, d_{19} \partial_z \mathcal{T}_{txz} + \tfrac{B_{0} \, d_{20}}{2} \partial_z \mathcal{T}_{txz} - B_{0} \, d_{21} \partial_z \mathcal{T}_{txz} - B_{0} \, {\zeta}_{1} \partial_z \mathcal{T}_{txz} + \tfrac{B_{0} \, {\zeta}_{3}}{2} \partial_z \mathcal{T}_{txz} \\
  % &\quad + B_{0} \, {\zeta}_{6} \partial_z \mathcal{T}_{txz} - \tfrac{1}{2} B_{0} \, d_{19} \partial_z \mathcal{T}_{xtz} + \tfrac{B_{0} \, d_{20}}{2} \partial_z \mathcal{T}_{xtz} - B_{0} \, d_{21} \partial_z \mathcal{T}_{xtz} - B_{0} \, {\zeta}_{1} \partial_z \mathcal{T}_{xtz} + \tfrac{B_{0} \, {\zeta}_{3}}{2} \partial_z \mathcal{T}_{xtz} + B_{0} \, {\zeta}_{6} \partial_z \mathcal{T}_{xtz} \\
  % &\quad + \tfrac{3 \, B_{0} \, \delta_{1}}{4} \partial_z \mathcal{T}_{xxy} + \tfrac{B_{0} \, \zeta_{1}}{2} \partial_z \mathcal{T}_{xxy} + \tfrac{1}{2 \, \kappa^2} \partial_z^2 \mathcal{H}_{tt} - \tfrac{1}{2 \, \kappa^2} \partial_z^2 \mathcal{H}_{xx} - \tfrac{1}{2} B_{0} \, d_{19} \partial_t \mathcal{T}_{ttx} + \tfrac{B_{0} \, d_{20}}{2} \partial_t \mathcal{T}_{ttx} \\
  % &\quad - 2 \, B_{0} \, d_{21} \partial_t \mathcal{T}_{ttx} + B_{0} \, {\zeta}_{2} \partial_t \mathcal{T}_{ttx} + \tfrac{B_{0} \, {\zeta}_{3}}{2} \partial_t \mathcal{T}_{ttx} + 2 \, B_{0} \, {\zeta}_{5} \partial_t \mathcal{T}_{ttx} + B_{0} \, {\zeta}_{6} \partial_t \mathcal{T}_{ttx} - B_{0} \, \delta_{1} \partial_t \mathcal{T}_{ytz} + B_{0} \, \zeta_{2} \partial_t \mathcal{T}_{ytz} \\
  % &\quad + \tfrac{B_{0} \, \zeta_{3}}{2} \partial_t \mathcal{T}_{ytz} - \tfrac{1}{2} B_{0} \, d_{19} \partial_t \mathcal{T}_{yxy} + \tfrac{B_{0} \, d_{20}}{2} \partial_t \mathcal{T}_{yxy} - 2 \, B_{0} \, d_{21} \partial_t \mathcal{T}_{yxy} + B_{0} \, {\zeta}_{2} \partial_t \mathcal{T}_{yxy} + \tfrac{B_{0} \, {\zeta}_{3}}{2} \partial_t \mathcal{T}_{yxy} \\
  % &\quad + 2 \, B_{0} \, {\zeta}_{5} \partial_t \mathcal{T}_{yxy} + B_{0} \, {\zeta}_{6} \partial_t \mathcal{T}_{yxy} + B_{0} \, \delta_{1} \partial_t \mathcal{T}_{tyz} - B_{0} \, \zeta_{2} \partial_t \mathcal{T}_{tyz} - \tfrac{1}{2} B_{0} \, \zeta_{3} \partial_t \mathcal{T}_{tyz} - \tfrac{1}{\kappa^{2}} \partial_z \dot{\mathcal{H}}_{tz} + \tfrac{1}{2 \, \kappa^2} \partial_t^2 \mathcal{H}_{xx} \\
  % &\quad + \tfrac{1}{2 \, \kappa^2} \partial_t^2 \mathcal{H}_{zz}
  % \end{aligned} \\[1ex]
  % \begin{aligned}
  %   -\tfrac{1}{\kappa^{2}} \, \partial_t^{2} \mathcal{H}_{yz} &= \tfrac{B_{0}^2}{2} \mathcal{H}_{yz} - B_{0} \, \delta_{1} \partial_t \mathcal{T}_{yty} + B_{0} \, \zeta_{2} \partial_t \mathcal{T}_{yty} + \tfrac{B_{0} \, \zeta_{3}}{2} \partial_t \mathcal{T}_{yty} + \tfrac{B_{0} \, d_{19}}{2} \partial_t \mathcal{T}_{yxz} - \tfrac{1}{2} B_{0} \, d_{20} \partial_t \mathcal{T}_{yxz} \\
  % &\quad + 2 \, B_{0} \, d_{21} \partial_t \mathcal{T}_{yxz} - B_{0} \, {\zeta}_{2} \partial_t \mathcal{T}_{yxz} - \tfrac{1}{2} B_{0} \, {\zeta}_{3} \partial_t \mathcal{T}_{yxz} - 2 \, B_{0} \, {\zeta}_{5} \partial_t \mathcal{T}_{yxz} - B_{0} \, {\zeta}_{6} \partial_t \mathcal{T}_{yxz} + B_{0} \, \delta_{1} \partial_t \mathcal{T}_{ztz} - B_{0} \, \zeta_{2} \partial_t \mathcal{T}_{ztz} \\
  % &\quad - \tfrac{1}{2} B_{0} \, \zeta_{3} \partial_t \mathcal{T}_{ztz} + \tfrac{B_{0} \, d_{19}}{2} \partial_t \mathcal{T}_{zxy} - \tfrac{1}{2} B_{0} \, d_{20} \partial_t \mathcal{T}_{zxy} + 2 \, B_{0} \, d_{21} \partial_t \mathcal{T}_{zxy} - B_{0} \, {\zeta}_{2} \partial_t \mathcal{T}_{zxy} - \tfrac{1}{2} B_{0} \, {\zeta}_{3} \partial_t \mathcal{T}_{zxy} \\
  % &\quad - 2 \, B_{0} \, {\zeta}_{5} \partial_t \mathcal{T}_{zxy} - B_{0} \, {\zeta}_{6} \partial_t \mathcal{T}_{zxy} - \tfrac{1}{\kappa^{2}} \partial_z \dot{\mathcal{H}}_{ty}
  % \end{aligned} \\[1ex]
  % \begin{aligned}
  %   0 &= -\tfrac{1}{8} B_{0}^2 \mathcal{H}_{tt} + \tfrac{B_{0}^2}{8} \mathcal{H}_{xx} - \tfrac{1}{8} B_{0}^2 \mathcal{H}_{yy} - \tfrac{3 \, B_{0}^2}{8} \mathcal{H}_{zz} - \tfrac{1}{2} B_{0} \partial_z \mathcal{A}_{y} + \tfrac{3 \, B_{0} \, \delta_{1}}{4} \partial_z \mathcal{T}_{tty} + \tfrac{B_{0} \, \zeta_{1}}{2} \partial_z \mathcal{T}_{tty} - B_{0} \, d_{21} \partial_z \mathcal{T}_{ztx} \\
  % &\quad + B_{0} \, {\zeta}_{1} \partial_z \mathcal{T}_{ztx} + B_{0} \, {\zeta}_{2} \partial_z \mathcal{T}_{ztx} + 2 \, B_{0} \, {\zeta}_{5} \partial_z \mathcal{T}_{ztx} + \tfrac{B_{0} \, \delta_{1}}{4} \partial_z \mathcal{T}_{zyz} - \tfrac{1}{2} B_{0} \, \zeta_{1} \partial_z \mathcal{T}_{zyz} - B_{0} \, \zeta_{2} \partial_z \mathcal{T}_{zyz} \\
  % &\quad - \tfrac{1}{2} B_{0} \, \zeta_{3} \partial_z \mathcal{T}_{zyz} - \tfrac{1}{2} B_{0} \, d_{19} \partial_z \mathcal{T}_{txz} + \tfrac{B_{0} \, d_{20}}{2} \partial_z \mathcal{T}_{txz} - B_{0} \, d_{21} \partial_z \mathcal{T}_{txz} - B_{0} \, {\zeta}_{1} \partial_z \mathcal{T}_{txz} + \tfrac{B_{0} \, {\zeta}_{3}}{2} \partial_z \mathcal{T}_{txz} \\
  % &\quad + B_{0} \, {\zeta}_{6} \partial_z \mathcal{T}_{txz} - \tfrac{1}{2} B_{0} \, d_{19} \partial_z \mathcal{T}_{xtz} + \tfrac{B_{0} \, d_{20}}{2} \partial_z \mathcal{T}_{xtz} - B_{0} \, d_{21} \partial_z \mathcal{T}_{xtz} - B_{0} \, {\zeta}_{1} \partial_z \mathcal{T}_{xtz} + \tfrac{B_{0} \, {\zeta}_{3}}{2} \partial_z \mathcal{T}_{xtz} + B_{0} \, {\zeta}_{6} \partial_z \mathcal{T}_{xtz} \\
  % &\quad + \tfrac{3 \, B_{0} \, \delta_{1}}{4} \partial_z \mathcal{T}_{xxy} + \tfrac{B_{0} \, \zeta_{1}}{2} \partial_z \mathcal{T}_{xxy} - \tfrac{1}{2} B_{0} \, d_{19} \partial_t \mathcal{T}_{ttx} + \tfrac{B_{0} \, d_{20}}{2} \partial_t \mathcal{T}_{ttx} - 2 \, B_{0} \, d_{21} \partial_t \mathcal{T}_{ttx} + B_{0} \, {\zeta}_{2} \partial_t \mathcal{T}_{ttx} \\
  % &\quad + \tfrac{B_{0} \, {\zeta}_{3}}{2} \partial_t \mathcal{T}_{ttx} + 2 \, B_{0} \, {\zeta}_{5} \partial_t \mathcal{T}_{ttx} + B_{0} \, {\zeta}_{6} \partial_t \mathcal{T}_{ttx} + B_{0} \, \delta_{1} \partial_t \mathcal{T}_{zty} - B_{0} \, \zeta_{2} \partial_t \mathcal{T}_{zty} - \tfrac{1}{2} B_{0} \, \zeta_{3} \partial_t \mathcal{T}_{zty} \\
  % &\quad - \tfrac{1}{2} B_{0} \, d_{19} \partial_t \mathcal{T}_{zxz} + \tfrac{B_{0} \, d_{20}}{2} \partial_t \mathcal{T}_{zxz} - 2 \, B_{0} \, d_{21} \partial_t \mathcal{T}_{zxz} + B_{0} \, {\zeta}_{2} \partial_t \mathcal{T}_{zxz} + \tfrac{B_{0} \, {\zeta}_{3}}{2} \partial_t \mathcal{T}_{zxz} + 2 \, B_{0} \, {\zeta}_{5} \partial_t \mathcal{T}_{zxz} \\
  % &\quad + B_{0} \, {\zeta}_{6} \partial_t \mathcal{T}_{zxz} + B_{0} \, \delta_{1} \partial_t \mathcal{T}_{tyz} - B_{0} \, \zeta_{2} \partial_t \mathcal{T}_{tyz} - \tfrac{1}{2} B_{0} \, \zeta_{3} \partial_t \mathcal{T}_{tyz} + \tfrac{1}{2 \, \kappa^2} \partial_t^2 \mathcal{H}_{xx} + \tfrac{1}{2 \, \kappa^2} \partial_t^2 \mathcal{H}_{yy}
  % \end{aligned} \\[1ex]
  % \begin{aligned}
  %   \left(-\chi - \xi\right) \, \partial_t^{2} \mathcal{T}_{ttx} &= -4 \, \beta_{1} \mathcal{T}_{ttx} - 2 \, \beta_{2} \mathcal{T}_{ttx} - 2 \, \beta_{3} \mathcal{T}_{ttx} - \tfrac{1}{2 \, \kappa^2} \mathcal{T}_{ttx} - \tfrac{3 \, B_{0} \, \delta_{1}}{2} \mathcal{T}_{xyz} - B_{0} \, d_{19} \mathcal{T}_{yty} + B_{0} \, d_{20} \mathcal{T}_{yty} - 6 \, B_{0} \, d_{21} \mathcal{T}_{yty} \\
  % &\quad - 2 \, d_{17} \mathcal{T}_{ytz} + 2 \, \beta_{3} \mathcal{T}_{yxy} - \tfrac{3}{\kappa^2} \mathcal{T}_{yxy} - \tfrac{3 \, B_{0} \, \delta_{1}}{2} \mathcal{T}_{yxz} + 2 \, d_{17} \mathcal{T}_{zty} - B_{0} \, d_{19} \mathcal{T}_{ztz} + B_{0} \, d_{20} \mathcal{T}_{ztz} - 6 \, B_{0} \, d_{21} \mathcal{T}_{ztz} + \tfrac{3 \, B_{0} \, \delta_{1}}{2} \mathcal{T}_{zxy} \\
  % &\quad + 2 \, \beta_{3} \mathcal{T}_{zxz} - \tfrac{3}{\kappa^2} \mathcal{T}_{zxz} - 8 \, d_{14} \mathcal{T}_{tyz} + 8 \, d_{15} \mathcal{T}_{tyz} - 2 \, d_{17} \mathcal{T}_{tyz} + \tfrac{3 \, B_{0} \, \delta_{1}}{2} \partial_z \mathcal{H}_{xy} + B_{0} \, \zeta_{1} \partial_z \mathcal{H}_{xy} + \tfrac{3 \, \delta_{1}}{2} \partial_z^2 \mathcal{A}_{x} + \zeta_{1} \partial_z^2 \mathcal{A}_{x} \\
  % &\quad + 3 \, \chi \partial_z^2 \mathcal{T}_{ttx} - \xi \partial_z^2 \mathcal{T}_{ttx} - 3 \, \chi \partial_z^2 \mathcal{T}_{yxy} + \xi \partial_z^2 \mathcal{T}_{yxy} - \chi \partial_z^2 \mathcal{T}_{zxz} + \xi \partial_z^2 \mathcal{T}_{zxz} - \tfrac{1}{2} B_{0} \, d_{19} \partial_t \mathcal{H}_{yy} + \tfrac{B_{0} \, d_{20}}{2} \partial_t \mathcal{H}_{yy} \\
  % &\quad - 2 \, B_{0} \, d_{21} \partial_t \mathcal{H}_{yy} + B_{0} \, {\zeta}_{2} \partial_t \mathcal{H}_{yy} + \tfrac{B_{0} \, {\zeta}_{3}}{2} \partial_t \mathcal{H}_{yy} + 2 \, B_{0} \, {\zeta}_{5} \partial_t \mathcal{H}_{yy} + B_{0} \, {\zeta}_{6} \partial_t \mathcal{H}_{yy} - \tfrac{1}{2} B_{0} \, d_{19} \partial_t \mathcal{H}_{zz} + \tfrac{B_{0} \, d_{20}}{2} \partial_t \mathcal{H}_{zz} \\
  % &\quad - 2 \, B_{0} \, d_{21} \partial_t \mathcal{H}_{zz} + B_{0} \, {\zeta}_{2} \partial_t \mathcal{H}_{zz} + \tfrac{B_{0} \, {\zeta}_{3}}{2} \partial_t \mathcal{H}_{zz} + 2 \, B_{0} \, {\zeta}_{5} \partial_t \mathcal{H}_{zz} + B_{0} \, {\zeta}_{6} \partial_t \mathcal{H}_{zz} - d_{19} \partial_z \dot{\mathcal{A}}_{y} + d_{20} \partial_z \dot{\mathcal{A}}_{y} \\
  % &\quad - 4 \, d_{21} \partial_z \dot{\mathcal{A}}_{y} + 2 \, {\zeta}_{2} \partial_z \dot{\mathcal{A}}_{y} + {\zeta}_{3} \partial_z \dot{\mathcal{A}}_{y} + 4 \, {\zeta}_{5} \partial_z \dot{\mathcal{A}}_{y} + 2 \, {\zeta}_{6} \partial_z \dot{\mathcal{A}}_{y} + 2 \, \chi \partial_z \dot{\mathcal{T}}_{ztx} + 2 \, \chi \partial_z \dot{\mathcal{T}}_{txz} \\
  % &\quad + \tfrac{\delta_{1}}{2} \partial_t^2 \mathcal{A}_{x} - \zeta_{1} \partial_t^2 \mathcal{A}_{x} - 2 \, \zeta_{2} \partial_t^2 \mathcal{A}_{x} - \zeta_{3} \partial_t^2 \mathcal{A}_{x} + \chi \partial_t^2 \mathcal{T}_{yxy} - \xi \partial_t^2 \mathcal{T}_{yxy} + \chi \partial_t^2 \mathcal{T}_{zxz} - \xi \partial_t^2 \mathcal{T}_{zxz}
  % \end{aligned} \\[1ex]
  % \begin{aligned}
  %   \left(-\chi - \xi\right) \, \partial_t^{2} \mathcal{T}_{tty} &= -4 \, \beta_{1} \mathcal{T}_{tty} - 2 \, \beta_{2} \mathcal{T}_{tty} - 2 \, \beta_{3} \mathcal{T}_{tty} - \tfrac{1}{2 \, \kappa^2} \mathcal{T}_{tty} + \tfrac{3 \, B_{0} \, d_{19}}{2} \mathcal{T}_{ytx} - \tfrac{3 \, B_{0} \, d_{20}}{2} \mathcal{T}_{ytx} + 4 \, B_{0} \, d_{21} \mathcal{T}_{ytx} \\
  % &\quad - 3 \, B_{0} \, \delta_{1} \mathcal{T}_{yyz} - 2 \, d_{17} \mathcal{T}_{ztx} + 2 \, \beta_{3} \mathcal{T}_{zyz} - \tfrac{3}{\kappa^2} \mathcal{T}_{zyz} - \tfrac{1}{2} B_{0} \, d_{19} \mathcal{T}_{txy} + \tfrac{B_{0} \, d_{20}}{2} \mathcal{T}_{txy} + 2 \, B_{0} \, d_{21} \mathcal{T}_{txy} + 8 \, d_{14} \mathcal{T}_{txz} - 8 \, d_{15} \mathcal{T}_{txz} + 2 \, d_{17} \mathcal{T}_{txz} \\
  % &\quad - \tfrac{1}{2} B_{0} \, d_{19} \mathcal{T}_{xty} + \tfrac{B_{0} \, d_{20}}{2} \mathcal{T}_{xty} + 2 \, B_{0} \, d_{21} \mathcal{T}_{xty} + 2 \, d_{17} \mathcal{T}_{xtz} - 2 \, \beta_{3} \mathcal{T}_{xxy} + \tfrac{3}{\kappa^2} \mathcal{T}_{xxy} - \tfrac{1}{4} B_{0} \, \delta_{1} \partial_z \mathcal{H}_{tt} + \tfrac{B_{0} \, \zeta_{1}}{2} \partial_z \mathcal{H}_{tt} \\
  % &\quad + B_{0} \, \zeta_{2} \partial_z \mathcal{H}_{tt} + \tfrac{B_{0} \, \zeta_{3}}{2} \partial_z \mathcal{H}_{tt} - \tfrac{3 \, B_{0} \, \delta_{1}}{4} \partial_z \mathcal{H}_{xx} - \tfrac{1}{2} B_{0} \, \zeta_{1} \partial_z \mathcal{H}_{xx} + \tfrac{3 \, B_{0} \, \delta_{1}}{4} \partial_z \mathcal{H}_{yy} + \tfrac{B_{0} \, \zeta_{1}}{2} \partial_z \mathcal{H}_{yy} \\
  % &\quad + \tfrac{3 \, B_{0} \, \delta_{1}}{4} \partial_z \mathcal{H}_{zz} + \tfrac{B_{0} \, \zeta_{1}}{2} \partial_z \mathcal{H}_{zz} + \tfrac{3 \, \delta_{1}}{2} \partial_z^2 \mathcal{A}_{y} + \zeta_{1} \partial_z^2 \mathcal{A}_{y} + 3 \, \chi \partial_z^2 \mathcal{T}_{tty} - \xi \partial_z^2 \mathcal{T}_{tty} - \chi \partial_z^2 \mathcal{T}_{zyz} + \xi \partial_z^2 \mathcal{T}_{zyz} \\
  % &\quad + 3 \, \chi \partial_z^2 \mathcal{T}_{xxy} - \xi \partial_z^2 \mathcal{T}_{xxy} + \tfrac{B_{0} \, \delta_{1}}{2} \partial_t \mathcal{H}_{tz} - B_{0} \, \zeta_{1} \partial_t \mathcal{H}_{tz} - 2 \, B_{0} \, \zeta_{2} \partial_t \mathcal{H}_{tz} - B_{0} \, \zeta_{3} \partial_t \mathcal{H}_{tz} + \tfrac{B_{0} \, d_{19}}{2} \partial_t \mathcal{H}_{xy} \\
  % &\quad - \tfrac{1}{2} B_{0} \, d_{20} \partial_t \mathcal{H}_{xy} + 2 \, B_{0} \, d_{21} \partial_t \mathcal{H}_{xy} - B_{0} \, {\zeta}_{2} \partial_t \mathcal{H}_{xy} - \tfrac{1}{2} B_{0} \, {\zeta}_{3} \partial_t \mathcal{H}_{xy} - 2 \, B_{0} \, {\zeta}_{5} \partial_t \mathcal{H}_{xy} - B_{0} \, {\zeta}_{6} \partial_t \mathcal{H}_{xy} + d_{19} \partial_z \dot{\mathcal{A}}_{x} \\
  % &\quad - d_{20} \partial_z \dot{\mathcal{A}}_{x} + 4 \, d_{21} \partial_z \dot{\mathcal{A}}_{x} - 2 \, {\zeta}_{2} \partial_z \dot{\mathcal{A}}_{x} - {\zeta}_{3} \partial_z \dot{\mathcal{A}}_{x} - 4 \, {\zeta}_{5} \partial_z \dot{\mathcal{A}}_{x} - 2 \, {\zeta}_{6} \partial_z \dot{\mathcal{A}}_{x} + 2 \, \chi \partial_z \dot{\mathcal{T}}_{zty} + 2 \, \chi \partial_z \dot{\mathcal{T}}_{tyz} \\
  % &\quad + \tfrac{\delta_{1}}{2} \partial_t^2 \mathcal{A}_{y} - \zeta_{1} \partial_t^2 \mathcal{A}_{y} - 2 \, \zeta_{2} \partial_t^2 \mathcal{A}_{y} - \zeta_{3} \partial_t^2 \mathcal{A}_{y} + \chi \partial_t^2 \mathcal{T}_{zyz} - \xi \partial_t^2 \mathcal{T}_{zyz} - \chi \partial_t^2 \mathcal{T}_{xxy} + \xi \partial_t^2 \mathcal{T}_{xxy}
  % \end{aligned} \\[1ex]
  % \begin{aligned}
  %   \left(3 \, \chi - \xi\right) \, \partial_t^{2} \mathcal{T}_{xxz} &= -4 \, \beta_{1} \mathcal{T}_{xxz} - 2 \, \beta_{2} \mathcal{T}_{xxz} - 2 \, \beta_{3} \mathcal{T}_{xxz} - \tfrac{1}{2 \, \kappa^2} \mathcal{T}_{xxz} + 2 \, d_{17} \mathcal{T}_{ytx} - 2 \, \beta_{3} \mathcal{T}_{yyz} + \tfrac{3}{\kappa^2} \mathcal{T}_{yyz} + \tfrac{3 \, B_{0} \, d_{19}}{2} \mathcal{T}_{ztx} \\
  % &\quad - \tfrac{3 \, B_{0} \, d_{20}}{2} \mathcal{T}_{ztx} + 4 \, B_{0} \, d_{21} \mathcal{T}_{ztx} - 2 \, \beta_{3} \mathcal{T}_{ttz} + \tfrac{3}{\kappa^2} \mathcal{T}_{ttz} - 3 \, B_{0} \, \delta_{1} \mathcal{T}_{zyz} - 2 \, d_{17} \mathcal{T}_{txy} - \tfrac{1}{2} B_{0} \, d_{19} \mathcal{T}_{txz} + \tfrac{B_{0} \, d_{20}}{2} \mathcal{T}_{txz} + 2 \, B_{0} \, d_{21} \mathcal{T}_{txz} \\
  % &\quad - 8 \, d_{14} \mathcal{T}_{xty} + 8 \, d_{15} \mathcal{T}_{xty} - 2 \, d_{17} \mathcal{T}_{xty} - \tfrac{1}{2} B_{0} \, d_{19} \mathcal{T}_{xtz} + \tfrac{B_{0} \, d_{20}}{2} \mathcal{T}_{xtz} + 2 \, B_{0} \, d_{21} \mathcal{T}_{xtz} + \tfrac{B_{0} \, d_{19}}{2} \partial_z \mathcal{H}_{tx} - \tfrac{1}{2} B_{0} \, d_{20} \partial_z \mathcal{H}_{tx} \\
  % &\quad + 2 \, B_{0} \, d_{21} \partial_z \mathcal{H}_{tx} - B_{0} \, {\zeta}_{2} \partial_z \mathcal{H}_{tx} - \tfrac{1}{2} B_{0} \, {\zeta}_{3} \partial_z \mathcal{H}_{tx} - 2 \, B_{0} \, {\zeta}_{5} \partial_z \mathcal{H}_{tx} - B_{0} \, {\zeta}_{6} \partial_z \mathcal{H}_{tx} + \tfrac{3 \, B_{0} \, \delta_{1}}{2} \partial_t \mathcal{H}_{ty} + B_{0} \, \zeta_{1} \partial_t \mathcal{H}_{ty} \\
  % &\quad - \tfrac{1}{2} B_{0} \, d_{19} \partial_t \mathcal{H}_{xz} + \tfrac{B_{0} \, d_{20}}{2} \partial_t \mathcal{H}_{xz} - 2 \, B_{0} \, d_{21} \partial_t \mathcal{H}_{xz} + B_{0} \, {\zeta}_{2} \partial_t \mathcal{H}_{xz} + \tfrac{B_{0} \, {\zeta}_{3}}{2} \partial_t \mathcal{H}_{xz} + 2 \, B_{0} \, {\zeta}_{5} \partial_t \mathcal{H}_{xz} + B_{0} \, {\zeta}_{6} \partial_t \mathcal{H}_{xz} \\
  % &\quad + \tfrac{3 \, \delta_{1}}{2} \partial_z \dot{\mathcal{A}}_{t} + \zeta_{1} \partial_z \dot{\mathcal{A}}_{t} - 3 \, \chi \partial_z \dot{\mathcal{T}}_{yty} + \xi \partial_z \dot{\mathcal{T}}_{yty} - \chi \partial_z \dot{\mathcal{T}}_{ztz} + \xi \partial_z \dot{\mathcal{T}}_{ztz} - 3 \, \chi \partial_z \dot{\mathcal{T}}_{xtx} + \xi \partial_z \dot{\mathcal{T}}_{xtx} \\
  % &\quad - \tfrac{3 \, \delta_{1}}{2} \partial_t^2 \mathcal{A}_{z} - \zeta_{1} \partial_t^2 \mathcal{A}_{z} - 3 \, \chi \partial_t^2 \mathcal{T}_{yyz} + \xi \partial_t^2 \mathcal{T}_{yyz} - \chi \partial_t^2 \mathcal{T}_{ttz} + \xi \partial_t^2 \mathcal{T}_{ttz}
  % \end{aligned} \\[1ex]
  % \begin{aligned}
  %   0 &= \tfrac{3 \, B_{0} \, \delta_{1}}{2} \mathcal{T}_{ttx} + 4 \, \beta_{1} \mathcal{T}_{xyz} + \tfrac{2}{\kappa^2} \mathcal{T}_{xyz} - 2 \, d_{17} \mathcal{T}_{yty} + \tfrac{3 \, B_{0} \, d_{19}}{2} \mathcal{T}_{ytz} - \tfrac{3 \, B_{0} \, d_{20}}{2} \mathcal{T}_{ytz} + B_{0} \, d_{21} \mathcal{T}_{ytz} - 2 \, B_{0} \, \delta_{1} \mathcal{T}_{yxy} + 2 \, \beta_{2} \mathcal{T}_{yxz} \\
  % &\quad + \tfrac{3}{2 \, \kappa^2} \mathcal{T}_{yxz} - \tfrac{3 \, B_{0} \, d_{19}}{2} \mathcal{T}_{zty} + \tfrac{3 \, B_{0} \, d_{20}}{2} \mathcal{T}_{zty} - B_{0} \, d_{21} \mathcal{T}_{zty} - 2 \, d_{17} \mathcal{T}_{ztz} - 2 \, \beta_{2} \mathcal{T}_{zxy} - \tfrac{3}{2 \, \kappa^2} \mathcal{T}_{zxy} - 2 \, B_{0} \, \delta_{1} \mathcal{T}_{zxz} - 8 \, d_{14} \mathcal{T}_{xtx} \\
  % &\quad + 8 \, d_{15} \mathcal{T}_{xtx} - 2 \, d_{17} \mathcal{T}_{xtx} - B_{0} \, d_{19} \partial_z \mathcal{H}_{ty} + B_{0} \, d_{20} \partial_z \mathcal{H}_{ty} - 2 \, B_{0} \, d_{21} \partial_z \mathcal{H}_{ty} - 2 \, B_{0} \, {\zeta}_{1} \partial_z \mathcal{H}_{ty} + B_{0} \, {\zeta}_{3} \partial_z \mathcal{H}_{ty} + 2 \, B_{0} \, {\zeta}_{6} \partial_z \mathcal{H}_{ty} \\
  % &\quad - d_{19} \partial_z^2 \mathcal{A}_{t} + d_{20} \partial_z^2 \mathcal{A}_{t} - 2 \, d_{21} \partial_z^2 \mathcal{A}_{t} - 2 \, {\zeta}_{1} \partial_z^2 \mathcal{A}_{t} + {\zeta}_{3} \partial_z^2 \mathcal{A}_{t} + 2 \, {\zeta}_{6} \partial_z^2 \mathcal{A}_{t} + d_{19} \partial_z \dot{\mathcal{A}}_{z} - d_{20} \partial_z \dot{\mathcal{A}}_{z} + 2 \, d_{21} \partial_z \dot{\mathcal{A}}_{z} \\
  % &\quad + 2 \, {\zeta}_{1} \partial_z \dot{\mathcal{A}}_{z} - {\zeta}_{3} \partial_z \dot{\mathcal{A}}_{z} - 2 \, {\zeta}_{6} \partial_z \dot{\mathcal{A}}_{z}
  % \end{aligned} \\[1ex]
  % \begin{aligned}
  %   0 &= -\tfrac{3 \, B_{0} \, d_{19}}{2} \mathcal{T}_{tty} + \tfrac{3 \, B_{0} \, d_{20}}{2} \mathcal{T}_{tty} - 4 \, B_{0} \, d_{21} \mathcal{T}_{tty} - 2 \, d_{17} \mathcal{T}_{xxz} - 4 \, \beta_{1} \mathcal{T}_{ytx} - \tfrac{2}{\kappa^2} \mathcal{T}_{ytx} - 8 \, d_{14} \mathcal{T}_{yyz} + 8 \, d_{15} \mathcal{T}_{yyz} - 2 \, d_{17} \mathcal{T}_{yyz} - 2 \, d_{17} \mathcal{T}_{ttz} \\
  % &\quad + 3 \, B_{0} \, d_{21} \mathcal{T}_{zyz} - 2 \, \beta_{2} \mathcal{T}_{txy} - \tfrac{3}{2 \, \kappa^2} \mathcal{T}_{txy} - \tfrac{1}{2} B_{0} \, \delta_{1} \mathcal{T}_{txz} - 2 \, \beta_{2} \mathcal{T}_{xty} - \tfrac{3}{2 \, \kappa^2} \mathcal{T}_{xty} - \tfrac{1}{2} B_{0} \, \delta_{1} \mathcal{T}_{xtz} - \tfrac{3 \, B_{0} \, d_{19}}{2} \mathcal{T}_{xxy} \\
  % &\quad + \tfrac{3 \, B_{0} \, d_{20}}{2} \mathcal{T}_{xxy} - 4 \, B_{0} \, d_{21} \mathcal{T}_{xxy} + B_{0} \, d_{19} \partial_t \mathcal{H}_{ty} - B_{0} \, d_{20} \partial_t \mathcal{H}_{ty} + 2 \, B_{0} \, d_{21} \partial_t \mathcal{H}_{ty} + 2 \, B_{0} \, {\zeta}_{1} \partial_t \mathcal{H}_{ty} - B_{0} \, {\zeta}_{3} \partial_t \mathcal{H}_{ty} \\
  % &\quad - 2 \, B_{0} \, {\zeta}_{6} \partial_t \mathcal{H}_{ty} + d_{19} \partial_z \dot{\mathcal{A}}_{t} - d_{20} \partial_z \dot{\mathcal{A}}_{t} + 2 \, d_{21} \partial_z \dot{\mathcal{A}}_{t} + 2 \, {\zeta}_{1} \partial_z \dot{\mathcal{A}}_{t} - {\zeta}_{3} \partial_z \dot{\mathcal{A}}_{t} - 2 \, {\zeta}_{6} \partial_z \dot{\mathcal{A}}_{t} - d_{19} \partial_t^2 \mathcal{A}_{z} \\
  % &\quad + d_{20} \partial_t^2 \mathcal{A}_{z} - 2 \, d_{21} \partial_t^2 \mathcal{A}_{z} - 2 \, {\zeta}_{1} \partial_t^2 \mathcal{A}_{z} + {\zeta}_{3} \partial_t^2 \mathcal{A}_{z} + 2 \, {\zeta}_{6} \partial_t^2 \mathcal{A}_{z}
  % \end{aligned} \\[1ex]
  % \begin{aligned}
  %   0 &= B_{0} \, d_{19} \mathcal{T}_{ttx} - B_{0} \, d_{20} \mathcal{T}_{ttx} + 6 \, B_{0} \, d_{21} \mathcal{T}_{ttx} - 2 \, d_{17} \mathcal{T}_{xyz} - 4 \, \beta_{1} \mathcal{T}_{yty} - 2 \, \beta_{2} \mathcal{T}_{yty} - 2 \, \beta_{3} \mathcal{T}_{yty} - \tfrac{1}{2 \, \kappa^2} \mathcal{T}_{yty} + B_{0} \, \delta_{1} \mathcal{T}_{ytz} + 8 \, d_{14} \mathcal{T}_{yxz} \\
  % &\quad - 8 \, d_{15} \mathcal{T}_{yxz} + 2 \, d_{17} \mathcal{T}_{yxz} - B_{0} \, \delta_{1} \mathcal{T}_{zty} - 2 \, \beta_{3} \mathcal{T}_{ztz} + \tfrac{3}{\kappa^2} \mathcal{T}_{ztz} - 2 \, d_{17} \mathcal{T}_{zxy} - 2 \, B_{0} \, \delta_{1} \mathcal{T}_{tyz} - 2 \, \beta_{3} \mathcal{T}_{xtx} + \tfrac{3}{\kappa^2} \mathcal{T}_{xtx} \\
  % &\quad - \tfrac{1}{2} B_{0} \, \delta_{1} \partial_z \mathcal{H}_{ty} - B_{0} \, \zeta_{1} \partial_z \mathcal{H}_{ty} - B_{0} \, \zeta_{2} \partial_z \mathcal{H}_{ty} - \tfrac{1}{2} B_{0} \, \zeta_{3} \partial_z \mathcal{H}_{ty} - \tfrac{3 \, \delta_{1}}{2} \partial_z^2 \mathcal{A}_{t} - \zeta_{1} \partial_z^2 \mathcal{A}_{t} + 3 \, \chi \partial_z^2 \mathcal{T}_{yty} \\
  % &\quad - \xi \partial_z^2 \mathcal{T}_{yty} + \chi \partial_z^2 \mathcal{T}_{ztz} - \xi \partial_z^2 \mathcal{T}_{ztz} + 3 \, \chi \partial_z^2 \mathcal{T}_{xtx} - \xi \partial_z^2 \mathcal{T}_{xtx} - B_{0} \, \delta_{1} \partial_t \mathcal{H}_{yz} + B_{0} \, \zeta_{2} \partial_t \mathcal{H}_{yz} + \tfrac{B_{0} \, \zeta_{3}}{2} \partial_t \mathcal{H}_{yz} \\
  % &\quad + \tfrac{3 \, \delta_{1}}{2} \partial_z \dot{\mathcal{A}}_{z} + \zeta_{1} \partial_z \dot{\mathcal{A}}_{z} + 3 \, \chi \partial_z \dot{\mathcal{T}}_{xxz} - \xi \partial_z \dot{\mathcal{T}}_{xxz} + 3 \, \chi \partial_z \dot{\mathcal{T}}_{yyz} - \xi \partial_z \dot{\mathcal{T}}_{yyz} + \chi \partial_z \dot{\mathcal{T}}_{ttz} - \xi \partial_z \dot{\mathcal{T}}_{ttz}
  % \end{aligned} \\[1ex]
  % \begin{aligned}
  %   0 &= 2 \, d_{17} \mathcal{T}_{ttx} + \tfrac{3 \, B_{0} \, d_{19}}{2} \mathcal{T}_{xyz} - \tfrac{3 \, B_{0} \, d_{20}}{2} \mathcal{T}_{xyz} + B_{0} \, d_{21} \mathcal{T}_{xyz} + B_{0} \, \delta_{1} \mathcal{T}_{yty} - 4 \, \beta_{1} \mathcal{T}_{ytz} - \tfrac{2}{\kappa^2} \mathcal{T}_{ytz} - 8 \, d_{14} \mathcal{T}_{yxy} + 8 \, d_{15} \mathcal{T}_{yxy} - 2 \, d_{17} \mathcal{T}_{yxy} \\
  % &\quad - 2 \, \beta_{2} \mathcal{T}_{zty} - \tfrac{3}{2 \, \kappa^2} \mathcal{T}_{zty} + B_{0} \, \delta_{1} \mathcal{T}_{ztz} - 2 \, d_{17} \mathcal{T}_{zxz} + 2 \, \beta_{2} \mathcal{T}_{tyz} + \tfrac{3}{2 \, \kappa^2} \mathcal{T}_{tyz} + \tfrac{3 \, B_{0} \, \delta_{1}}{2} \mathcal{T}_{xtx} - \tfrac{1}{2} B_{0} \, d_{19} \partial_z \mathcal{H}_{xy} \\
  % &\quad + \tfrac{B_{0} \, d_{20}}{2} \partial_z \mathcal{H}_{xy} - 2 \, B_{0} \, {\zeta}_{1} \partial_z \mathcal{H}_{xy} - B_{0} \, {\zeta}_{2} \partial_z \mathcal{H}_{xy} + \tfrac{B_{0} \, {\zeta}_{3}}{2} \partial_z \mathcal{H}_{xy} - 2 \, B_{0} \, {\zeta}_{5} \partial_z \mathcal{H}_{xy} + B_{0} \, {\zeta}_{6} \partial_z \mathcal{H}_{xy} - d_{19} \partial_z^2 \mathcal{A}_{x} \\
  % &\quad + d_{20} \partial_z^2 \mathcal{A}_{x} - 2 \, d_{21} \partial_z^2 \mathcal{A}_{x} - 2 \, {\zeta}_{1} \partial_z^2 \mathcal{A}_{x} + {\zeta}_{3} \partial_z^2 \mathcal{A}_{x} + 2 \, {\zeta}_{6} \partial_z^2 \mathcal{A}_{x} + B_{0} \, \delta_{1} \partial_t \mathcal{H}_{yy} - B_{0} \, \zeta_{2} \partial_t \mathcal{H}_{yy} - \tfrac{1}{2} B_{0} \, \zeta_{3} \partial_t \mathcal{H}_{yy} \\
  % &\quad + d_{19} \partial_t^2 \mathcal{A}_{x} - d_{20} \partial_t^2 \mathcal{A}_{x} + 2 \, d_{21} \partial_t^2 \mathcal{A}_{x} + 2 \, {\zeta}_{1} \partial_t^2 \mathcal{A}_{x} - {\zeta}_{3} \partial_t^2 \mathcal{A}_{x} - 2 \, {\zeta}_{6} \partial_t^2 \mathcal{A}_{x}
  % \end{aligned} \\[1ex]
  % \begin{aligned}
  %   \left(3 \, \chi - \xi\right) \, \partial_t^{2} \mathcal{T}_{yxy} &= 2 \, \beta_{3} \mathcal{T}_{ttx} - \tfrac{3}{\kappa^2} \mathcal{T}_{ttx} + 2 \, B_{0} \, \delta_{1} \mathcal{T}_{xyz} + 8 \, d_{14} \mathcal{T}_{ytz} - 8 \, d_{15} \mathcal{T}_{ytz} + 2 \, d_{17} \mathcal{T}_{ytz} - 4 \, \beta_{1} \mathcal{T}_{yxy} - 2 \, \beta_{2} \mathcal{T}_{yxy} - 2 \, \beta_{3} \mathcal{T}_{yxy} \\
  % &\quad - \tfrac{1}{2 \, \kappa^2} \mathcal{T}_{yxy} + B_{0} \, \delta_{1} \mathcal{T}_{yxz} - 2 \, d_{17} \mathcal{T}_{zty} - B_{0} \, \delta_{1} \mathcal{T}_{zxy} - 2 \, \beta_{3} \mathcal{T}_{zxz} + \tfrac{3}{\kappa^2} \mathcal{T}_{zxz} + 2 \, d_{17} \mathcal{T}_{tyz} - B_{0} \, d_{19} \mathcal{T}_{xtx} + B_{0} \, d_{20} \mathcal{T}_{xtx} - 6 \, B_{0} \, d_{21} \mathcal{T}_{xtx} \\
  % &\quad - \tfrac{1}{2} B_{0} \, \delta_{1} \partial_z \mathcal{H}_{xy} - B_{0} \, \zeta_{1} \partial_z \mathcal{H}_{xy} - B_{0} \, \zeta_{2} \partial_z \mathcal{H}_{xy} - \tfrac{1}{2} B_{0} \, \zeta_{3} \partial_z \mathcal{H}_{xy} - \tfrac{3 \, \delta_{1}}{2} \partial_z^2 \mathcal{A}_{x} - \zeta_{1} \partial_z^2 \mathcal{A}_{x} - 3 \, \chi \partial_z^2 \mathcal{T}_{ttx} \\
  % &\quad + \xi \partial_z^2 \mathcal{T}_{ttx} + 3 \, \chi \partial_z^2 \mathcal{T}_{yxy} - \xi \partial_z^2 \mathcal{T}_{yxy} + \chi \partial_z^2 \mathcal{T}_{zxz} - \xi \partial_z^2 \mathcal{T}_{zxz} - \tfrac{1}{2} B_{0} \, d_{19} \partial_t \mathcal{H}_{yy} + \tfrac{B_{0} \, d_{20}}{2} \partial_t \mathcal{H}_{yy} - 2 \, B_{0} \, d_{21} \partial_t \mathcal{H}_{yy} \\
  % &\quad + B_{0} \, {\zeta}_{2} \partial_t \mathcal{H}_{yy} + \tfrac{B_{0} \, {\zeta}_{3}}{2} \partial_t \mathcal{H}_{yy} + 2 \, B_{0} \, {\zeta}_{5} \partial_t \mathcal{H}_{yy} + B_{0} \, {\zeta}_{6} \partial_t \mathcal{H}_{yy} - 2 \, \chi \partial_z \dot{\mathcal{T}}_{ztx} - 2 \, \chi \partial_z \dot{\mathcal{T}}_{txz} + \tfrac{3 \, \delta_{1}}{2} \partial_t^2 \mathcal{A}_{x} \\
  % &\quad + \zeta_{1} \partial_t^2 \mathcal{A}_{x} + \chi \partial_t^2 \mathcal{T}_{ttx} - \xi \partial_t^2 \mathcal{T}_{ttx} - 3 \, \chi \partial_t^2 \mathcal{T}_{zxz} + \xi \partial_t^2 \mathcal{T}_{zxz}
  % \end{aligned} \\[1ex]
  % \begin{aligned}
  %   0 &= \tfrac{3 \, B_{0} \, \delta_{1}}{2} \mathcal{T}_{ttx} + 2 \, \beta_{2} \mathcal{T}_{xyz} + \tfrac{3}{2 \, \kappa^2} \mathcal{T}_{xyz} + 8 \, d_{14} \mathcal{T}_{yty} - 8 \, d_{15} \mathcal{T}_{yty} + 2 \, d_{17} \mathcal{T}_{yty} - B_{0} \, \delta_{1} \mathcal{T}_{yxy} + 4 \, \beta_{1} \mathcal{T}_{yxz} + \tfrac{2}{\kappa^2} \mathcal{T}_{yxz} + 2 \, d_{17} \mathcal{T}_{ztz} \\
  % &\quad + 2 \, \beta_{2} \mathcal{T}_{zxy} + \tfrac{3}{2 \, \kappa^2} \mathcal{T}_{zxy} - B_{0} \, \delta_{1} \mathcal{T}_{zxz} + \tfrac{3 \, B_{0} \, d_{19}}{2} \mathcal{T}_{tyz} - \tfrac{3 \, B_{0} \, d_{20}}{2} \mathcal{T}_{tyz} + B_{0} \, d_{21} \mathcal{T}_{tyz} + 2 \, d_{17} \mathcal{T}_{xtx} + \tfrac{B_{0} \, d_{19}}{2} \partial_z \mathcal{H}_{ty} \\
  % &\quad - \tfrac{1}{2} B_{0} \, d_{20} \partial_z \mathcal{H}_{ty} + 2 \, B_{0} \, {\zeta}_{1} \partial_z \mathcal{H}_{ty} + B_{0} \, {\zeta}_{2} \partial_z \mathcal{H}_{ty} - \tfrac{1}{2} B_{0} \, {\zeta}_{3} \partial_z \mathcal{H}_{ty} + 2 \, B_{0} \, {\zeta}_{5} \partial_z \mathcal{H}_{ty} - B_{0} \, {\zeta}_{6} \partial_z \mathcal{H}_{ty} + d_{19} \partial_z^2 \mathcal{A}_{t} \\
  % &\quad - d_{20} \partial_z^2 \mathcal{A}_{t} + 2 \, d_{21} \partial_z^2 \mathcal{A}_{t} + 2 \, {\zeta}_{1} \partial_z^2 \mathcal{A}_{t} - {\zeta}_{3} \partial_z^2 \mathcal{A}_{t} - 2 \, {\zeta}_{6} \partial_z^2 \mathcal{A}_{t} + \tfrac{B_{0} \, d_{19}}{2} \partial_t \mathcal{H}_{yz} - \tfrac{1}{2} B_{0} \, d_{20} \partial_t \mathcal{H}_{yz} + 2 \, B_{0} \, d_{21} \partial_t \mathcal{H}_{yz} \\
  % &\quad - B_{0} \, {\zeta}_{2} \partial_t \mathcal{H}_{yz} - \tfrac{1}{2} B_{0} \, {\zeta}_{3} \partial_t \mathcal{H}_{yz} - 2 \, B_{0} \, {\zeta}_{5} \partial_t \mathcal{H}_{yz} - B_{0} \, {\zeta}_{6} \partial_t \mathcal{H}_{yz} - d_{19} \partial_z \dot{\mathcal{A}}_{z} + d_{20} \partial_z \dot{\mathcal{A}}_{z} - 2 \, d_{21} \partial_z \dot{\mathcal{A}}_{z} \\
  % &\quad - 2 \, {\zeta}_{1} \partial_z \dot{\mathcal{A}}_{z} + {\zeta}_{3} \partial_z \dot{\mathcal{A}}_{z} + 2 \, {\zeta}_{6} \partial_z \dot{\mathcal{A}}_{z}
  % \end{aligned} \\[1ex]
  % \begin{aligned}
  %   \left(3 \, \chi - \xi\right) \, \partial_t^{2} \mathcal{T}_{yyz} &= -3 \, B_{0} \, \delta_{1} \mathcal{T}_{tty} - 2 \, \beta_{3} \mathcal{T}_{xxz} + \tfrac{3}{\kappa^2} \mathcal{T}_{xxz} + 8 \, d_{14} \mathcal{T}_{ytx} - 8 \, d_{15} \mathcal{T}_{ytx} + 2 \, d_{17} \mathcal{T}_{ytx} - 4 \, \beta_{1} \mathcal{T}_{yyz} - 2 \, \beta_{2} \mathcal{T}_{yyz} - 2 \, \beta_{3} \mathcal{T}_{yyz} \\
  % &\quad - \tfrac{1}{2 \, \kappa^2} \mathcal{T}_{yyz} + 3 \, B_{0} \, d_{21} \mathcal{T}_{ztx} - 2 \, \beta_{3} \mathcal{T}_{ttz} + \tfrac{3}{\kappa^2} \mathcal{T}_{ttz} - 2 \, d_{17} \mathcal{T}_{txy} + B_{0} \, d_{19} \mathcal{T}_{txz} - B_{0} \, d_{20} \mathcal{T}_{txz} + 3 \, B_{0} \, d_{21} \mathcal{T}_{txz} - 2 \, d_{17} \mathcal{T}_{xty} + B_{0} \, d_{19} \mathcal{T}_{xtz} \\
  % &\quad - B_{0} \, d_{20} \mathcal{T}_{xtz} + 3 \, B_{0} \, d_{21} \mathcal{T}_{xtz} - 3 \, B_{0} \, \delta_{1} \mathcal{T}_{xxy} + \tfrac{3 \, B_{0} \, \delta_{1}}{2} \partial_t \mathcal{H}_{ty} + B_{0} \, \zeta_{1} \partial_t \mathcal{H}_{ty} + \tfrac{3 \, \delta_{1}}{2} \partial_z \dot{\mathcal{A}}_{t} + \zeta_{1} \partial_z \dot{\mathcal{A}}_{t} - 3 \, \chi \partial_z \dot{\mathcal{T}}_{yty} \\
  % &\quad + \xi \partial_z \dot{\mathcal{T}}_{yty} - \chi \partial_z \dot{\mathcal{T}}_{ztz} + \xi \partial_z \dot{\mathcal{T}}_{ztz} - 3 \, \chi \partial_z \dot{\mathcal{T}}_{xtx} + \xi \partial_z \dot{\mathcal{T}}_{xtx} - \tfrac{3 \, \delta_{1}}{2} \partial_t^2 \mathcal{A}_{z} - \zeta_{1} \partial_t^2 \mathcal{A}_{z} - 3 \, \chi \partial_t^2 \mathcal{T}_{xxz} + \xi \partial_t^2 \mathcal{T}_{xxz} \\
  % &\quad - \chi \partial_t^2 \mathcal{T}_{ttz} + \xi \partial_t^2 \mathcal{T}_{ttz}
  % \end{aligned} \\[1ex]
  % \begin{aligned}
  %   0 &= 2 \, d_{17} \mathcal{T}_{tty} - \tfrac{3 \, B_{0} \, d_{19}}{2} \mathcal{T}_{xxz} + \tfrac{3 \, B_{0} \, d_{20}}{2} \mathcal{T}_{xxz} - 4 \, B_{0} \, d_{21} \mathcal{T}_{xxz} - 3 \, B_{0} \, d_{21} \mathcal{T}_{yyz} - 4 \, \beta_{1} \mathcal{T}_{ztx} - \tfrac{2}{\kappa^2} \mathcal{T}_{ztx} - \tfrac{3 \, B_{0} \, d_{19}}{2} \mathcal{T}_{ttz} \\
  % &\quad + \tfrac{3 \, B_{0} \, d_{20}}{2} \mathcal{T}_{ttz} - 4 \, B_{0} \, d_{21} \mathcal{T}_{ttz} - 8 \, d_{14} \mathcal{T}_{zyz} + 8 \, d_{15} \mathcal{T}_{zyz} - 2 \, d_{17} \mathcal{T}_{zyz} + \tfrac{B_{0} \, \delta_{1}}{2} \mathcal{T}_{txy} - 2 \, \beta_{2} \mathcal{T}_{txz} - \tfrac{3}{2 \, \kappa^2} \mathcal{T}_{txz} + \tfrac{B_{0} \, \delta_{1}}{2} \mathcal{T}_{xty} \\
  % &\quad - 2 \, \beta_{2} \mathcal{T}_{xtz} - \tfrac{3}{2 \, \kappa^2} \mathcal{T}_{xtz} + 2 \, d_{17} \mathcal{T}_{xxy} - \tfrac{1}{2} B_{0} \, d_{19} \partial_z \mathcal{H}_{tt} + \tfrac{B_{0} \, d_{20}}{2} \partial_z \mathcal{H}_{tt} - B_{0} \, d_{21} \partial_z \mathcal{H}_{tt} - B_{0} \, {\zeta}_{1} \partial_z \mathcal{H}_{tt} + \tfrac{B_{0} \, {\zeta}_{3}}{2} \partial_z \mathcal{H}_{tt} \\
  % &\quad + B_{0} \, {\zeta}_{6} \partial_z \mathcal{H}_{tt} + \tfrac{B_{0} \, d_{19}}{2} \partial_z \mathcal{H}_{xx} - \tfrac{1}{2} B_{0} \, d_{20} \partial_z \mathcal{H}_{xx} + B_{0} \, d_{21} \partial_z \mathcal{H}_{xx} + B_{0} \, {\zeta}_{1} \partial_z \mathcal{H}_{xx} - \tfrac{1}{2} B_{0} \, {\zeta}_{3} \partial_z \mathcal{H}_{xx} - B_{0} \, {\zeta}_{6} \partial_z \mathcal{H}_{xx} \\
  % &\quad + B_{0} \, d_{21} \partial_z \mathcal{H}_{yy} - B_{0} \, {\zeta}_{1} \partial_z \mathcal{H}_{yy} - B_{0} \, {\zeta}_{2} \partial_z \mathcal{H}_{yy} - 2 \, B_{0} \, {\zeta}_{5} \partial_z \mathcal{H}_{yy} + B_{0} \, d_{21} \partial_z \mathcal{H}_{zz} - B_{0} \, {\zeta}_{1} \partial_z \mathcal{H}_{zz} - B_{0} \, {\zeta}_{2} \partial_z \mathcal{H}_{zz} \\
  % &\quad - 2 \, B_{0} \, {\zeta}_{5} \partial_z \mathcal{H}_{zz} + 2 \, d_{21} \partial_z^2 \mathcal{A}_{y} - 2 \, {\zeta}_{1} \partial_z^2 \mathcal{A}_{y} - 2 \, {\zeta}_{2} \partial_z^2 \mathcal{A}_{y} - 4 \, {\zeta}_{5} \partial_z^2 \mathcal{A}_{y} + B_{0} \, d_{19} \partial_t \mathcal{H}_{tz} - B_{0} \, d_{20} \partial_t \mathcal{H}_{tz} + 2 \, B_{0} \, d_{21} \partial_t \mathcal{H}_{tz} \\
  % &\quad + 2 \, B_{0} \, {\zeta}_{1} \partial_t \mathcal{H}_{tz} - B_{0} \, {\zeta}_{3} \partial_t \mathcal{H}_{tz} - 2 \, B_{0} \, {\zeta}_{6} \partial_t \mathcal{H}_{tz} - 2 \, \delta_{1} \partial_z \dot{\mathcal{A}}_{x} + 2 \, \zeta_{2} \partial_z \dot{\mathcal{A}}_{x} + \zeta_{3} \partial_z \dot{\mathcal{A}}_{x} - 2 \, \chi \partial_z \dot{\mathcal{T}}_{ttx} \\
  % &\quad + 2 \, \chi \partial_z \dot{\mathcal{T}}_{yxy} + 2 \, \chi \partial_z \dot{\mathcal{T}}_{zxz} + d_{19} \partial_t^2 \mathcal{A}_{y} - d_{20} \partial_t^2 \mathcal{A}_{y} + 2 \, d_{21} \partial_t^2 \mathcal{A}_{y} + 2 \, {\zeta}_{1} \partial_t^2 \mathcal{A}_{y} - {\zeta}_{3} \partial_t^2 \mathcal{A}_{y} - 2 \, {\zeta}_{6} \partial_t^2 \mathcal{A}_{y}
  % \end{aligned} \\[1ex]
  % \begin{aligned}
  %   0 &= -2 \, d_{17} \mathcal{T}_{ttx} - \tfrac{3 \, B_{0} \, d_{19}}{2} \mathcal{T}_{xyz} + \tfrac{3 \, B_{0} \, d_{20}}{2} \mathcal{T}_{xyz} - B_{0} \, d_{21} \mathcal{T}_{xyz} - B_{0} \, \delta_{1} \mathcal{T}_{yty} - 2 \, \beta_{2} \mathcal{T}_{ytz} - \tfrac{3}{2 \, \kappa^2} \mathcal{T}_{ytz} + 2 \, d_{17} \mathcal{T}_{yxy} - 4 \, \beta_{1} \mathcal{T}_{zty} \\
  % &\quad - \tfrac{2}{\kappa^2} \mathcal{T}_{zty} - B_{0} \, \delta_{1} \mathcal{T}_{ztz} + 8 \, d_{14} \mathcal{T}_{zxz} - 8 \, d_{15} \mathcal{T}_{zxz} + 2 \, d_{17} \mathcal{T}_{zxz} - 2 \, \beta_{2} \mathcal{T}_{tyz} - \tfrac{3}{2 \, \kappa^2} \mathcal{T}_{tyz} - \tfrac{3 \, B_{0} \, \delta_{1}}{2} \mathcal{T}_{xtx} + \tfrac{B_{0} \, d_{19}}{2} \partial_z \mathcal{H}_{xy} \\
  % &\quad - \tfrac{1}{2} B_{0} \, d_{20} \partial_z \mathcal{H}_{xy} + 2 \, B_{0} \, {\zeta}_{1} \partial_z \mathcal{H}_{xy} + B_{0} \, {\zeta}_{2} \partial_z \mathcal{H}_{xy} - \tfrac{1}{2} B_{0} \, {\zeta}_{3} \partial_z \mathcal{H}_{xy} + 2 \, B_{0} \, {\zeta}_{5} \partial_z \mathcal{H}_{xy} - B_{0} \, {\zeta}_{6} \partial_z \mathcal{H}_{xy} - 2 \, d_{21} \partial_z^2 \mathcal{A}_{x} \\
  % &\quad + 2 \, {\zeta}_{1} \partial_z^2 \mathcal{A}_{x} + 2 \, {\zeta}_{2} \partial_z^2 \mathcal{A}_{x} + 4 \, {\zeta}_{5} \partial_z^2 \mathcal{A}_{x} - B_{0} \, \delta_{1} \partial_t \mathcal{H}_{zz} + B_{0} \, \zeta_{2} \partial_t \mathcal{H}_{zz} + \tfrac{B_{0} \, \zeta_{3}}{2} \partial_t \mathcal{H}_{zz} - 2 \, \delta_{1} \partial_z \dot{\mathcal{A}}_{y} \\
  % &\quad + 2 \, \zeta_{2} \partial_z \dot{\mathcal{A}}_{y} + \zeta_{3} \partial_z \dot{\mathcal{A}}_{y} - 2 \, \chi \partial_z \dot{\mathcal{T}}_{tty} + 2 \, \chi \partial_z \dot{\mathcal{T}}_{zyz} - 2 \, \chi \partial_z \dot{\mathcal{T}}_{xxy} - d_{19} \partial_t^2 \mathcal{A}_{x} + d_{20} \partial_t^2 \mathcal{A}_{x} - 2 \, d_{21} \partial_t^2 \mathcal{A}_{x} \\
  % &\quad - 2 \, {\zeta}_{1} \partial_t^2 \mathcal{A}_{x} + {\zeta}_{3} \partial_t^2 \mathcal{A}_{x} + 2 \, {\zeta}_{6} \partial_t^2 \mathcal{A}_{x}
  % \end{aligned} \\[1ex]
  % \begin{aligned}
  %   \left(-\chi - \xi\right) \, \partial_t^{2} \mathcal{T}_{ttz} &= -2 \, \beta_{3} \mathcal{T}_{xxz} + \tfrac{3}{\kappa^2} \mathcal{T}_{xxz} + 2 \, d_{17} \mathcal{T}_{ytx} - 2 \, \beta_{3} \mathcal{T}_{yyz} + \tfrac{3}{\kappa^2} \mathcal{T}_{yyz} + \tfrac{3 \, B_{0} \, d_{19}}{2} \mathcal{T}_{ztx} - \tfrac{3 \, B_{0} \, d_{20}}{2} \mathcal{T}_{ztx} + 4 \, B_{0} \, d_{21} \mathcal{T}_{ztx} \\
  % &\quad - 4 \, \beta_{1} \mathcal{T}_{ttz} - 2 \, \beta_{2} \mathcal{T}_{ttz} - 2 \, \beta_{3} \mathcal{T}_{ttz} - \tfrac{1}{2 \, \kappa^2} \mathcal{T}_{ttz} - 3 \, B_{0} \, \delta_{1} \mathcal{T}_{zyz} - 8 \, d_{14} \mathcal{T}_{txy} + 8 \, d_{15} \mathcal{T}_{txy} - 2 \, d_{17} \mathcal{T}_{txy} - \tfrac{1}{2} B_{0} \, d_{19} \mathcal{T}_{txz} + \tfrac{B_{0} \, d_{20}}{2} \mathcal{T}_{txz} \\
  % &\quad + 2 \, B_{0} \, d_{21} \mathcal{T}_{txz} - 2 \, d_{17} \mathcal{T}_{xty} - \tfrac{1}{2} B_{0} \, d_{19} \mathcal{T}_{xtz} + \tfrac{B_{0} \, d_{20}}{2} \mathcal{T}_{xtz} + 2 \, B_{0} \, d_{21} \mathcal{T}_{xtz} - \tfrac{1}{2} B_{0} \, d_{19} \partial_z \mathcal{H}_{tx} + \tfrac{B_{0} \, d_{20}}{2} \partial_z \mathcal{H}_{tx} - 2 \, B_{0} \, d_{21} \partial_z \mathcal{H}_{tx} \\
  % &\quad + B_{0} \, {\zeta}_{2} \partial_z \mathcal{H}_{tx} + \tfrac{B_{0} \, {\zeta}_{3}}{2} \partial_z \mathcal{H}_{tx} + 2 \, B_{0} \, {\zeta}_{5} \partial_z \mathcal{H}_{tx} + B_{0} \, {\zeta}_{6} \partial_z \mathcal{H}_{tx} - \tfrac{1}{2} B_{0} \, \delta_{1} \partial_t \mathcal{H}_{ty} + B_{0} \, \zeta_{1} \partial_t \mathcal{H}_{ty} + 2 \, B_{0} \, \zeta_{2} \partial_t \mathcal{H}_{ty} \\
  % &\quad + B_{0} \, \zeta_{3} \partial_t \mathcal{H}_{ty} + \tfrac{B_{0} \, d_{19}}{2} \partial_t \mathcal{H}_{xz} - \tfrac{1}{2} B_{0} \, d_{20} \partial_t \mathcal{H}_{xz} + 2 \, B_{0} \, d_{21} \partial_t \mathcal{H}_{xz} - B_{0} \, {\zeta}_{2} \partial_t \mathcal{H}_{xz} - \tfrac{1}{2} B_{0} \, {\zeta}_{3} \partial_t \mathcal{H}_{xz} - 2 \, B_{0} \, {\zeta}_{5} \partial_t \mathcal{H}_{xz} \\
  % &\quad - B_{0} \, {\zeta}_{6} \partial_t \mathcal{H}_{xz} - \tfrac{1}{2} \delta_{1} \partial_z \dot{\mathcal{A}}_{t} + \zeta_{1} \partial_z \dot{\mathcal{A}}_{t} + 2 \, \zeta_{2} \partial_z \dot{\mathcal{A}}_{t} + \zeta_{3} \partial_z \dot{\mathcal{A}}_{t} - \chi \partial_z \dot{\mathcal{T}}_{yty} + \xi \partial_z \dot{\mathcal{T}}_{yty} + \chi \partial_z \dot{\mathcal{T}}_{ztz} \\
  % &\quad + \xi \partial_z \dot{\mathcal{T}}_{ztz} - \chi \partial_z \dot{\mathcal{T}}_{xtx} + \xi \partial_z \dot{\mathcal{T}}_{xtx} + \tfrac{\delta_{1}}{2} \partial_t^2 \mathcal{A}_{z} - \zeta_{1} \partial_t^2 \mathcal{A}_{z} - 2 \, \zeta_{2} \partial_t^2 \mathcal{A}_{z} - \zeta_{3} \partial_t^2 \mathcal{A}_{z} - \chi \partial_t^2 \mathcal{T}_{xxz} + \xi \partial_t^2 \mathcal{T}_{xxz} \\
  % &\quad - \chi \partial_t^2 \mathcal{T}_{yyz} + \xi \partial_t^2 \mathcal{T}_{yyz}
  % \end{aligned} \\[1ex]
  % \begin{aligned}
  %   0 &= B_{0} \, d_{19} \mathcal{T}_{ttx} - B_{0} \, d_{20} \mathcal{T}_{ttx} + 6 \, B_{0} \, d_{21} \mathcal{T}_{ttx} - 2 \, d_{17} \mathcal{T}_{xyz} - 2 \, \beta_{3} \mathcal{T}_{yty} + \tfrac{3}{\kappa^2} \mathcal{T}_{yty} + B_{0} \, \delta_{1} \mathcal{T}_{ytz} + 2 \, d_{17} \mathcal{T}_{yxz} - B_{0} \, \delta_{1} \mathcal{T}_{zty} - 4 \, \beta_{1} \mathcal{T}_{ztz} \\
  % &\quad - 2 \, \beta_{2} \mathcal{T}_{ztz} - 2 \, \beta_{3} \mathcal{T}_{ztz} - \tfrac{1}{2 \, \kappa^2} \mathcal{T}_{ztz} - 8 \, d_{14} \mathcal{T}_{zxy} + 8 \, d_{15} \mathcal{T}_{zxy} - 2 \, d_{17} \mathcal{T}_{zxy} - 2 \, B_{0} \, \delta_{1} \mathcal{T}_{tyz} - 2 \, \beta_{3} \mathcal{T}_{xtx} + \tfrac{3}{\kappa^2} \mathcal{T}_{xtx} \\
  % &\quad - \tfrac{1}{2} B_{0} \, \delta_{1} \partial_z \mathcal{H}_{ty} - B_{0} \, \zeta_{1} \partial_z \mathcal{H}_{ty} - B_{0} \, \zeta_{2} \partial_z \mathcal{H}_{ty} - \tfrac{1}{2} B_{0} \, \zeta_{3} \partial_z \mathcal{H}_{ty} + \tfrac{\delta_{1}}{2} \partial_z^2 \mathcal{A}_{t} - \zeta_{1} \partial_z^2 \mathcal{A}_{t} - 2 \, \zeta_{2} \partial_z^2 \mathcal{A}_{t} \\
  % &\quad - \zeta_{3} \partial_z^2 \mathcal{A}_{t} + \chi \partial_z^2 \mathcal{T}_{yty} - \xi \partial_z^2 \mathcal{T}_{yty} - \chi \partial_z^2 \mathcal{T}_{ztz} - \xi \partial_z^2 \mathcal{T}_{ztz} + \chi \partial_z^2 \mathcal{T}_{xtx} - \xi \partial_z^2 \mathcal{T}_{xtx} + B_{0} \, \delta_{1} \partial_t \mathcal{H}_{yz} - B_{0} \, \zeta_{2} \partial_t \mathcal{H}_{yz} \\
  % &\quad - \tfrac{1}{2} B_{0} \, \zeta_{3} \partial_t \mathcal{H}_{yz} - \tfrac{1}{2} \delta_{1} \partial_z \dot{\mathcal{A}}_{z} + \zeta_{1} \partial_z \dot{\mathcal{A}}_{z} + 2 \, \zeta_{2} \partial_z \dot{\mathcal{A}}_{z} + \zeta_{3} \partial_z \dot{\mathcal{A}}_{z} + \chi \partial_z \dot{\mathcal{T}}_{xxz} - \xi \partial_z \dot{\mathcal{T}}_{xxz} + \chi \partial_z \dot{\mathcal{T}}_{yyz} \\
  % &\quad - \xi \partial_z \dot{\mathcal{T}}_{yyz} - \chi \partial_z \dot{\mathcal{T}}_{ttz} - \xi \partial_z \dot{\mathcal{T}}_{ttz}
  % \end{aligned} \\[1ex]
  % \begin{aligned}
  %   0 &= -\tfrac{3 \, B_{0} \, \delta_{1}}{2} \mathcal{T}_{ttx} - 2 \, \beta_{2} \mathcal{T}_{xyz} - \tfrac{3}{2 \, \kappa^2} \mathcal{T}_{xyz} - 2 \, d_{17} \mathcal{T}_{yty} + B_{0} \, \delta_{1} \mathcal{T}_{yxy} + 2 \, \beta_{2} \mathcal{T}_{yxz} + \tfrac{3}{2 \, \kappa^2} \mathcal{T}_{yxz} - 8 \, d_{14} \mathcal{T}_{ztz} + 8 \, d_{15} \mathcal{T}_{ztz} - 2 \, d_{17} \mathcal{T}_{ztz} \\
  % &\quad + 4 \, \beta_{1} \mathcal{T}_{zxy} + \tfrac{2}{\kappa^2} \mathcal{T}_{zxy} + B_{0} \, \delta_{1} \mathcal{T}_{zxz} - \tfrac{3 \, B_{0} \, d_{19}}{2} \mathcal{T}_{tyz} + \tfrac{3 \, B_{0} \, d_{20}}{2} \mathcal{T}_{tyz} - B_{0} \, d_{21} \mathcal{T}_{tyz} - 2 \, d_{17} \mathcal{T}_{xtx} - \tfrac{1}{2} B_{0} \, d_{19} \partial_z \mathcal{H}_{ty} \\
  % &\quad + \tfrac{B_{0} \, d_{20}}{2} \partial_z \mathcal{H}_{ty} - 2 \, B_{0} \, {\zeta}_{1} \partial_z \mathcal{H}_{ty} - B_{0} \, {\zeta}_{2} \partial_z \mathcal{H}_{ty} + \tfrac{B_{0} \, {\zeta}_{3}}{2} \partial_z \mathcal{H}_{ty} - 2 \, B_{0} \, {\zeta}_{5} \partial_z \mathcal{H}_{ty} + B_{0} \, {\zeta}_{6} \partial_z \mathcal{H}_{ty} + 2 \, d_{21} \partial_z^2 \mathcal{A}_{t} \\
  % &\quad - 2 \, {\zeta}_{1} \partial_z^2 \mathcal{A}_{t} - 2 \, {\zeta}_{2} \partial_z^2 \mathcal{A}_{t} - 4 \, {\zeta}_{5} \partial_z^2 \mathcal{A}_{t} + \tfrac{B_{0} \, d_{19}}{2} \partial_t \mathcal{H}_{yz} - \tfrac{1}{2} B_{0} \, d_{20} \partial_t \mathcal{H}_{yz} + 2 \, B_{0} \, d_{21} \partial_t \mathcal{H}_{yz} - B_{0} \, {\zeta}_{2} \partial_t \mathcal{H}_{yz} \\
  % &\quad - \tfrac{1}{2} B_{0} \, {\zeta}_{3} \partial_t \mathcal{H}_{yz} - 2 \, B_{0} \, {\zeta}_{5} \partial_t \mathcal{H}_{yz} - B_{0} \, {\zeta}_{6} \partial_t \mathcal{H}_{yz} - 2 \, d_{21} \partial_z \dot{\mathcal{A}}_{z} + 2 \, {\zeta}_{1} \partial_z \dot{\mathcal{A}}_{z} + 2 \, {\zeta}_{2} \partial_z \dot{\mathcal{A}}_{z} + 4 \, {\zeta}_{5} \partial_z \dot{\mathcal{A}}_{z}
  % \end{aligned} \\[1ex]
  % \begin{aligned}
  %   \left(3 \, \chi - \xi\right) \, \partial_t^{2} \mathcal{T}_{zxz} &= 2 \, \beta_{3} \mathcal{T}_{ttx} - \tfrac{3}{\kappa^2} \mathcal{T}_{ttx} + 2 \, B_{0} \, \delta_{1} \mathcal{T}_{xyz} + 2 \, d_{17} \mathcal{T}_{ytz} - 2 \, \beta_{3} \mathcal{T}_{yxy} + \tfrac{3}{\kappa^2} \mathcal{T}_{yxy} + B_{0} \, \delta_{1} \mathcal{T}_{yxz} - 8 \, d_{14} \mathcal{T}_{zty} + 8 \, d_{15} \mathcal{T}_{zty} \\
  % &\quad - 2 \, d_{17} \mathcal{T}_{zty} - B_{0} \, \delta_{1} \mathcal{T}_{zxy} - 4 \, \beta_{1} \mathcal{T}_{zxz} - 2 \, \beta_{2} \mathcal{T}_{zxz} - 2 \, \beta_{3} \mathcal{T}_{zxz} - \tfrac{1}{2 \, \kappa^2} \mathcal{T}_{zxz} + 2 \, d_{17} \mathcal{T}_{tyz} - B_{0} \, d_{19} \mathcal{T}_{xtx} + B_{0} \, d_{20} \mathcal{T}_{xtx} - 6 \, B_{0} \, d_{21} \mathcal{T}_{xtx} \\
  % &\quad - \tfrac{1}{2} B_{0} \, \delta_{1} \partial_z \mathcal{H}_{xy} - B_{0} \, \zeta_{1} \partial_z \mathcal{H}_{xy} - B_{0} \, \zeta_{2} \partial_z \mathcal{H}_{xy} - \tfrac{1}{2} B_{0} \, \zeta_{3} \partial_z \mathcal{H}_{xy} + \tfrac{\delta_{1}}{2} \partial_z^2 \mathcal{A}_{x} - \zeta_{1} \partial_z^2 \mathcal{A}_{x} - 2 \, \zeta_{2} \partial_z^2 \mathcal{A}_{x} \\
  % &\quad - \zeta_{3} \partial_z^2 \mathcal{A}_{x} - \chi \partial_z^2 \mathcal{T}_{ttx} + \xi \partial_z^2 \mathcal{T}_{ttx} + \chi \partial_z^2 \mathcal{T}_{yxy} - \xi \partial_z^2 \mathcal{T}_{yxy} - \chi \partial_z^2 \mathcal{T}_{zxz} - \xi \partial_z^2 \mathcal{T}_{zxz} - \tfrac{1}{2} B_{0} \, d_{19} \partial_t \mathcal{H}_{zz} + \tfrac{B_{0} \, d_{20}}{2} \partial_t \mathcal{H}_{zz} \\
  % &\quad - 2 \, B_{0} \, d_{21} \partial_t \mathcal{H}_{zz} + B_{0} \, {\zeta}_{2} \partial_t \mathcal{H}_{zz} + \tfrac{B_{0} \, {\zeta}_{3}}{2} \partial_t \mathcal{H}_{zz} + 2 \, B_{0} \, {\zeta}_{5} \partial_t \mathcal{H}_{zz} + B_{0} \, {\zeta}_{6} \partial_t \mathcal{H}_{zz} - d_{19} \partial_z \dot{\mathcal{A}}_{y} + d_{20} \partial_z \dot{\mathcal{A}}_{y} \\
  % &\quad - 4 \, d_{21} \partial_z \dot{\mathcal{A}}_{y} + 2 \, {\zeta}_{2} \partial_z \dot{\mathcal{A}}_{y} + {\zeta}_{3} \partial_z \dot{\mathcal{A}}_{y} + 4 \, {\zeta}_{5} \partial_z \dot{\mathcal{A}}_{y} + 2 \, {\zeta}_{6} \partial_z \dot{\mathcal{A}}_{y} - 2 \, \chi \partial_z \dot{\mathcal{T}}_{ztx} - 2 \, \chi \partial_z \dot{\mathcal{T}}_{txz} \\
  % &\quad + \tfrac{3 \, \delta_{1}}{2} \partial_t^2 \mathcal{A}_{x} + \zeta_{1} \partial_t^2 \mathcal{A}_{x} + \chi \partial_t^2 \mathcal{T}_{ttx} - \xi \partial_t^2 \mathcal{T}_{ttx} - 3 \, \chi \partial_t^2 \mathcal{T}_{yxy} + \xi \partial_t^2 \mathcal{T}_{yxy}
  % \end{aligned} \\[1ex]
  % \begin{aligned}
  %   \left(3 \, \chi - \xi\right) \, \partial_t^{2} \mathcal{T}_{zyz} &= 2 \, \beta_{3} \mathcal{T}_{tty} - \tfrac{3}{\kappa^2} \mathcal{T}_{tty} - 3 \, B_{0} \, \delta_{1} \mathcal{T}_{xxz} - 3 \, B_{0} \, d_{21} \mathcal{T}_{ytx} + 8 \, d_{14} \mathcal{T}_{ztx} - 8 \, d_{15} \mathcal{T}_{ztx} + 2 \, d_{17} \mathcal{T}_{ztx} - 3 \, B_{0} \, \delta_{1} \mathcal{T}_{ttz} \\
  % &\quad - 4 \, \beta_{1} \mathcal{T}_{zyz} - 2 \, \beta_{2} \mathcal{T}_{zyz} - 2 \, \beta_{3} \mathcal{T}_{zyz} - \tfrac{1}{2 \, \kappa^2} \mathcal{T}_{zyz} - B_{0} \, d_{19} \mathcal{T}_{txy} + B_{0} \, d_{20} \mathcal{T}_{txy} - 3 \, B_{0} \, d_{21} \mathcal{T}_{txy} - 2 \, d_{17} \mathcal{T}_{txz} - B_{0} \, d_{19} \mathcal{T}_{xty} + B_{0} \, d_{20} \mathcal{T}_{xty} \\
  % &\quad - 3 \, B_{0} \, d_{21} \mathcal{T}_{xty} - 2 \, d_{17} \mathcal{T}_{xtz} + 2 \, \beta_{3} \mathcal{T}_{xxy} - \tfrac{3}{\kappa^2} \mathcal{T}_{xxy} - \tfrac{3 \, B_{0} \, \delta_{1}}{4} \partial_z \mathcal{H}_{tt} - \tfrac{1}{2} B_{0} \, \zeta_{1} \partial_z \mathcal{H}_{tt} + \tfrac{3 \, B_{0} \, \delta_{1}}{4} \partial_z \mathcal{H}_{xx} \\
  % &\quad + \tfrac{B_{0} \, \zeta_{1}}{2} \partial_z \mathcal{H}_{xx} + \tfrac{B_{0} \, \delta_{1}}{4} \partial_z \mathcal{H}_{yy} - \tfrac{1}{2} B_{0} \, \zeta_{1} \partial_z \mathcal{H}_{yy} - B_{0} \, \zeta_{2} \partial_z \mathcal{H}_{yy} - \tfrac{1}{2} B_{0} \, \zeta_{3} \partial_z \mathcal{H}_{yy} + \tfrac{B_{0} \, \delta_{1}}{4} \partial_z \mathcal{H}_{zz} \\
  % &\quad - \tfrac{1}{2} B_{0} \, \zeta_{1} \partial_z \mathcal{H}_{zz} - B_{0} \, \zeta_{2} \partial_z \mathcal{H}_{zz} - \tfrac{1}{2} B_{0} \, \zeta_{3} \partial_z \mathcal{H}_{zz} + \tfrac{\delta_{1}}{2} \partial_z^2 \mathcal{A}_{y} - \zeta_{1} \partial_z^2 \mathcal{A}_{y} - 2 \, \zeta_{2} \partial_z^2 \mathcal{A}_{y} - \zeta_{3} \partial_z^2 \mathcal{A}_{y} - \chi \partial_z^2 \mathcal{T}_{tty} \\
  % &\quad + \xi \partial_z^2 \mathcal{T}_{tty} - \chi \partial_z^2 \mathcal{T}_{zyz} - \xi \partial_z^2 \mathcal{T}_{zyz} - \chi \partial_z^2 \mathcal{T}_{xxy} + \xi \partial_z^2 \mathcal{T}_{xxy} + \tfrac{3 \, B_{0} \, \delta_{1}}{2} \partial_t \mathcal{H}_{tz} + B_{0} \, \zeta_{1} \partial_t \mathcal{H}_{tz} + d_{19} \partial_z \dot{\mathcal{A}}_{x} - d_{20} \partial_z \dot{\mathcal{A}}_{x} \\
  % &\quad + 4 \, d_{21} \partial_z \dot{\mathcal{A}}_{x} - 2 \, {\zeta}_{2} \partial_z \dot{\mathcal{A}}_{x} - {\zeta}_{3} \partial_z \dot{\mathcal{A}}_{x} - 4 \, {\zeta}_{5} \partial_z \dot{\mathcal{A}}_{x} - 2 \, {\zeta}_{6} \partial_z \dot{\mathcal{A}}_{x} - 2 \, \chi \partial_z \dot{\mathcal{T}}_{zty} - 2 \, \chi \partial_z \dot{\mathcal{T}}_{tyz} \\
  % &\quad + \tfrac{3 \, \delta_{1}}{2} \partial_t^2 \mathcal{A}_{y} + \zeta_{1} \partial_t^2 \mathcal{A}_{y} + \chi \partial_t^2 \mathcal{T}_{tty} - \xi \partial_t^2 \mathcal{T}_{tty} + 3 \, \chi \partial_t^2 \mathcal{T}_{xxy} - \xi \partial_t^2 \mathcal{T}_{xxy}
  % \end{aligned} \\[1ex]
  % \begin{aligned}
  %   0 &= \tfrac{B_{0} \, d_{19}}{2} \mathcal{T}_{tty} - \tfrac{1}{2} B_{0} \, d_{20} \mathcal{T}_{tty} - 2 \, B_{0} \, d_{21} \mathcal{T}_{tty} + 2 \, d_{17} \mathcal{T}_{xxz} - 2 \, \beta_{2} \mathcal{T}_{ytx} - \tfrac{3}{2 \, \kappa^2} \mathcal{T}_{ytx} + 2 \, d_{17} \mathcal{T}_{yyz} + \tfrac{B_{0} \, \delta_{1}}{2} \mathcal{T}_{ztx} + 8 \, d_{14} \mathcal{T}_{ttz} \\
  % &\quad - 8 \, d_{15} \mathcal{T}_{ttz} + 2 \, d_{17} \mathcal{T}_{ttz} + B_{0} \, d_{19} \mathcal{T}_{zyz} - B_{0} \, d_{20} \mathcal{T}_{zyz} + 3 \, B_{0} \, d_{21} \mathcal{T}_{zyz} - 4 \, \beta_{1} \mathcal{T}_{txy} - \tfrac{2}{\kappa^2} \mathcal{T}_{txy} + 2 \, \beta_{2} \mathcal{T}_{xty} + \tfrac{3}{2 \, \kappa^2} \mathcal{T}_{xty} + \tfrac{B_{0} \, d_{19}}{2} \mathcal{T}_{xxy} \\
  % &\quad - \tfrac{1}{2} B_{0} \, d_{20} \mathcal{T}_{xxy} - 2 \, B_{0} \, d_{21} \mathcal{T}_{xxy} - B_{0} \, \delta_{1} \partial_z \mathcal{H}_{tx} + B_{0} \, \zeta_{2} \partial_z \mathcal{H}_{tx} + \tfrac{B_{0} \, \zeta_{3}}{2} \partial_z \mathcal{H}_{tx} + 2 \, B_{0} \, d_{21} \partial_t \mathcal{H}_{ty} - 2 \, B_{0} \, {\zeta}_{1} \partial_t \mathcal{H}_{ty} \\
  % &\quad - 2 \, B_{0} \, {\zeta}_{2} \partial_t \mathcal{H}_{ty} - 4 \, B_{0} \, {\zeta}_{5} \partial_t \mathcal{H}_{ty} + B_{0} \, \delta_{1} \partial_t \mathcal{H}_{xz} - B_{0} \, \zeta_{2} \partial_t \mathcal{H}_{xz} - \tfrac{1}{2} B_{0} \, \zeta_{3} \partial_t \mathcal{H}_{xz} + 2 \, d_{21} \partial_z \dot{\mathcal{A}}_{t} - 2 \, {\zeta}_{1} \partial_z \dot{\mathcal{A}}_{t} \\
  % &\quad - 2 \, {\zeta}_{2} \partial_z \dot{\mathcal{A}}_{t} - 4 \, {\zeta}_{5} \partial_z \dot{\mathcal{A}}_{t} - 2 \, d_{21} \partial_t^2 \mathcal{A}_{z} + 2 \, {\zeta}_{1} \partial_t^2 \mathcal{A}_{z} + 2 \, {\zeta}_{2} \partial_t^2 \mathcal{A}_{z} + 4 \, {\zeta}_{5} \partial_t^2 \mathcal{A}_{z}
  % \end{aligned} \\[1ex]
  % \begin{aligned}
  %   0 &= -8 \, d_{14} \mathcal{T}_{tty} + 8 \, d_{15} \mathcal{T}_{tty} - 2 \, d_{17} \mathcal{T}_{tty} + \tfrac{B_{0} \, d_{19}}{2} \mathcal{T}_{xxz} - \tfrac{1}{2} B_{0} \, d_{20} \mathcal{T}_{xxz} - 2 \, B_{0} \, d_{21} \mathcal{T}_{xxz} - \tfrac{1}{2} B_{0} \, \delta_{1} \mathcal{T}_{ytx} - B_{0} \, d_{19} \mathcal{T}_{yyz} + B_{0} \, d_{20} \mathcal{T}_{yyz} \\
  % &\quad - 3 \, B_{0} \, d_{21} \mathcal{T}_{yyz} - 2 \, \beta_{2} \mathcal{T}_{ztx} - \tfrac{3}{2 \, \kappa^2} \mathcal{T}_{ztx} + \tfrac{B_{0} \, d_{19}}{2} \mathcal{T}_{ttz} - \tfrac{1}{2} B_{0} \, d_{20} \mathcal{T}_{ttz} - 2 \, B_{0} \, d_{21} \mathcal{T}_{ttz} + 2 \, d_{17} \mathcal{T}_{zyz} - 4 \, \beta_{1} \mathcal{T}_{txz} - \tfrac{2}{\kappa^2} \mathcal{T}_{txz} \\
  % &\quad + 2 \, \beta_{2} \mathcal{T}_{xtz} + \tfrac{3}{2 \, \kappa^2} \mathcal{T}_{xtz} - 2 \, d_{17} \mathcal{T}_{xxy} - B_{0} \, d_{21} \partial_z \mathcal{H}_{tt} + B_{0} \, {\zeta}_{1} \partial_z \mathcal{H}_{tt} + B_{0} \, {\zeta}_{2} \partial_z \mathcal{H}_{tt} + 2 \, B_{0} \, {\zeta}_{5} \partial_z \mathcal{H}_{tt} - \tfrac{1}{2} B_{0} \, d_{19} \partial_z \mathcal{H}_{xx} \\
  % &\quad + \tfrac{B_{0} \, d_{20}}{2} \partial_z \mathcal{H}_{xx} - B_{0} \, d_{21} \partial_z \mathcal{H}_{xx} - B_{0} \, {\zeta}_{1} \partial_z \mathcal{H}_{xx} + \tfrac{B_{0} \, {\zeta}_{3}}{2} \partial_z \mathcal{H}_{xx} + B_{0} \, {\zeta}_{6} \partial_z \mathcal{H}_{xx} + \tfrac{B_{0} \, d_{19}}{2} \partial_z \mathcal{H}_{yy} - \tfrac{1}{2} B_{0} \, d_{20} \partial_z \mathcal{H}_{yy} \\
  % &\quad + B_{0} \, d_{21} \partial_z \mathcal{H}_{yy} + B_{0} \, {\zeta}_{1} \partial_z \mathcal{H}_{yy} - \tfrac{1}{2} B_{0} \, {\zeta}_{3} \partial_z \mathcal{H}_{yy} - B_{0} \, {\zeta}_{6} \partial_z \mathcal{H}_{yy} + \tfrac{B_{0} \, d_{19}}{2} \partial_z \mathcal{H}_{zz} - \tfrac{1}{2} B_{0} \, d_{20} \partial_z \mathcal{H}_{zz} + B_{0} \, d_{21} \partial_z \mathcal{H}_{zz} \\
  % &\quad + B_{0} \, {\zeta}_{1} \partial_z \mathcal{H}_{zz} - \tfrac{1}{2} B_{0} \, {\zeta}_{3} \partial_z \mathcal{H}_{zz} - B_{0} \, {\zeta}_{6} \partial_z \mathcal{H}_{zz} + d_{19} \partial_z^2 \mathcal{A}_{y} - d_{20} \partial_z^2 \mathcal{A}_{y} + 2 \, d_{21} \partial_z^2 \mathcal{A}_{y} + 2 \, {\zeta}_{1} \partial_z^2 \mathcal{A}_{y} - {\zeta}_{3} \partial_z^2 \mathcal{A}_{y} \\
  % &\quad - 2 \, {\zeta}_{6} \partial_z^2 \mathcal{A}_{y} + 2 \, B_{0} \, d_{21} \partial_t \mathcal{H}_{tz} - 2 \, B_{0} \, {\zeta}_{1} \partial_t \mathcal{H}_{tz} - 2 \, B_{0} \, {\zeta}_{2} \partial_t \mathcal{H}_{tz} - 4 \, B_{0} \, {\zeta}_{5} \partial_t \mathcal{H}_{tz} - B_{0} \, \delta_{1} \partial_t \mathcal{H}_{xy} + B_{0} \, \zeta_{2} \partial_t \mathcal{H}_{xy} \\
  % &\quad + \tfrac{B_{0} \, \zeta_{3}}{2} \partial_t \mathcal{H}_{xy} - 2 \, \delta_{1} \partial_z \dot{\mathcal{A}}_{x} + 2 \, \zeta_{2} \partial_z \dot{\mathcal{A}}_{x} + \zeta_{3} \partial_z \dot{\mathcal{A}}_{x} - 2 \, \chi \partial_z \dot{\mathcal{T}}_{ttx} + 2 \, \chi \partial_z \dot{\mathcal{T}}_{yxy} + 2 \, \chi \partial_z \dot{\mathcal{T}}_{zxz} + 2 \, d_{21} \partial_t^2 \mathcal{A}_{y} \\
  % &\quad - 2 \, {\zeta}_{1} \partial_t^2 \mathcal{A}_{y} - 2 \, {\zeta}_{2} \partial_t^2 \mathcal{A}_{y} - 4 \, {\zeta}_{5} \partial_t^2 \mathcal{A}_{y}
  % \end{aligned} \\[1ex]
  % \begin{aligned}
  %   0 &= 8 \, d_{14} \mathcal{T}_{ttx} - 8 \, d_{15} \mathcal{T}_{ttx} + 2 \, d_{17} \mathcal{T}_{ttx} - 2 \, B_{0} \, \delta_{1} \mathcal{T}_{yty} + 2 \, \beta_{2} \mathcal{T}_{ytz} + \tfrac{3}{2 \, \kappa^2} \mathcal{T}_{ytz} - 2 \, d_{17} \mathcal{T}_{yxy} + \tfrac{3 \, B_{0} \, d_{19}}{2} \mathcal{T}_{yxz} - \tfrac{3 \, B_{0} \, d_{20}}{2} \mathcal{T}_{yxz} + B_{0} \, d_{21} \mathcal{T}_{yxz} \\
  % &\quad - 2 \, \beta_{2} \mathcal{T}_{zty} - \tfrac{3}{2 \, \kappa^2} \mathcal{T}_{zty} - 2 \, B_{0} \, \delta_{1} \mathcal{T}_{ztz} - \tfrac{3 \, B_{0} \, d_{19}}{2} \mathcal{T}_{zxy} + \tfrac{3 \, B_{0} \, d_{20}}{2} \mathcal{T}_{zxy} - B_{0} \, d_{21} \mathcal{T}_{zxy} - 2 \, d_{17} \mathcal{T}_{zxz} - 4 \, \beta_{1} \mathcal{T}_{tyz} - \tfrac{2}{\kappa^2} \mathcal{T}_{tyz} \\
  % &\quad - \tfrac{3 \, B_{0} \, \delta_{1}}{2} \mathcal{T}_{xtx} - B_{0} \, d_{19} \partial_z \mathcal{H}_{xy} + B_{0} \, d_{20} \partial_z \mathcal{H}_{xy} - 2 \, B_{0} \, d_{21} \partial_z \mathcal{H}_{xy} - 2 \, B_{0} \, {\zeta}_{1} \partial_z \mathcal{H}_{xy} + B_{0} \, {\zeta}_{3} \partial_z \mathcal{H}_{xy} + 2 \, B_{0} \, {\zeta}_{6} \partial_z \mathcal{H}_{xy} \\
  % &\quad - d_{19} \partial_z^2 \mathcal{A}_{x} + d_{20} \partial_z^2 \mathcal{A}_{x} - 2 \, d_{21} \partial_z^2 \mathcal{A}_{x} - 2 \, {\zeta}_{1} \partial_z^2 \mathcal{A}_{x} + {\zeta}_{3} \partial_z^2 \mathcal{A}_{x} + 2 \, {\zeta}_{6} \partial_z^2 \mathcal{A}_{x} - B_{0} \, \delta_{1} \partial_t \mathcal{H}_{yy} + B_{0} \, \zeta_{2} \partial_t \mathcal{H}_{yy} \\
  % &\quad + \tfrac{B_{0} \, \zeta_{3}}{2} \partial_t \mathcal{H}_{yy} - B_{0} \, \delta_{1} \partial_t \mathcal{H}_{zz} + B_{0} \, \zeta_{2} \partial_t \mathcal{H}_{zz} + \tfrac{B_{0} \, \zeta_{3}}{2} \partial_t \mathcal{H}_{zz} - 2 \, \delta_{1} \partial_z \dot{\mathcal{A}}_{y} + 2 \, \zeta_{2} \partial_z \dot{\mathcal{A}}_{y} + \zeta_{3} \partial_z \dot{\mathcal{A}}_{y} \\
  % &\quad - 2 \, \chi \partial_z \dot{\mathcal{T}}_{tty} + 2 \, \chi \partial_z \dot{\mathcal{T}}_{zyz} - 2 \, \chi \partial_z \dot{\mathcal{T}}_{xxy} - 2 \, d_{21} \partial_t^2 \mathcal{A}_{x} + 2 \, {\zeta}_{1} \partial_t^2 \mathcal{A}_{x} + 2 \, {\zeta}_{2} \partial_t^2 \mathcal{A}_{x} + 4 \, {\zeta}_{5} \partial_t^2 \mathcal{A}_{x}
  % \end{aligned} \\[1ex]
  % \begin{aligned}
  %   0 &= -8 \, d_{14} \mathcal{T}_{xyz} + 8 \, d_{15} \mathcal{T}_{xyz} - 2 \, d_{17} \mathcal{T}_{xyz} - 2 \, \beta_{3} \mathcal{T}_{yty} + \tfrac{3}{\kappa^2} \mathcal{T}_{yty} + \tfrac{3 \, B_{0} \, \delta_{1}}{2} \mathcal{T}_{ytz} + B_{0} \, d_{19} \mathcal{T}_{yxy} - B_{0} \, d_{20} \mathcal{T}_{yxy} + 6 \, B_{0} \, d_{21} \mathcal{T}_{yxy} + 2 \, d_{17} \mathcal{T}_{yxz} \\
  % &\quad - \tfrac{3 \, B_{0} \, \delta_{1}}{2} \mathcal{T}_{zty} - 2 \, \beta_{3} \mathcal{T}_{ztz} + \tfrac{3}{\kappa^2} \mathcal{T}_{ztz} - 2 \, d_{17} \mathcal{T}_{zxy} + B_{0} \, d_{19} \mathcal{T}_{zxz} - B_{0} \, d_{20} \mathcal{T}_{zxz} + 6 \, B_{0} \, d_{21} \mathcal{T}_{zxz} - \tfrac{3 \, B_{0} \, \delta_{1}}{2} \mathcal{T}_{tyz} - 4 \, \beta_{1} \mathcal{T}_{xtx} \\
  % &\quad - 2 \, \beta_{2} \mathcal{T}_{xtx} - 2 \, \beta_{3} \mathcal{T}_{xtx} - \tfrac{1}{2 \, \kappa^2} \mathcal{T}_{xtx} - \tfrac{3 \, B_{0} \, \delta_{1}}{2} \partial_z \mathcal{H}_{ty} - B_{0} \, \zeta_{1} \partial_z \mathcal{H}_{ty} - \tfrac{3 \, \delta_{1}}{2} \partial_z^2 \mathcal{A}_{t} - \zeta_{1} \partial_z^2 \mathcal{A}_{t} + 3 \, \chi \partial_z^2 \mathcal{T}_{yty} \\
  % &\quad - \xi \partial_z^2 \mathcal{T}_{yty} + \chi \partial_z^2 \mathcal{T}_{ztz} - \xi \partial_z^2 \mathcal{T}_{ztz} + 3 \, \chi \partial_z^2 \mathcal{T}_{xtx} - \xi \partial_z^2 \mathcal{T}_{xtx} + \tfrac{3 \, \delta_{1}}{2} \partial_z \dot{\mathcal{A}}_{z} + \zeta_{1} \partial_z \dot{\mathcal{A}}_{z} + 3 \, \chi \partial_z \dot{\mathcal{T}}_{xxz} - \xi \partial_z \dot{\mathcal{T}}_{xxz} \\
  % &\quad + 3 \, \chi \partial_z \dot{\mathcal{T}}_{yyz} - \xi \partial_z \dot{\mathcal{T}}_{yyz} + \chi \partial_z \dot{\mathcal{T}}_{ttz} - \xi \partial_z \dot{\mathcal{T}}_{ttz}
  % \end{aligned} \\[1ex]
  % \begin{aligned}
  %   0 &= \tfrac{B_{0} \, d_{19}}{2} \mathcal{T}_{tty} - \tfrac{1}{2} B_{0} \, d_{20} \mathcal{T}_{tty} - 2 \, B_{0} \, d_{21} \mathcal{T}_{tty} + 8 \, d_{14} \mathcal{T}_{xxz} - 8 \, d_{15} \mathcal{T}_{xxz} + 2 \, d_{17} \mathcal{T}_{xxz} - 2 \, \beta_{2} \mathcal{T}_{ytx} - \tfrac{3}{2 \, \kappa^2} \mathcal{T}_{ytx} + 2 \, d_{17} \mathcal{T}_{yyz} \\
  % &\quad + \tfrac{B_{0} \, \delta_{1}}{2} \mathcal{T}_{ztx} + 2 \, d_{17} \mathcal{T}_{ttz} + B_{0} \, d_{19} \mathcal{T}_{zyz} - B_{0} \, d_{20} \mathcal{T}_{zyz} + 3 \, B_{0} \, d_{21} \mathcal{T}_{zyz} + 2 \, \beta_{2} \mathcal{T}_{txy} + \tfrac{3}{2 \, \kappa^2} \mathcal{T}_{txy} - 4 \, \beta_{1} \mathcal{T}_{xty} - \tfrac{2}{\kappa^2} \mathcal{T}_{xty} \\
  % &\quad + \tfrac{B_{0} \, d_{19}}{2} \mathcal{T}_{xxy} - \tfrac{1}{2} B_{0} \, d_{20} \mathcal{T}_{xxy} - 2 \, B_{0} \, d_{21} \mathcal{T}_{xxy} + B_{0} \, \delta_{1} \partial_z \mathcal{H}_{tx} - B_{0} \, \zeta_{2} \partial_z \mathcal{H}_{tx} - \tfrac{1}{2} B_{0} \, \zeta_{3} \partial_z \mathcal{H}_{tx} - B_{0} \, d_{19} \partial_t \mathcal{H}_{ty} + B_{0} \, d_{20} \partial_t \mathcal{H}_{ty} \\
  % &\quad - 2 \, B_{0} \, d_{21} \partial_t \mathcal{H}_{ty} - 2 \, B_{0} \, {\zeta}_{1} \partial_t \mathcal{H}_{ty} + B_{0} \, {\zeta}_{3} \partial_t \mathcal{H}_{ty} + 2 \, B_{0} \, {\zeta}_{6} \partial_t \mathcal{H}_{ty} - B_{0} \, \delta_{1} \partial_t \mathcal{H}_{xz} + B_{0} \, \zeta_{2} \partial_t \mathcal{H}_{xz} + \tfrac{B_{0} \, \zeta_{3}}{2} \partial_t \mathcal{H}_{xz} \\
  % &\quad - d_{19} \partial_z \dot{\mathcal{A}}_{t} + d_{20} \partial_z \dot{\mathcal{A}}_{t} - 2 \, d_{21} \partial_z \dot{\mathcal{A}}_{t} - 2 \, {\zeta}_{1} \partial_z \dot{\mathcal{A}}_{t} + {\zeta}_{3} \partial_z \dot{\mathcal{A}}_{t} + 2 \, {\zeta}_{6} \partial_z \dot{\mathcal{A}}_{t} + d_{19} \partial_t^2 \mathcal{A}_{z} - d_{20} \partial_t^2 \mathcal{A}_{z} \\
  % &\quad + 2 \, d_{21} \partial_t^2 \mathcal{A}_{z} + 2 \, {\zeta}_{1} \partial_t^2 \mathcal{A}_{z} - {\zeta}_{3} \partial_t^2 \mathcal{A}_{z} - 2 \, {\zeta}_{6} \partial_t^2 \mathcal{A}_{z}
  % \end{aligned} \\[1ex]
  % \begin{aligned}
  %   0 &= -2 \, d_{17} \mathcal{T}_{tty} + \tfrac{B_{0} \, d_{19}}{2} \mathcal{T}_{xxz} - \tfrac{1}{2} B_{0} \, d_{20} \mathcal{T}_{xxz} - 2 \, B_{0} \, d_{21} \mathcal{T}_{xxz} - \tfrac{1}{2} B_{0} \, \delta_{1} \mathcal{T}_{ytx} - B_{0} \, d_{19} \mathcal{T}_{yyz} + B_{0} \, d_{20} \mathcal{T}_{yyz} - 3 \, B_{0} \, d_{21} \mathcal{T}_{yyz} - 2 \, \beta_{2} \mathcal{T}_{ztx} \\
  % &\quad - \tfrac{3}{2 \, \kappa^2} \mathcal{T}_{ztx} + \tfrac{B_{0} \, d_{19}}{2} \mathcal{T}_{ttz} - \tfrac{1}{2} B_{0} \, d_{20} \mathcal{T}_{ttz} - 2 \, B_{0} \, d_{21} \mathcal{T}_{ttz} + 2 \, d_{17} \mathcal{T}_{zyz} + 2 \, \beta_{2} \mathcal{T}_{txz} + \tfrac{3}{2 \, \kappa^2} \mathcal{T}_{txz} - 4 \, \beta_{1} \mathcal{T}_{xtz} - \tfrac{2}{\kappa^2} \mathcal{T}_{xtz} \\
  % &\quad - 8 \, d_{14} \mathcal{T}_{xxy} + 8 \, d_{15} \mathcal{T}_{xxy} - 2 \, d_{17} \mathcal{T}_{xxy} + \tfrac{B_{0} \, d_{19}}{2} \partial_z \mathcal{H}_{tt} - \tfrac{1}{2} B_{0} \, d_{20} \partial_z \mathcal{H}_{tt} + B_{0} \, d_{21} \partial_z \mathcal{H}_{tt} + B_{0} \, {\zeta}_{1} \partial_z \mathcal{H}_{tt} - \tfrac{1}{2} B_{0} \, {\zeta}_{3} \partial_z \mathcal{H}_{tt} \\
  % &\quad - B_{0} \, {\zeta}_{6} \partial_z \mathcal{H}_{tt} + B_{0} \, d_{21} \partial_z \mathcal{H}_{xx} - B_{0} \, {\zeta}_{1} \partial_z \mathcal{H}_{xx} - B_{0} \, {\zeta}_{2} \partial_z \mathcal{H}_{xx} - 2 \, B_{0} \, {\zeta}_{5} \partial_z \mathcal{H}_{xx} + \tfrac{B_{0} \, d_{19}}{2} \partial_z \mathcal{H}_{yy} - \tfrac{1}{2} B_{0} \, d_{20} \partial_z \mathcal{H}_{yy} \\
  % &\quad + B_{0} \, d_{21} \partial_z \mathcal{H}_{yy} + B_{0} \, {\zeta}_{1} \partial_z \mathcal{H}_{yy} - \tfrac{1}{2} B_{0} \, {\zeta}_{3} \partial_z \mathcal{H}_{yy} - B_{0} \, {\zeta}_{6} \partial_z \mathcal{H}_{yy} + \tfrac{B_{0} \, d_{19}}{2} \partial_z \mathcal{H}_{zz} - \tfrac{1}{2} B_{0} \, d_{20} \partial_z \mathcal{H}_{zz} + B_{0} \, d_{21} \partial_z \mathcal{H}_{zz} \\
  % &\quad + B_{0} \, {\zeta}_{1} \partial_z \mathcal{H}_{zz} - \tfrac{1}{2} B_{0} \, {\zeta}_{3} \partial_z \mathcal{H}_{zz} - B_{0} \, {\zeta}_{6} \partial_z \mathcal{H}_{zz} + d_{19} \partial_z^2 \mathcal{A}_{y} - d_{20} \partial_z^2 \mathcal{A}_{y} + 2 \, d_{21} \partial_z^2 \mathcal{A}_{y} + 2 \, {\zeta}_{1} \partial_z^2 \mathcal{A}_{y} - {\zeta}_{3} \partial_z^2 \mathcal{A}_{y} \\
  % &\quad - 2 \, {\zeta}_{6} \partial_z^2 \mathcal{A}_{y} - B_{0} \, d_{19} \partial_t \mathcal{H}_{tz} + B_{0} \, d_{20} \partial_t \mathcal{H}_{tz} - 2 \, B_{0} \, d_{21} \partial_t \mathcal{H}_{tz} - 2 \, B_{0} \, {\zeta}_{1} \partial_t \mathcal{H}_{tz} + B_{0} \, {\zeta}_{3} \partial_t \mathcal{H}_{tz} + 2 \, B_{0} \, {\zeta}_{6} \partial_t \mathcal{H}_{tz} \\
  % &\quad + B_{0} \, \delta_{1} \partial_t \mathcal{H}_{xy} - B_{0} \, \zeta_{2} \partial_t \mathcal{H}_{xy} - \tfrac{1}{2} B_{0} \, \zeta_{3} \partial_t \mathcal{H}_{xy} - d_{19} \partial_t^2 \mathcal{A}_{y} + d_{20} \partial_t^2 \mathcal{A}_{y} - 2 \, d_{21} \partial_t^2 \mathcal{A}_{y} - 2 \, {\zeta}_{1} \partial_t^2 \mathcal{A}_{y} + {\zeta}_{3} \partial_t^2 \mathcal{A}_{y} \\
  % &\quad + 2 \, {\zeta}_{6} \partial_t^2 \mathcal{A}_{y}
  % \end{aligned} \\[1ex]
  % \begin{aligned}
  %   \left(3 \, \chi - \xi\right) \, \partial_t^{2} \mathcal{T}_{xxy} &= -2 \, \beta_{3} \mathcal{T}_{tty} + \tfrac{3}{\kappa^2} \mathcal{T}_{tty} + \tfrac{3 \, B_{0} \, d_{19}}{2} \mathcal{T}_{ytx} - \tfrac{3 \, B_{0} \, d_{20}}{2} \mathcal{T}_{ytx} + 4 \, B_{0} \, d_{21} \mathcal{T}_{ytx} - 3 \, B_{0} \, \delta_{1} \mathcal{T}_{yyz} - 2 \, d_{17} \mathcal{T}_{ztx} \\
  % &\quad + 2 \, \beta_{3} \mathcal{T}_{zyz} - \tfrac{3}{\kappa^2} \mathcal{T}_{zyz} - \tfrac{1}{2} B_{0} \, d_{19} \mathcal{T}_{txy} + \tfrac{B_{0} \, d_{20}}{2} \mathcal{T}_{txy} + 2 \, B_{0} \, d_{21} \mathcal{T}_{txy} + 2 \, d_{17} \mathcal{T}_{txz} - \tfrac{1}{2} B_{0} \, d_{19} \mathcal{T}_{xty} + \tfrac{B_{0} \, d_{20}}{2} \mathcal{T}_{xty} + 2 \, B_{0} \, d_{21} \mathcal{T}_{xty} \\
  % &\quad + 8 \, d_{14} \mathcal{T}_{xtz} - 8 \, d_{15} \mathcal{T}_{xtz} + 2 \, d_{17} \mathcal{T}_{xtz} - 4 \, \beta_{1} \mathcal{T}_{xxy} - 2 \, \beta_{2} \mathcal{T}_{xxy} - 2 \, \beta_{3} \mathcal{T}_{xxy} - \tfrac{1}{2 \, \kappa^2} \mathcal{T}_{xxy} + \tfrac{3 \, B_{0} \, \delta_{1}}{4} \partial_z \mathcal{H}_{tt} + \tfrac{B_{0} \, \zeta_{1}}{2} \partial_z \mathcal{H}_{tt} \\
  % &\quad + \tfrac{B_{0} \, \delta_{1}}{4} \partial_z \mathcal{H}_{xx} - \tfrac{1}{2} B_{0} \, \zeta_{1} \partial_z \mathcal{H}_{xx} - B_{0} \, \zeta_{2} \partial_z \mathcal{H}_{xx} - \tfrac{1}{2} B_{0} \, \zeta_{3} \partial_z \mathcal{H}_{xx} + \tfrac{3 \, B_{0} \, \delta_{1}}{4} \partial_z \mathcal{H}_{yy} + \tfrac{B_{0} \, \zeta_{1}}{2} \partial_z \mathcal{H}_{yy} \\
  % &\quad + \tfrac{3 \, B_{0} \, \delta_{1}}{4} \partial_z \mathcal{H}_{zz} + \tfrac{B_{0} \, \zeta_{1}}{2} \partial_z \mathcal{H}_{zz} + \tfrac{3 \, \delta_{1}}{2} \partial_z^2 \mathcal{A}_{y} + \zeta_{1} \partial_z^2 \mathcal{A}_{y} + 3 \, \chi \partial_z^2 \mathcal{T}_{tty} - \xi \partial_z^2 \mathcal{T}_{tty} - \chi \partial_z^2 \mathcal{T}_{zyz} + \xi \partial_z^2 \mathcal{T}_{zyz} \\
  % &\quad + 3 \, \chi \partial_z^2 \mathcal{T}_{xxy} - \xi \partial_z^2 \mathcal{T}_{xxy} - \tfrac{3 \, B_{0} \, \delta_{1}}{2} \partial_t \mathcal{H}_{tz} - B_{0} \, \zeta_{1} \partial_t \mathcal{H}_{tz} - \tfrac{1}{2} B_{0} \, d_{19} \partial_t \mathcal{H}_{xy} + \tfrac{B_{0} \, d_{20}}{2} \partial_t \mathcal{H}_{xy} - 2 \, B_{0} \, d_{21} \partial_t \mathcal{H}_{xy} \\
  % &\quad + B_{0} \, {\zeta}_{2} \partial_t \mathcal{H}_{xy} + \tfrac{B_{0} \, {\zeta}_{3}}{2} \partial_t \mathcal{H}_{xy} + 2 \, B_{0} \, {\zeta}_{5} \partial_t \mathcal{H}_{xy} + B_{0} \, {\zeta}_{6} \partial_t \mathcal{H}_{xy} + 2 \, \chi \partial_z \dot{\mathcal{T}}_{zty} + 2 \, \chi \partial_z \dot{\mathcal{T}}_{tyz} - \tfrac{3 \, \delta_{1}}{2} \partial_t^2 \mathcal{A}_{y} \\
  % &\quad - \zeta_{1} \partial_t^2 \mathcal{A}_{y} - \chi \partial_t^2 \mathcal{T}_{tty} + \xi \partial_t^2 \mathcal{T}_{tty} + 3 \, \chi \partial_t^2 \mathcal{T}_{zyz} - \xi \partial_t^2 \mathcal{T}_{zyz}
  % \end{aligned} \\[1ex]
  \begin{aligned}
    \mathscr{H} &\supset -\tfrac{1}{16} B_{0}^2 \mathcal{H}_{tt}^2 + \tfrac{B_{0}^2}{4} \mathcal{H}_{tx}^2 - \tfrac{1}{4} B_{0}^2 \mathcal{H}_{ty}^2 - \tfrac{1}{4} B_{0}^2 \mathcal{H}_{tz}^2 - \tfrac{1}{8} B_{0}^2 \mathcal{H}_{tt} \, \mathcal{H}_{xx} - \tfrac{1}{16} B_{0}^2 \mathcal{H}_{xx}^2 + \tfrac{B_{0}^2}{4} \mathcal{H}_{xy}^2 + \tfrac{B_{0}^2}{4} \mathcal{H}_{xz}^2 \\
  &\quad + \tfrac{B_{0}^2}{8} \mathcal{H}_{tt} \, \mathcal{H}_{yy} - \tfrac{1}{8} B_{0}^2 \mathcal{H}_{xx} \, \mathcal{H}_{yy} + \tfrac{3 \, B_{0}^2}{16} \mathcal{H}_{yy}^2 + \tfrac{B_{0}^2}{4} \mathcal{H}_{yz}^2 + \tfrac{B_{0}^2}{8} \mathcal{H}_{tt} \, \mathcal{H}_{zz} - \tfrac{1}{8} B_{0}^2 \mathcal{H}_{xx} \, \mathcal{H}_{zz} + \tfrac{B_{0}^2}{8} \mathcal{H}_{yy} \, \mathcal{H}_{zz} \\
  &\quad + \tfrac{3 \, B_{0}^2}{16} \mathcal{H}_{zz}^2 - \tfrac{1}{2} \partial_z \mathcal{A}_{t}^2 + \tfrac{1}{2} \partial_z \mathcal{A}_{x}^2 + \tfrac{1}{2} \partial_z \mathcal{A}_{y}^2 + \tfrac{1}{2} \partial_z \mathcal{A}_{z}^2 - \tfrac{1}{2 \, \kappa^2} \partial_z \mathcal{H}_{tx}^2 - \tfrac{1}{2 \, \kappa^2} \partial_z \mathcal{H}_{ty}^2 \\
  &\quad + \tfrac{1}{2 \, \kappa^2} \partial_z \mathcal{H}_{tt} \, \partial_z \mathcal{H}_{xx} + \tfrac{1}{2 \, \kappa^2} \partial_z \mathcal{H}_{xy}^2 + \tfrac{1}{2 \, \kappa^2} \partial_z \mathcal{H}_{tt} \, \partial_z \mathcal{H}_{yy} - \tfrac{1}{2 \, \kappa^2} \partial_z \mathcal{H}_{xx} \, \partial_z \mathcal{H}_{yy} - \tfrac{1}{2} \dot{\mathcal{A}}_{t}^2 \\
  &\quad + \tfrac{1}{2} \dot{\mathcal{A}}_{x}^2 + \tfrac{1}{2} \dot{\mathcal{A}}_{y}^2 + \tfrac{1}{2} \dot{\mathcal{A}}_{z}^2 + \tfrac{1}{2 \, \kappa^2} \dot{\mathcal{H}}_{xy}^2 + \tfrac{1}{2 \, \kappa^2} \dot{\mathcal{H}}_{xz}^2 - \tfrac{1}{2 \, \kappa^2} \dot{\mathcal{H}}_{xx} \, \dot{\mathcal{H}}_{yy} + \tfrac{1}{2 \, \kappa^2} \dot{\mathcal{H}}_{yz}^2 \\
  &\quad - \tfrac{1}{2 \, \kappa^2} \dot{\mathcal{H}}_{xx} \, \dot{\mathcal{H}}_{zz} - \tfrac{1}{2 \, \kappa^2} \dot{\mathcal{H}}_{yy} \, \dot{\mathcal{H}}_{zz} - \tfrac{2}{\kappa^2} \mathcal{H}_{xz} \, \partial_t \partial_z \mathcal{H}_{tx} - \tfrac{2}{\kappa^2} \mathcal{H}_{yz} \, \partial_t \partial_z \mathcal{H}_{ty} \\
  &\quad + \tfrac{1}{\kappa^{2}} \mathcal{H}_{tt} \, \partial_t \partial_z \mathcal{H}_{tz} + \tfrac{1}{\kappa^{2}} \mathcal{H}_{xx} \, \partial_t \partial_z \mathcal{H}_{tz} + \tfrac{1}{\kappa^{2}} \mathcal{H}_{yy} \, \partial_t \partial_z \mathcal{H}_{tz} - \tfrac{1}{\kappa^{2}} \mathcal{H}_{zz} \, \partial_t \partial_z \mathcal{H}_{tz} \\
  &\quad + \tfrac{2}{\kappa^2} \mathcal{H}_{tz} \, \partial_t \partial_z \mathcal{H}_{xx} - \tfrac{2}{\kappa^2} \mathcal{H}_{tx} \, \partial_t \partial_z \mathcal{H}_{xz} + \tfrac{2}{\kappa^2} \mathcal{H}_{tz} \, \partial_t \partial_z \mathcal{H}_{yy} - \tfrac{2}{\kappa^2} \mathcal{H}_{ty} \, \partial_t \partial_z \mathcal{H}_{yz}
  \end{aligned}

\end{gathered}
$}
\medskip


% --- E.8  Dark-photon-plasma analogue -------------------------------
\subsection{Dark-photon-plasma}%
\label{SymbolicDarkPhotonPlasma}
A kinetic-mixing vertex
$\delta_{m}\,F_{\mu\nu}\tilde{F}^{\mu\nu}$ (with $\tilde{F}$ the
torsion-trace field strength) and a photon mass term
$m_{A}^{2}$ are added in \cref{DarkPhotonPlasmaLagrangian}.
$m_{A}^{2}$ enters the diagonal photon kinetic term;
$\delta_{m}$ enters the off-diagonal photon--torsion entries
(\cref{SectorSurvey}). The symbolic output comprises
\DarkPhotonPlasmaEomCount{} equations of motion
(\DarkPhotonPlasmaConstraintCount{} constraints) with
\DarkPhotonPlasmaTermCount{} RHS terms.

\medskip
\noindent\adjustbox{max width=\linewidth}{$\displaystyle
\setlength{\jot}{1pt}%
\begin{gathered}
% !TEX root = ../../../main.tex
  \begin{aligned}
    \mathscr{L} &= \tfrac{1}{\kappa^2} \mathcal{R} - \alpha_{3} \tensor{T}{^b_b_a} \tensor{g}{^a^d} \tensor{T}{^c_c_d} - \tfrac{1}{4} \xi \tensor{\tilde{F}}{_a_b} \tensor{g}{^a^c} \tensor{g}{^b^d} \tensor{\tilde{F}}{_c_d} + \delta_m \tensor{\mathcal{F}}{_a_b} \tensor{g}{^a^c} \tensor{g}{^b^d} \tensor{\tilde{F}}{_c_d} - \tfrac{1}{4} \tensor{\mathcal{F}}{_a_b} \tensor{g}{^a^c} \tensor{g}{^b^d} \tensor{\mathcal{F}}{_c_d} - \tfrac{1}{2} m_{A}^{2} \tensor{a}{_a} \tensor{g}{^a^b} \tensor{a}{_b}
  \end{aligned} \\[1ex]
  \begin{aligned}
    \partial_t^{2} A_{t} &= m_{A}^{2} A_{t} - \partial_z^2 A_{t} -2 \, \delta_m \partial_z^2 \mathcal{T}_{yty} -2 \, \delta_m \partial_z^2 \mathcal{T}_{ztz} -2 \, \delta_m \partial_z^2 \mathcal{T}_{xtx} -2 \, \delta_m \partial_z \dot{\mathcal{T}}_{xxz} -2 \, \delta_m \partial_z \dot{\mathcal{T}}_{yyz} -2 \, \delta_m \partial_z \dot{\mathcal{T}}_{ttz} \\
    \partial_t^{2} A_{x} &= -m_{A}^{2} A_{x} + B_{0} \partial_z h_{xy} + \partial_z^2 A_{x} - 2 \, \delta_m \partial_z^2 \mathcal{T}_{ttx} + 2 \, \delta_m \partial_z^2 \mathcal{T}_{yxy} + 2 \, \delta_m \partial_z^2 \mathcal{T}_{zxz} + 2 \, \delta_m \partial_t^2 \mathcal{T}_{ttx} - 2 \, \delta_m \partial_t^2 \mathcal{T}_{yxy} - 2 \, \delta_m \partial_t^2 \mathcal{T}_{zxz} \\
    \partial_t^{2} A_{y} &= -m_{A}^{2} A_{y} + B_{0} \partial_z h_{yy} + \partial_z^2 A_{y} - 2 \, \delta_m \partial_z^2 \mathcal{T}_{tty} + 2 \, \delta_m \partial_z^2 \mathcal{T}_{zyz} - 2 \, \delta_m \partial_z^2 \mathcal{T}_{xxy} + 2 \, \delta_m \partial_t^2 \mathcal{T}_{tty} - 2 \, \delta_m \partial_t^2 \mathcal{T}_{zyz} + 2 \, \delta_m \partial_t^2 \mathcal{T}_{xxy} \\
    \partial_t^{2} A_{z} &= -m_{A}^{2} A_{z} + \partial_z^2 A_{z} + 2 \, \delta_m \partial_z \dot{\mathcal{T}}_{yty} + 2 \, \delta_m \partial_z \dot{\mathcal{T}}_{ztz} + 2 \, \delta_m \partial_z \dot{\mathcal{T}}_{xtx} + 2 \, \delta_m \partial_t^2 \mathcal{T}_{xxz} + 2 \, \delta_m \partial_t^2 \mathcal{T}_{yyz} + 2 \, \delta_m \partial_t^2 \mathcal{T}_{ttz}
  \end{aligned} \\[1ex]
  \begin{aligned}
    -\tfrac{1}{\kappa^{2}} \, \partial_t^{2} h_{xy} &= \tfrac{B_{0}^2}{2} h_{xy} + B_{0} \partial_z A_{x} -2 \, B_{0} \, \delta_m \partial_z \mathcal{T}_{ttx} + 2 \, B_{0} \, \delta_m \partial_z \mathcal{T}_{yxy} + 2 \, B_{0} \, \delta_m \partial_z \mathcal{T}_{zxz} -\tfrac{1}{\kappa^{2}} \partial_z^2 h_{xy} \\
    -\tfrac{1}{\kappa^{2}} \, \partial_t^{2} h_{yy} &= \tfrac{B_{0}^2}{2} h_{yy} + B_{0} \partial_z A_{y} -2 \, B_{0} \, \delta_m \partial_z \mathcal{T}_{tty} + 2 \, B_{0} \, \delta_m \partial_z \mathcal{T}_{zyz} -2 \, B_{0} \, \delta_m \partial_z \mathcal{T}_{xxy} -\tfrac{1}{\kappa^{2}} \partial_z^2 h_{yy}
  \end{aligned} \\[1ex]
  \begin{aligned}
    -\xi \, \partial_t^{2} \mathcal{T}_{ttx} &= 2 \, \alpha_{3} \mathcal{T}_{ttx} - 2 \, \alpha_{3} \mathcal{T}_{yxy} - 2 \, \alpha_{3} \mathcal{T}_{zxz} + 2 \, B_{0} \, \delta_m \partial_z h_{xy} + 2 \, \delta_m \partial_z^2 A_{x} - \xi \partial_z^2 \mathcal{T}_{ttx} + \xi \partial_z^2 \mathcal{T}_{yxy} + \xi \partial_z^2 \mathcal{T}_{zxz} - 2 \, \delta_m \partial_t^2 A_{x} - \xi \partial_t^2 \mathcal{T}_{yxy} - \xi \partial_t^2 \mathcal{T}_{zxz} \\
    -\xi \, \partial_t^{2} \mathcal{T}_{tty} &= 2 \, \alpha_{3} \mathcal{T}_{tty} - 2 \, \alpha_{3} \mathcal{T}_{zyz} + 2 \, \alpha_{3} \mathcal{T}_{xxy} + 2 \, B_{0} \, \delta_m \partial_z h_{yy} + 2 \, \delta_m \partial_z^2 A_{y} - \xi \partial_z^2 \mathcal{T}_{tty} + \xi \partial_z^2 \mathcal{T}_{zyz} - \xi \partial_z^2 \mathcal{T}_{xxy} - 2 \, \delta_m \partial_t^2 A_{y} - \xi \partial_t^2 \mathcal{T}_{zyz} + \xi \partial_t^2 \mathcal{T}_{xxy} \\
    -\xi \, \partial_t^{2} \mathcal{T}_{xxz} &= 2 \, \alpha_{3} \mathcal{T}_{xxz} + 2 \, \alpha_{3} \mathcal{T}_{yyz} + 2 \, \alpha_{3} \mathcal{T}_{ttz} + 2 \, \delta_m \partial_z \dot{A}_{t} + \xi \partial_z \dot{\mathcal{T}}_{yty} + \xi \partial_z \dot{\mathcal{T}}_{ztz} + \xi \partial_z \dot{\mathcal{T}}_{xtx} - 2 \, \delta_m \partial_t^2 A_{z} + \xi \partial_t^2 \mathcal{T}_{yyz} + \xi \partial_t^2 \mathcal{T}_{ttz} \\
    0 &= 2 \, \alpha_{3} \mathcal{T}_{yty} + 2 \, \alpha_{3} \mathcal{T}_{ztz} + 2 \, \alpha_{3} \mathcal{T}_{xtx} - 2 \, \delta_m \partial_z^2 A_{t} - \xi \partial_z^2 \mathcal{T}_{yty} - \xi \partial_z^2 \mathcal{T}_{ztz} - \xi \partial_z^2 \mathcal{T}_{xtx} + 2 \, \delta_m \partial_z \dot{A}_{z} - \xi \partial_z \dot{\mathcal{T}}_{xxz} - \xi \partial_z \dot{\mathcal{T}}_{yyz} - \xi \partial_z \dot{\mathcal{T}}_{ttz} \\
    -\xi \, \partial_t^{2} \mathcal{T}_{yxy} &= -2 \, \alpha_{3} \mathcal{T}_{ttx} + 2 \, \alpha_{3} \mathcal{T}_{yxy} + 2 \, \alpha_{3} \mathcal{T}_{zxz} - 2 \, B_{0} \, \delta_m \partial_z h_{xy} - 2 \, \delta_m \partial_z^2 A_{x} + \xi \partial_z^2 \mathcal{T}_{ttx} - \xi \partial_z^2 \mathcal{T}_{yxy} - \xi \partial_z^2 \mathcal{T}_{zxz} + 2 \, \delta_m \partial_t^2 A_{x} - \xi \partial_t^2 \mathcal{T}_{ttx} + \xi \partial_t^2 \mathcal{T}_{zxz} \\
    -\xi \, \partial_t^{2} \mathcal{T}_{yyz} &= 2 \, \alpha_{3} \mathcal{T}_{xxz} + 2 \, \alpha_{3} \mathcal{T}_{yyz} + 2 \, \alpha_{3} \mathcal{T}_{ttz} + 2 \, \delta_m \partial_z \dot{A}_{t} + \xi \partial_z \dot{\mathcal{T}}_{yty} + \xi \partial_z \dot{\mathcal{T}}_{ztz} + \xi \partial_z \dot{\mathcal{T}}_{xtx} - 2 \, \delta_m \partial_t^2 A_{z} + \xi \partial_t^2 \mathcal{T}_{xxz} + \xi \partial_t^2 \mathcal{T}_{ttz} \\
    -\xi \, \partial_t^{2} \mathcal{T}_{ttz} &= 2 \, \alpha_{3} \mathcal{T}_{xxz} + 2 \, \alpha_{3} \mathcal{T}_{yyz} + 2 \, \alpha_{3} \mathcal{T}_{ttz} + 2 \, \delta_m \partial_z \dot{A}_{t} + \xi \partial_z \dot{\mathcal{T}}_{yty} + \xi \partial_z \dot{\mathcal{T}}_{ztz} + \xi \partial_z \dot{\mathcal{T}}_{xtx} - 2 \, \delta_m \partial_t^2 A_{z} + \xi \partial_t^2 \mathcal{T}_{xxz} + \xi \partial_t^2 \mathcal{T}_{yyz} \\
    0 &= 2 \, \alpha_{3} \mathcal{T}_{yty} + 2 \, \alpha_{3} \mathcal{T}_{ztz} + 2 \, \alpha_{3} \mathcal{T}_{xtx} - 2 \, \delta_m \partial_z^2 A_{t} - \xi \partial_z^2 \mathcal{T}_{yty} - \xi \partial_z^2 \mathcal{T}_{ztz} - \xi \partial_z^2 \mathcal{T}_{xtx} + 2 \, \delta_m \partial_z \dot{A}_{z} - \xi \partial_z \dot{\mathcal{T}}_{xxz} - \xi \partial_z \dot{\mathcal{T}}_{yyz} - \xi \partial_z \dot{\mathcal{T}}_{ttz} \\
    -\xi \, \partial_t^{2} \mathcal{T}_{zxz} &= -2 \, \alpha_{3} \mathcal{T}_{ttx} + 2 \, \alpha_{3} \mathcal{T}_{yxy} + 2 \, \alpha_{3} \mathcal{T}_{zxz} - 2 \, B_{0} \, \delta_m \partial_z h_{xy} - 2 \, \delta_m \partial_z^2 A_{x} + \xi \partial_z^2 \mathcal{T}_{ttx} - \xi \partial_z^2 \mathcal{T}_{yxy} - \xi \partial_z^2 \mathcal{T}_{zxz} + 2 \, \delta_m \partial_t^2 A_{x} - \xi \partial_t^2 \mathcal{T}_{ttx} + \xi \partial_t^2 \mathcal{T}_{yxy} \\
    -\xi \, \partial_t^{2} \mathcal{T}_{zyz} &= -2 \, \alpha_{3} \mathcal{T}_{tty} + 2 \, \alpha_{3} \mathcal{T}_{zyz} - 2 \, \alpha_{3} \mathcal{T}_{xxy} - 2 \, B_{0} \, \delta_m \partial_z h_{yy} - 2 \, \delta_m \partial_z^2 A_{y} + \xi \partial_z^2 \mathcal{T}_{tty} - \xi \partial_z^2 \mathcal{T}_{zyz} + \xi \partial_z^2 \mathcal{T}_{xxy} + 2 \, \delta_m \partial_t^2 A_{y} - \xi \partial_t^2 \mathcal{T}_{tty} - \xi \partial_t^2 \mathcal{T}_{xxy} \\
    0 &= 2 \, \alpha_{3} \mathcal{T}_{yty} + 2 \, \alpha_{3} \mathcal{T}_{ztz} + 2 \, \alpha_{3} \mathcal{T}_{xtx} - 2 \, \delta_m \partial_z^2 A_{t} - \xi \partial_z^2 \mathcal{T}_{yty} - \xi \partial_z^2 \mathcal{T}_{ztz} - \xi \partial_z^2 \mathcal{T}_{xtx} + 2 \, \delta_m \partial_z \dot{A}_{z} - \xi \partial_z \dot{\mathcal{T}}_{xxz} - \xi \partial_z \dot{\mathcal{T}}_{yyz} - \xi \partial_z \dot{\mathcal{T}}_{ttz} \\
    -\xi \, \partial_t^{2} \mathcal{T}_{xxy} &= 2 \, \alpha_{3} \mathcal{T}_{tty} - 2 \, \alpha_{3} \mathcal{T}_{zyz} + 2 \, \alpha_{3} \mathcal{T}_{xxy} + 2 \, B_{0} \, \delta_m \partial_z h_{yy} + 2 \, \delta_m \partial_z^2 A_{y} - \xi \partial_z^2 \mathcal{T}_{tty} + \xi \partial_z^2 \mathcal{T}_{zyz} - \xi \partial_z^2 \mathcal{T}_{xxy} - 2 \, \delta_m \partial_t^2 A_{y} + \xi \partial_t^2 \mathcal{T}_{tty} - \xi \partial_t^2 \mathcal{T}_{zyz}
  \end{aligned} \\[1ex]
  \begin{aligned}
    \mathscr{H} &\supset -\tfrac{1}{2} m_{A}^{2} A_{t}^2 + \tfrac{m_{A}^{2}}{2} A_{x}^2 + \tfrac{m_{A}^{2}}{2} A_{y}^2 + \tfrac{m_{A}^{2}}{2} A_{z}^2 + \tfrac{B_{0}^2}{4} h_{xy}^2 + \tfrac{B_{0}^2}{4} h_{yy}^2 - \tfrac{1}{2} \partial_z A_{t}^2 + \tfrac{1}{2} \partial_z A_{x}^2 \\
  &\quad + \tfrac{1}{2} \partial_z A_{y}^2 + \tfrac{1}{2} \partial_z A_{z}^2 + \tfrac{1}{2 \, \kappa^2} \partial_z h_{xy}^2 + \tfrac{1}{2 \, \kappa^2} \partial_z h_{yy}^2 - \tfrac{1}{2} \dot{A}_{t}^2 + \tfrac{1}{2} \dot{A}_{x}^2 + \tfrac{1}{2} \dot{A}_{y}^2 + \tfrac{1}{2} \dot{A}_{z}^2 \\
  &\quad + \tfrac{1}{2 \, \kappa^2} \dot{h}_{xy}^2 + \tfrac{1}{2 \, \kappa^2} \dot{h}_{yy}^2
  \end{aligned}

\end{gathered}
$}
\medskip

% The $\delta_{m}$ kinetic-mixing vertex produces off-diagonal
% photon--torsion entries explicitly visible in the
% $A \leftrightarrow T$ block of the kinetic matrix below;
% the photon-mass term $m_A^{2}$ enters the diagonal $A_\mu$ cells.

% \medskip
% \noindent\adjustbox{max width=\linewidth}{$\displaystyle
% \setlength{\jot}{0.5pt}%
% \mathcal{K}(\partial_t,\partial_z) = \begin{bmatrix}
  \partial_t^2 - m_{A}^{2} + \partial_z^2 & 0 & 0 & 0 & 0 & 0 & 0 & 0 & 0 & 2 \, \delta_m \, \partial_z^2 & 0 & 0 & 0 & 2 \, \delta_m \, \partial_z^2 & 0 & 0 & 2 \, \delta_m \, \partial_z^2 & 0 \\
  0 & \partial_t^2 + m_{A}^{2} - \partial_z^2 & 0 & 0 & -B_{0} \, \partial_z & 0 & 2 \, \delta_m \, \partial_z^2 - 2 \, \delta_m \, \partial_t^2 & 0 & 0 & 0 & -2 \, \delta_m \, \partial_z^2 + 2 \, \delta_m \, \partial_t^2 & 0 & 0 & 0 & -2 \, \delta_m \, \partial_z^2 + 2 \, \delta_m \, \partial_t^2 & 0 & 0 & 0 \\
  0 & 0 & \partial_t^2 + m_{A}^{2} - \partial_z^2 & 0 & 0 & -B_{0} \, \partial_z & 0 & 2 \, \delta_m \, \partial_z^2 - 2 \, \delta_m \, \partial_t^2 & 0 & 0 & 0 & 0 & 0 & 0 & 0 & -2 \, \delta_m \, \partial_z^2 + 2 \, \delta_m \, \partial_t^2 & 0 & 2 \, \delta_m \, \partial_z^2 - 2 \, \delta_m \, \partial_t^2 \\
  0 & 0 & 0 & \partial_t^2 + m_{A}^{2} - \partial_z^2 & 0 & 0 & 0 & 0 & -2 \, \delta_m \, \partial_t^2 & 0 & 0 & -2 \, \delta_m \, \partial_t^2 & -2 \, \delta_m \, \partial_t^2 & 0 & 0 & 0 & 0 & 0 \\
  0 & -B_{0} \, \partial_z & 0 & 0 & -\tfrac{1}{\kappa^{2}} \, \partial_t^2 - \tfrac{B_{0}^2}{2} + \tfrac{1}{\kappa^{2}} \, \partial_z^2 & 0 & 2 \, B_{0} \, \delta_m \, \partial_z & 0 & 0 & 0 & -2 \, B_{0} \, \delta_m \, \partial_z & 0 & 0 & 0 & -2 \, B_{0} \, \delta_m \, \partial_z & 0 & 0 & 0 \\
  0 & 0 & -B_{0} \, \partial_z & 0 & 0 & -\tfrac{1}{\kappa^{2}} \, \partial_t^2 - \tfrac{B_{0}^2}{2} + \tfrac{1}{\kappa^{2}} \, \partial_z^2 & 0 & 2 \, B_{0} \, \delta_m \, \partial_z & 0 & 0 & 0 & 0 & 0 & 0 & 0 & -2 \, B_{0} \, \delta_m \, \partial_z & 0 & 2 \, B_{0} \, \delta_m \, \partial_z \\
  0 & -2 \, \delta_m \, \partial_z^2 + 2 \, \delta_m \, \partial_t^2 & 0 & 0 & -2 \, B_{0} \, \delta_m \, \partial_z & 0 & -\xi \, \partial_t^2 - 2 \, \alpha_{3} + \xi \, \partial_z^2 & 0 & 0 & 0 & 2 \, \alpha_{3} - \xi \, \partial_z^2 + \xi \, \partial_t^2 & 0 & 0 & 0 & 2 \, \alpha_{3} - \xi \, \partial_z^2 + \xi \, \partial_t^2 & 0 & 0 & 0 \\
  0 & 0 & -2 \, \delta_m \, \partial_z^2 + 2 \, \delta_m \, \partial_t^2 & 0 & 0 & -2 \, B_{0} \, \delta_m \, \partial_z & 0 & -\xi \, \partial_t^2 - 2 \, \alpha_{3} + \xi \, \partial_z^2 & 0 & 0 & 0 & 0 & 0 & 0 & 0 & 2 \, \alpha_{3} - \xi \, \partial_z^2 + \xi \, \partial_t^2 & 0 & -2 \, \alpha_{3} + \xi \, \partial_z^2 - \xi \, \partial_t^2 \\
  0 & 0 & 0 & 2 \, \delta_m \, \partial_t^2 & 0 & 0 & 0 & 0 & -\xi \, \partial_t^2 - 2 \, \alpha_{3} & 0 & 0 & -2 \, \alpha_{3} - \xi \, \partial_t^2 & -2 \, \alpha_{3} - \xi \, \partial_t^2 & 0 & 0 & 0 & 0 & 0 \\
  2 \, \delta_m \, \partial_z^2 & 0 & 0 & 0 & 0 & 0 & 0 & 0 & 0 & 1 - 2 \, \alpha_{3} + \xi \, \partial_z^2 & 0 & 0 & 0 & -2 \, \alpha_{3} + \xi \, \partial_z^2 & 0 & 0 & -2 \, \alpha_{3} + \xi \, \partial_z^2 & 0 \\
  0 & 2 \, \delta_m \, \partial_z^2 - 2 \, \delta_m \, \partial_t^2 & 0 & 0 & 2 \, B_{0} \, \delta_m \, \partial_z & 0 & 2 \, \alpha_{3} - \xi \, \partial_z^2 + \xi \, \partial_t^2 & 0 & 0 & 0 & -\xi \, \partial_t^2 - 2 \, \alpha_{3} + \xi \, \partial_z^2 & 0 & 0 & 0 & -2 \, \alpha_{3} + \xi \, \partial_z^2 - \xi \, \partial_t^2 & 0 & 0 & 0 \\
  0 & 0 & 0 & 2 \, \delta_m \, \partial_t^2 & 0 & 0 & 0 & 0 & -2 \, \alpha_{3} - \xi \, \partial_t^2 & 0 & 0 & -\xi \, \partial_t^2 - 2 \, \alpha_{3} & -2 \, \alpha_{3} - \xi \, \partial_t^2 & 0 & 0 & 0 & 0 & 0 \\
  0 & 0 & 0 & 2 \, \delta_m \, \partial_t^2 & 0 & 0 & 0 & 0 & -2 \, \alpha_{3} - \xi \, \partial_t^2 & 0 & 0 & -2 \, \alpha_{3} - \xi \, \partial_t^2 & -\xi \, \partial_t^2 - 2 \, \alpha_{3} & 0 & 0 & 0 & 0 & 0 \\
  2 \, \delta_m \, \partial_z^2 & 0 & 0 & 0 & 0 & 0 & 0 & 0 & 0 & -2 \, \alpha_{3} + \xi \, \partial_z^2 & 0 & 0 & 0 & 1 - 2 \, \alpha_{3} + \xi \, \partial_z^2 & 0 & 0 & -2 \, \alpha_{3} + \xi \, \partial_z^2 & 0 \\
  0 & 2 \, \delta_m \, \partial_z^2 - 2 \, \delta_m \, \partial_t^2 & 0 & 0 & 2 \, B_{0} \, \delta_m \, \partial_z & 0 & 2 \, \alpha_{3} - \xi \, \partial_z^2 + \xi \, \partial_t^2 & 0 & 0 & 0 & -2 \, \alpha_{3} + \xi \, \partial_z^2 - \xi \, \partial_t^2 & 0 & 0 & 0 & -\xi \, \partial_t^2 - 2 \, \alpha_{3} + \xi \, \partial_z^2 & 0 & 0 & 0 \\
  0 & 0 & 2 \, \delta_m \, \partial_z^2 - 2 \, \delta_m \, \partial_t^2 & 0 & 0 & 2 \, B_{0} \, \delta_m \, \partial_z & 0 & 2 \, \alpha_{3} - \xi \, \partial_z^2 + \xi \, \partial_t^2 & 0 & 0 & 0 & 0 & 0 & 0 & 0 & -\xi \, \partial_t^2 - 2 \, \alpha_{3} + \xi \, \partial_z^2 & 0 & 2 \, \alpha_{3} - \xi \, \partial_z^2 + \xi \, \partial_t^2 \\
  2 \, \delta_m \, \partial_z^2 & 0 & 0 & 0 & 0 & 0 & 0 & 0 & 0 & -2 \, \alpha_{3} + \xi \, \partial_z^2 & 0 & 0 & 0 & -2 \, \alpha_{3} + \xi \, \partial_z^2 & 0 & 0 & 1 - 2 \, \alpha_{3} + \xi \, \partial_z^2 & 0 \\
  0 & 0 & -2 \, \delta_m \, \partial_z^2 + 2 \, \delta_m \, \partial_t^2 & 0 & 0 & -2 \, B_{0} \, \delta_m \, \partial_z & 0 & -2 \, \alpha_{3} + \xi \, \partial_z^2 - \xi \, \partial_t^2 & 0 & 0 & 0 & 0 & 0 & 0 & 0 & 2 \, \alpha_{3} - \xi \, \partial_z^2 + \xi \, \partial_t^2 & 0 & -\xi \, \partial_t^2 - 2 \, \alpha_{3} + \xi \, \partial_z^2
\end{bmatrix}

% $}
% \medskip
